% Options for packages loaded elsewhere
\PassOptionsToPackage{unicode}{hyperref}
\PassOptionsToPackage{hyphens}{url}
\documentclass[
  indonesian,
  12pt,
]{report}
\usepackage{xcolor}
\usepackage[a4paper,top=2.5cm,bottom=2.5cm,left=3cm,right=2.5cm]{geometry}
\usepackage{amsmath,amssymb}
\setcounter{secnumdepth}{5}
\usepackage{iftex}
\ifPDFTeX
  \usepackage[T1]{fontenc}
  \usepackage[utf8]{inputenc}
  \usepackage{textcomp} % provide euro and other symbols
\else % if luatex or xetex
  \usepackage{unicode-math} % this also loads fontspec
  \defaultfontfeatures{Scale=MatchLowercase}
  \defaultfontfeatures[\rmfamily]{Ligatures=TeX,Scale=1}
\fi
\usepackage{lmodern}
\ifPDFTeX\else
  % xetex/luatex font selection
\fi
% Use upquote if available, for straight quotes in verbatim environments
\IfFileExists{upquote.sty}{\usepackage{upquote}}{}
\IfFileExists{microtype.sty}{% use microtype if available
  \usepackage[]{microtype}
  \UseMicrotypeSet[protrusion]{basicmath} % disable protrusion for tt fonts
}{}
\usepackage{setspace}
\makeatletter
\@ifundefined{KOMAClassName}{% if non-KOMA class
  \IfFileExists{parskip.sty}{%
    \usepackage{parskip}
  }{% else
    \setlength{\parindent}{0pt}
    \setlength{\parskip}{6pt plus 2pt minus 1pt}}
}{% if KOMA class
  \KOMAoptions{parskip=half}}
\makeatother
\usepackage{color}
\usepackage{fancyvrb}
\newcommand{\VerbBar}{|}
\newcommand{\VERB}{\Verb[commandchars=\\\{\}]}
\DefineVerbatimEnvironment{Highlighting}{Verbatim}{commandchars=\\\{\}}
% Add ',fontsize=\small' for more characters per line
\newenvironment{Shaded}{}{}
\newcommand{\AlertTok}[1]{\textcolor[rgb]{1.00,0.00,0.00}{\textbf{#1}}}
\newcommand{\AnnotationTok}[1]{\textcolor[rgb]{0.38,0.63,0.69}{\textbf{\textit{#1}}}}
\newcommand{\AttributeTok}[1]{\textcolor[rgb]{0.49,0.56,0.16}{#1}}
\newcommand{\BaseNTok}[1]{\textcolor[rgb]{0.25,0.63,0.44}{#1}}
\newcommand{\BuiltInTok}[1]{\textcolor[rgb]{0.00,0.50,0.00}{#1}}
\newcommand{\CharTok}[1]{\textcolor[rgb]{0.25,0.44,0.63}{#1}}
\newcommand{\CommentTok}[1]{\textcolor[rgb]{0.38,0.63,0.69}{\textit{#1}}}
\newcommand{\CommentVarTok}[1]{\textcolor[rgb]{0.38,0.63,0.69}{\textbf{\textit{#1}}}}
\newcommand{\ConstantTok}[1]{\textcolor[rgb]{0.53,0.00,0.00}{#1}}
\newcommand{\ControlFlowTok}[1]{\textcolor[rgb]{0.00,0.44,0.13}{\textbf{#1}}}
\newcommand{\DataTypeTok}[1]{\textcolor[rgb]{0.56,0.13,0.00}{#1}}
\newcommand{\DecValTok}[1]{\textcolor[rgb]{0.25,0.63,0.44}{#1}}
\newcommand{\DocumentationTok}[1]{\textcolor[rgb]{0.73,0.13,0.13}{\textit{#1}}}
\newcommand{\ErrorTok}[1]{\textcolor[rgb]{1.00,0.00,0.00}{\textbf{#1}}}
\newcommand{\ExtensionTok}[1]{#1}
\newcommand{\FloatTok}[1]{\textcolor[rgb]{0.25,0.63,0.44}{#1}}
\newcommand{\FunctionTok}[1]{\textcolor[rgb]{0.02,0.16,0.49}{#1}}
\newcommand{\ImportTok}[1]{\textcolor[rgb]{0.00,0.50,0.00}{\textbf{#1}}}
\newcommand{\InformationTok}[1]{\textcolor[rgb]{0.38,0.63,0.69}{\textbf{\textit{#1}}}}
\newcommand{\KeywordTok}[1]{\textcolor[rgb]{0.00,0.44,0.13}{\textbf{#1}}}
\newcommand{\NormalTok}[1]{#1}
\newcommand{\OperatorTok}[1]{\textcolor[rgb]{0.40,0.40,0.40}{#1}}
\newcommand{\OtherTok}[1]{\textcolor[rgb]{0.00,0.44,0.13}{#1}}
\newcommand{\PreprocessorTok}[1]{\textcolor[rgb]{0.74,0.48,0.00}{#1}}
\newcommand{\RegionMarkerTok}[1]{#1}
\newcommand{\SpecialCharTok}[1]{\textcolor[rgb]{0.25,0.44,0.63}{#1}}
\newcommand{\SpecialStringTok}[1]{\textcolor[rgb]{0.73,0.40,0.53}{#1}}
\newcommand{\StringTok}[1]{\textcolor[rgb]{0.25,0.44,0.63}{#1}}
\newcommand{\VariableTok}[1]{\textcolor[rgb]{0.10,0.09,0.49}{#1}}
\newcommand{\VerbatimStringTok}[1]{\textcolor[rgb]{0.25,0.44,0.63}{#1}}
\newcommand{\WarningTok}[1]{\textcolor[rgb]{0.38,0.63,0.69}{\textbf{\textit{#1}}}}
\usepackage{longtable,booktabs,array}
\usepackage{caption}
\captionsetup[table]{skip=6pt}
\newcounter{none} % for unnumbered tables
\usepackage{calc} % for calculating minipage widths
% Correct order of tables after \paragraph or \subparagraph
\usepackage{etoolbox}
\makeatletter
\patchcmd\longtable{\par}{\if@noskipsec\mbox{}\fi\par}{}{}
\makeatother
% Allow footnotes in longtable head/foot
\IfFileExists{footnotehyper.sty}{\usepackage{footnotehyper}}{\usepackage{footnote}}
\makesavenoteenv{longtable}
\ifLuaTeX
\usepackage[bidi=basic,shorthands=off,provide=*]{babel}
\else
\usepackage[bidi=default,shorthands=off,provide=*]{babel}
\fi
\ifLuaTeX
  \usepackage{selnolig} % disable illegal ligatures
\fi
\setlength{\emergencystretch}{3em} % prevent overfull lines
\providecommand{\tightlist}{%
  \setlength{\itemsep}{0pt}\setlength{\parskip}{0pt}}
\usepackage{bookmark}
\IfFileExists{xurl.sty}{\usepackage{xurl}}{} % add URL line breaks if available
\urlstyle{same}
% fallback for those not using the hyperref driver hyperxmp:
\makeatletter
\@ifundefined{xmpquote}{\newcommand{\xmpquote}[1]{#1}}{}
\makeatother
\hypersetup{
  pdftitle={LAPORAN PENGEMBANGAN SISTEM INFORMASI PERPUSTAKAAN DIGITAL DESA (PUSTAKA PANGKALAN)},
  pdfauthor={Tim Pelaksana Kuliah Kerja Nyata (KKN) Tematik Desa Pangkalan},
  pdflang={id-ID},
  hidelinks,
  pdfcreator={LaTeX via pandoc}}

\title{LAPORAN PENGEMBANGAN SISTEM INFORMASI PERPUSTAKAAN DIGITAL DESA
(PUSTAKA PANGKALAN)}
\usepackage{etoolbox}
\makeatletter
\providecommand{\subtitle}[1]{% add subtitle to \maketitle
  \apptocmd{\@title}{\par {\large #1 \par}}{}{}
}
\makeatother
\subtitle{Implementasi Portal Literasi Digital dan Manajemen Sirkulasi
Buku Balai Desa Berbasis Jamstack, Next.js, dan Neon Serverless
PostgreSQL}
\author{Tim Pelaksana Kuliah Kerja Nyata (KKN) Tematik Desa Pangkalan}
\date{September 2026}

\begin{document}
\maketitle

{
\setcounter{tocdepth}{3}
\tableofcontents
}
\setstretch{1.15}
\newpage

\chapter*{LEMBAR PENGESAHAN}\label{lembar-pengesahan}
\addcontentsline{toc}{chapter}{LEMBAR PENGESAHAN}

\textbf{Judul Laporan}: Laporan Pengembangan Sistem Informasi
Perpustakaan Digital Desa (Pustaka Pangkalan)\\
\textbf{Bidang Fokus}: Sistem Informasi Pemerintahan Desa dan Literasi
Digital Pedesaan\\
\textbf{Lokasi Program}: Desa Pangkalan, Kecamatan Cikidang, Kabupaten
Sukabumi, Provinsi Jawa Barat\\
\textbf{Alamat Domain Akses}: https://perpus-pangkalan.vercel.app\\
\textbf{Repositori Kode Sumber}:
https://github.com/ZephyrGraphic/pustaka-pangkalan.git\\
\textbf{Waktu Pelaksanaan}: Periode Kuliah Kerja Nyata (KKN) Tematik
Semester Ganjil 2026

Naskah laporan hasil rancang bangun sistem informasi, pengujian
perangkat lunak, dan serah terima alih kelola teknologi ini telah
diperiksa, diujicobakan, dan disahkan oleh pihak-pihak terkait pada
tanggal yang tertera di bawah ini:

\vspace{1.2cm}

{\def\LTcaptype{none} % do not increment counter
\begin{longtable}[]{@{}
  >{\centering\arraybackslash}p{(\linewidth - 2\tabcolsep) * \real{0.5000}}
  >{\centering\arraybackslash}p{(\linewidth - 2\tabcolsep) * \real{0.5000}}@{}}
\toprule\noalign{}
\begin{minipage}[b]{\linewidth}\centering
Mengetahui, \textbf{Dosen Pembimbing Lapangan (DPL)}
\end{minipage} & \begin{minipage}[b]{\linewidth}\centering
Menyetujui, \textbf{Kepala Desa Pangkalan}
\end{minipage} \\
\midrule\noalign{}
\endhead
\bottomrule\noalign{}
\endlastfoot
\vspace{2.2cm} & \vspace{2.2cm} \\
\textbf{(
\ldots\ldots\ldots\ldots\ldots\ldots\ldots\ldots\ldots\ldots\ldots\ldots\ldots\ldots\ldots\ldots\ldots\ldots\ldots\ldots{}
)} NIP:
\ldots\ldots\ldots\ldots\ldots\ldots\ldots\ldots\ldots\ldots\ldots\ldots\ldots\ldots\ldots\ldots\ldots.
& \textbf{(
\ldots\ldots\ldots\ldots\ldots\ldots\ldots\ldots\ldots\ldots\ldots\ldots\ldots\ldots\ldots\ldots\ldots\ldots\ldots\ldots{}
)} NIP:
\ldots\ldots\ldots\ldots\ldots\ldots\ldots\ldots\ldots\ldots\ldots\ldots\ldots\ldots\ldots\ldots\ldots. \\
\end{longtable}
}

\vspace{1cm}

{\def\LTcaptype{none} % do not increment counter
\begin{longtable}[]{@{}
  >{\centering\arraybackslash}p{(\linewidth - 0\tabcolsep) * \real{1.0000}}@{}}
\toprule\noalign{}
\begin{minipage}[b]{\linewidth}\centering
Mengesahkan, \textbf{Koordinator Tim Mahasiswa KKN}
\end{minipage} \\
\midrule\noalign{}
\endhead
\bottomrule\noalign{}
\endlastfoot
\vspace{2.2cm} \\
\textbf{(
\ldots\ldots\ldots\ldots\ldots\ldots\ldots\ldots\ldots\ldots\ldots\ldots\ldots\ldots\ldots\ldots\ldots\ldots\ldots\ldots{}
)} NIM:
\ldots\ldots\ldots\ldots\ldots\ldots\ldots\ldots\ldots\ldots\ldots\ldots\ldots\ldots\ldots\ldots\ldots. \\
\end{longtable}
}

\newpage

\chapter*{PERNYATAAN KEASLIAN KARYA}\label{pernyataan-keaslian-karya}
\addcontentsline{toc}{chapter}{PERNYATAAN KEASLIAN KARYA}

Kami yang bertanda tangan di bawah ini, Tim Pelaksana Mahasiswa Kuliah
Kerja Nyata (KKN) Tematik Desa Pangkalan, dengan ini menyatakan secara
sadar dan sungguh-sungguh bahwa:

\begin{enumerate}
\def\labelenumi{\arabic{enumi}.}
\tightlist
\item
  Naskah laporan sistem informasi ini beserta seluruh artefak perangkat
  lunak yang dihasilkannya (kode sumber, skema basis data, antarmuka
  pengguna, dan naskah dokumentasi) merupakan karya asli hasil kerja
  mandiri tim kami di bawah bimbingan Dosen Pembimbing Lapangan dan
  arahan Pemerintah Desa Pangkalan, Kecamatan Cikidang, Kabupaten
  Sukabumi.
\item
  Seluruh sumber pustaka, kutipan teori, data lapangan, dan pustaka
  perangkat lunak \emph{open-source} yang digunakan telah dicantumkan
  dan dirujuk secara sah sesuai dengan kaidah penulisan karya ilmiah
  yang berlaku.
\item
  Seluruh fitur, arsitektur, dan metrik yang dilaporkan dalam dokumen
  ini bersumber langsung dari sistem terpasang aktual (\emph{as-built
  system}) tanpa memuat fitur rekaan (\emph{zero fake features}) serta
  sepenuhnya tunduk pada kebijakan perlindungan data rahasia (\emph{zero
  secrets policy}).
\end{enumerate}

Apabila di kemudian hari terbukti terdapat plagiarisme, fabrikasi data,
atau pelanggaran etika akademik dalam naskah ini, kami bersedia menerima
sanksi akademis dan administratif sesuai dengan peraturan
perundang-undangan yang berlaku.

\vspace{1.2cm}

\emph{Desa Pangkalan, September 2026}\\
\vspace{1.2cm} \textbf{( Tim Mahasiswa Pengembang Pustaka Pangkalan )}

\newpage

\chapter*{ABSTRAK}\label{abstrak}
\addcontentsline{toc}{chapter}{ABSTRAK}

\textbf{PENGEMBANGAN SISTEM INFORMASI PERPUSTAKAAN DIGITAL DESA (PUSTAKA
PANGKALAN) BERBASIS JAMSTACK DAN SERVERLESS POSTGRESQL DI DESA PANGKALAN
KECAMATAN CIKIDANG KABUPATEN SUKABUMI}

\emph{Oleh: Tim Pelaksana KKN Tematik Desa Pangkalan}

Keterbatasan ketersediaan bahan pustaka cetak dan pengelolaan sirkulasi
peminjaman buku yang masih manual di kantor Balai Desa Pangkalan,
Kecamatan Cikidang, Kabupaten Sukabumi, menjadi kendala utama dalam
mendorong peningkatan literasi masyarakat pedesaan yang tersebar di
empat dusun (Pangkalan, Cikajang, Pasir Arangan, dan Pasir Gombong).
Penelitian dan proyek perancangan perangkat lunak ini bertujuan
membangun \textbf{Pustaka Pangkalan}, sebuah sistem informasi
perpustakaan digital desa berbasis web yang memadukan portal modul
e-book interaktif per bab (\emph{chapter-by-chapter reader}) dengan
sistem otomasi sirkulasi peminjaman buku fisik balai desa. Sistem
dirancang menggunakan arsitektur \emph{Jamstack} modern berbasis
kerangka kerja Next.js 16 (React 19 dan TypeScript), basis data
relasional Neon Serverless PostgreSQL yang dikelola melalui Prisma ORM
7, serta diamankan melalui otentikasi NIK 16 digit dan PIN 6 digit
terenkripsi Bcrypt.

Metodologi pengembangan sistem yang diterapkan adalah \emph{Agile
Iterative Prototyping}, yang meliputi tahapan inisiasi lapangan,
analisis kebutuhan pengguna, perancangan sistem menggunakan pemodelan
UML (\emph{Use Case Diagram, Activity Diagram}) dan ERD, konstruksi
perangkat lunak, serta pengujian mutu komprehensif. Evaluasi sistem
dilakukan melalui \emph{Software Testing and Quality Assurance} (STQA)
berbasis pengujian \emph{black-box} dan \emph{regression testing} yang
mencakup 9 test suite dan 48 skenario uji fungsional. Hasil pengujian
menunjukkan tingkat kelulusan 100\% (\emph{zero defect}) dengan
responsivitas tinggi pada perangkat bergerak (\emph{mobile-first}).
Selain itu, sistem mengintegrasikan fitur gamifikasi literasi (poin baca
dan lencana prestasi), asisten cerdas \emph{Kades AI}, serta lokalisasi
dwi-bahasa (Bahasa Indonesia dan Basa Sunda) untuk melestarikan kearifan
lokal. Sistem ini telah sukses dideploy pada domain publik
https://perpus-pangkalan.vercel.app dan siap dioperasikan secara mandiri
oleh aparatur Balai Desa Pangkalan.

\textbf{Kata Kunci}: Sistem Informasi Perpustakaan, Perpustakaan Digital
Desa, Next.js, Serverless PostgreSQL, Jamstack, Literasi Pedesaan, Basa
Sunda, Kecamatan Cikidang.

\newpage

\chapter*{ABSTRACT}\label{abstract}
\addcontentsline{toc}{chapter}{ABSTRACT}

\textbf{DEVELOPMENT OF VILLAGE DIGITAL LIBRARY INFORMATION SYSTEM
(PUSTAKA PANGKALAN) BASED ON JAMSTACK AND SERVERLESS POSTGRESQL IN
PANGKALAN VILLAGE, CIKIDANG DISTRICT, SUKABUMI REGENCY}

\emph{By: Community Service Student Team of Pangkalan Village}

The acute shortage of printed reading resources and the conventional
paper-based circulation management at the Village Hall of Pangkalan,
Cikidang District, Sukabumi Regency, have hindered the advancement of
rural literacy across four distinct hamlets (Pangkalan, Cikajang, Pasir
Arangan, and Pasir Gombong). This research and development project
engineered \textbf{Pustaka Pangkalan}, a robust web-based village
digital library information system that synergizes a lightweight,
chapter-by-chapter interactive e-book reader with an automated physical
book circulation management system for village administrators. The
application was constructed leveraging modern Jamstack architecture
powered by Next.js 16 (React 19 and TypeScript), Neon Serverless
PostgreSQL interfaced via Prisma ORM 7, and secured utilizing 16-digit
National Identity Number (NIK) paired with Bcrypt-encrypted 6-digit PIN
authentication.

The system development methodology utilized an Agile Iterative
Prototyping framework, progressing systematically through field
observation, requirements elicitation, object-oriented UML modeling
(\emph{Use Case and Activity Diagrams}), relational database schema
formulation (ERD), full-stack implementation, and rigorous verification.
Software Testing and Quality Assurance (STQA) was carried out through
structured black-box and automated regression suites encompassing 9 test
suites and 48 distinct test cases. The verification yielded a 100\% pass
rate (zero defect), confirming resilient execution across varying
network bandwidths and mobile form factors. Furthermore, the platform
integrates gamified literacy incentives (reading points, streaks, and
milestone badges), an interactive localized conversational assistant
(\emph{Kades AI}), and dual-language localization (Indonesian and
Sundanese). The platform is officially deployed at
https://perpus-pangkalan.vercel.app and transitioned for long-term
municipal governance.

\textbf{Keywords}: Library Information System, Village Digital Library,
Next.js, Serverless PostgreSQL, Jamstack, Rural Literacy, Sundanese
Language, Cikidang District.

\newpage

\chapter*{KATA PENGANTAR}\label{kata-pengantar}
\addcontentsline{toc}{chapter}{KATA PENGANTAR}

Puji dan syukur senantiasa kami panjatkan ke hadirat Allah Subhanahu Wa
Ta'ala, Tuhan Yang Maha Kuasa, atas segala rahmat, hidayah, serta
pertolongan-Nya sehingga perancangan, pengembangan, pengujian, dan
penyusunan naskah \textbf{Laporan Pengembangan Sistem Informasi
Perpustakaan Digital Desa (Pustaka Pangkalan)} di Desa Pangkalan,
Kecamatan Cikidang, Kabupaten Sukabumi, Jawa Barat ini dapat
diselesaikan dengan tuntas dan paripurna.

Laporan ini disusun sebagai dokumen pertanggungjawaban akademis, teknis,
dan institusional atas pelaksanaan program Kuliah Kerja Nyata (KKN)
Tematik. Fokus program ini adalah menerapkan teknologi informasi modern
yang tepat guna untuk mengatasi permasalahan riil kesenjangan literasi
dan otomasi administrasi pemerintahan desa. Kehadiran Sistem Informasi
Pustaka Pangkalan diharapkan mampu menjembatani kendala geografis di
empat dusun terpisah, sehingga warga, pelajar, dan kader desa dapat
dengan mudah mengakses modul-modul ilmu terapan seperti pertanian
organik, budidaya perikanan, gizi keluarga pencegah stunting,
kewirausahaan UMKM, serta pelestarian kebudayaan dan sastra Sunda.

Keberhasilan penyusunan laporan dan pembangunan perangkat lunak ini
tidak terlepas dari bimbingan, kontribusi, dan dukungan yang tulus dari
berbagai pihak. Oleh karena itu, kami menyampaikan terima kasih dan
penghargaan yang mendalam kepada:

\begin{enumerate}
\def\labelenumi{\arabic{enumi}.}
\tightlist
\item
  \textbf{Dosen Pembimbing Lapangan (DPL)}, yang dengan penuh kesabaran
  memberikan arahan akademis, telaah metodologi, kritik konstruktif,
  serta bimbingan moral selama pelaksanaan KKN.
\item
  \textbf{Pemerintah Desa Pangkalan, Kecamatan Cikidang}, khususnya
  Bapak Kepala Desa beserta seluruh perangkat balai desa dan para Kepala
  Dusun (Dusun Pangkalan, Dusun Cikajang, Dusun Pasir Arangan, dan Dusun
  Pasir Gombong) atas sambutan hangat, fasilitas observasi, dan kerja
  sama kemitraan yang terjalin erat.
\item
  \textbf{Pengurus Karang Taruna dan Tokoh Pemuda Desa Pangkalan}, atas
  partisipasi aktif dalam pengujian fungsional aplikasi serta
  antusiasmenya menjadi penggerak budaya literasi digital di tingkat
  dusun.
\item
  \textbf{Masyarakat dan Pelajar Desa Pangkalan}, atas kesediaan
  meluangkan waktu dalam sesi pengumpulan data, pengisian survei
  kebutuhan, dan uji coba sistem.
\item
  \textbf{Rekan-rekan Mahasiswa Tim Pelaksana KKN}, yang telah
  memperlihatkan dedikasi, kekompakan, dan etos kerja yang tinggi dalam
  setiap tahapan analisis, perancangan, koding, hingga penyerahan
  sistem.
\end{enumerate}

Kami menyadari bahwa laporan ini masih memiliki ruang untuk perbaikan.
Oleh karena itu, segala bentuk kritik dan saran yang konstruktif sangat
kami harapkan guna penyempurnaan sistem di masa yang akan datang. Semoga
karya ini memberikan manfaat yang nyata dan berkelanjutan bagi
masyarakat Desa Pangkalan.

\vspace{1.2cm}

\emph{Desa Pangkalan, September 2026}\\
\vspace{0.8cm} \textbf{Tim Penyusun KKN Mahasiswa}

\newpage

\chapter*{DAFTAR SINGKATAN DAN
GLOSARIUM}\label{daftar-singkatan-dan-glosarium}
\addcontentsline{toc}{chapter}{DAFTAR SINGKATAN DAN GLOSARIUM}

{\def\LTcaptype{none} % do not increment counter
\begin{longtable}[]{@{}
  >{\raggedright\arraybackslash}p{(\linewidth - 2\tabcolsep) * \real{0.5000}}
  >{\raggedright\arraybackslash}p{(\linewidth - 2\tabcolsep) * \real{0.5000}}@{}}
\toprule\noalign{}
\begin{minipage}[b]{\linewidth}\raggedright
Singkatan / Istilah
\end{minipage} & \begin{minipage}[b]{\linewidth}\raggedright
Kepanjangan / Definisi Resmi
\end{minipage} \\
\midrule\noalign{}
\endhead
\bottomrule\noalign{}
\endlastfoot
\textbf{API} & \emph{Application Programming Interface}, antarmuka
protokol yang memungkinkan komunikasi pertukaran data antar-komponen
perangkat lunak secara terstandar. \\
\textbf{Bcrypt} & Fungsi \emph{hashing} kata sandi berbasis algoritma
\emph{Blowfish} yang dilengkapi \emph{salt} dan faktor biaya kerja
adaptif untuk mencegah serangan \emph{brute-force} dan \emph{rainbow
table}. \\
\textbf{CRUD} & \emph{Create, Read, Update, Delete}, empat operasi
fundamental dalam manipulasi data pada sistem manajemen basis data. \\
\textbf{DDC} & \emph{Dewey Decimal Classification}, sistem klasifikasi
perpustakaan internasional terstruktur berbasis hierarki numerik 10
kelas utama. \\
\textbf{DPL} & Dosen Pembimbing Lapangan, dosen perguruan tinggi yang
bertindak sebagai pembina dan penilai akademis kegiatan KKN
mahasiswa. \\
\textbf{ERD} & \emph{Entity-Relationship Diagram}, model grafis
konseptual yang menggambarkan entitas, atribut data, dan kardinalitas
relasi dalam basis data. \\
\textbf{GUI / UI} & \emph{Graphical User Interface}, antarmuka visual
grafis yang memfasilitasi interaksi pengguna dengan perangkat lunak
melalui elemen visual interaktif. \\
\textbf{Jamstack} & \emph{JavaScript, APIs, and Markup}, arsitektur
rekayasa web modern yang memisahkan lapisan penyajian (\emph{frontend})
dari lapisan logika data (\emph{backend}) guna mencapai performa dan
keamanan tinggi. \\
\textbf{JWT} & \emph{JSON Web Token}, standar terbuka (RFC 7519) yang
mendefinisikan cara ringkas dan mandiri untuk mentransmisikan informasi
sesi otentikasi antar pihak secara aman. \\
\textbf{KKN} & Kuliah Kerja Nyata, bentuk kegiatan pengabdian kepada
masyarakat yang wajib ditempuh oleh mahasiswa perguruan tinggi. \\
\textbf{NIK} & Nomor Induk Kependudukan, nomor identitas kependudukan
tunggal yang bersifat unik dan permanen sepanjang 16 digit angka di
Indonesia. \\
\textbf{ORM} & \emph{Object-Relational Mapping}, teknik pemetaan data
dalam rekayasa perangkat lunak yang mengonversi tipe data antar sistem
relasional dan bahasa berorientasi objek. \\
\textbf{Prisma} & \emph{Next-generation Object-Relational Mapper} untuk
TypeScript dan Node.js yang menyediakan skema deklaratif dan kueri data
berkeamanan tipe ketat (\emph{type-safe}). \\
\textbf{RBAC} & \emph{Role-Based Access Control}, metode kontrol akses
keamanan komputer yang membatasi hak akses sistem hanya kepada pengguna
berdasarkan peran wewenang resminya. \\
\textbf{REST} & \emph{Representational State Transfer}, gaya arsitektur
perangkat lunak untuk sistem terdistribusi yang memanfaatkan protokol
komunikasi HTTP tanpa status (\emph{stateless}). \\
\textbf{SDLC} & \emph{Software Development Life Cycle}, metodologi
terstruktur dalam rekayasa perangkat lunak untuk merencanakan,
menganalisis, merancang, membangun, menguji, dan memelihara sistem
informasi. \\
\textbf{STQA} & \emph{Software Testing and Quality Assurance}, disiplin
pengujian perangkat lunak terencana dan sistematis guna menjamin
kepatuhan sistem terhadap spesifikasi mutu yang telah ditetapkan. \\
\textbf{UML} & \emph{Unified Modeling Language}, bahasa visual standar
pemodelan perangkat lunak berorientasi objek yang mencakup diagram
struktural dan perilaku (\emph{behavioral}). \\
\textbf{Vercel} & Platform \emph{cloud deployment} dan komputasi
\emph{serverless edge} terdistribusi global yang dioptimalkan untuk
kerangka kerja web Next.js. \\
\end{longtable}
}

\newpage

\chapter{BAB I --- PENDAHULUAN}\label{bab-i-pendahuluan}

\section{1.1 Latar Belakang Masalah}\label{latar-belakang-masalah}

Pembangunan perdesaan merupakan fondasi penting dalam menopang ketahanan
nasional, stabilitas pangan, dan pelestarian nilai-nilai sosial budaya
bangsa Indonesia. Transformasi digital yang diamanatkan oleh pemerintah
pusat melalui konsep Desa Cerdas (\emph{Smart Village}) menuntut
pemanfaatan teknologi informasi dan komunikasi untuk meningkatkan
kualitas hidup, efisiensi pelayanan publik, dan keterbukaan akses ilmu
pengetahuan di tingkat desa.

Desa Pangkalan merupakan salah satu desa di wilayah administrasi
Kecamatan Cikidang, Kabupaten Sukabumi, Provinsi Jawa Barat. Secara
geografis, wilayah Desa Pangkalan memiliki karakteristik topografi
perbukitan dan agraris dengan permukiman penduduk yang tersebar di 4
(empat) wilayah dusun utama, yaitu: 1. \textbf{Dusun Pangkalan} (wilayah
pusat administrasi balai desa dan permukiman sentral); 2. \textbf{Dusun
Cikajang} (wilayah agraris dengan konsentrasi perkebunan dan pertanian
hortikultura); 3. \textbf{Dusun Pasir Arangan} (wilayah perbukitan
dengan kegiatan peternakan rakyat dan usaha pertanian lahan kering); 4.
\textbf{Dusun Pasir Gombong} (wilayah perbatasan dusun dengan potensi
kerajinan, anyaman bambu, dan perkebunan).

Sebagian besar mata pencaharian masyarakat Desa Pangkalan bertumpu pada
sektor pertanian tanaman pangan, perkebunan karet dan kelapa, peternakan
domba, serta usaha mikro, kecil, dan menengah (UMKM) pengolahan hasil
bumi. Dalam era keterbukaan informasi saat ini, penetrasi gawai pintar
(\emph{smartphone}) dan ketersediaan sinyal telekomunikasi seluler telah
menjangkau sebagian besar permukiman warga di keempat dusun tersebut.
Namun demikian, ketersediaan infrastruktur jaringan ini belum diimbangi
dengan ketersediaan konten edukasi dan sumber bacaan terapan yang
relevan dengan kebutuhan hidup masyarakat desa. Akses internet oleh
warga dan pemuda desa sebagian besar masih dihabiskan untuk hiburan
media sosial tanpa memberikan dampak peningkatan kapasitas ekonomi dan
keterampilan terapan.

Di sisi lain, kantor Balai Desa Pangkalan memiliki fasilitas
perpustakaan desa yang menyimpan sejumlah koleksi buku fisik. Namun,
pengelolaan perpustakaan balai desa tersebut selama ini menghadapi
berbagai kendala nyata: 1. \textbf{Keterbatasan Fisik dan Geografis}:
Warga yang berdomisili di dusun-dusun yang jauh dari kantor balai desa
harus menempuh perjalanan fisik yang cukup jauh hanya untuk memeriksa
ketersediaan buku atau meminjam buku bacaan. 2. \textbf{Pengelolaan
Sirkulasi Manual}: Pencatatan riwayat peminjaman dan pengembalian buku
masih menggunakan buku besar manual berbasis kertas (\emph{paper-based
record}). Hal ini menyebabkan tingginya risiko kehilangan catatan,
sulitnya memantau buku yang telah melampaui batas waktu pinjam
(\emph{overdue}), serta ketidakjelasan status inventaris eksemplar buku
di balai desa. 3. \textbf{Ketiadaan Modul Pembelajaran Terapan Digital}:
Koleksi buku cetak yang ada di balai desa didominasi oleh buku-buku umum
terbitan lama yang kurang relevan dengan isu-isu mendesak masyarakat
pedesaan saat ini, seperti: teknik pembuatan pupuk organik padat dan
cair, budidaya ikan air tawar sistem bioflok, panduan gizi seimbang
balita untuk pencegahan \emph{stunting}, pembukuan keuangan digital UMKM
desa, hingga materi pembelajaran bahasa dan aksara Sunda sebagai warisan
budaya lokal.

Menjawab permasalahan tersebut, Tim Mahasiswa Kuliah Kerja Nyata (KKN)
Tematik berinisiatif merancang dan mengimplementasikan sebuah solusi
rekayasa perangkat lunak terpadu yang diberi nama \textbf{Sistem
Informasi Perpustakaan Digital Desa (Pustaka Pangkalan)}. Sistem ini
menggabungkan portal baca digital (\emph{digital e-book reader}) yang
ringan diakses melalui peramban ponsel pintar warga dengan modul otomasi
manajemen sirkulasi peminjaman buku fisik yang dapat dioperasikan secara
tertib oleh aparatur Balai Desa Pangkalan.

\section{1.2 Perumusan Masalah}\label{perumusan-masalah}

Berdasarkan latar belakang masalah yang telah diuraikan, rumusan masalah
dalam penelitian dan pengembangan sistem informasi ini adalah: 1.
Bagaimana menganalisis dan memetakan kebutuhan sistem informasi
perpustakaan yang mampu melayani pembacaan materi digital per bab bagi
warga sekaligus memodernisasi tata kelola peminjaman buku fisik di Balai
Desa Pangkalan? 2. Bagaimana merancang arsitektur perangkat lunak
berbasis web modern (\emph{Jamstack, Next.js, dan Serverless
PostgreSQL}) yang berkinerja tinggi, aman, hemat sumber daya komputasi,
dan responsif terhadap perangkat bergerak (\emph{mobile-first})? 3.
Bagaimana mengimplementasikan mekanisme otentikasi yang ramah pengguna
pedesaan menggunakan kombinasi NIK 16 digit dan PIN 6 digit terenkripsi
tanpa mengorbankan standar keamanan data? 4. Bagaimana menguji dan
mengevaluasi keandalan, integritas skema basis data, dan fungsionalitas
sistem informasi menggunakan metode pengujian \emph{black-box} dan
\emph{regression testing} berstandar kualitas perangkat lunak?

\section{1.3 Batasan Masalah}\label{batasan-masalah}

Agar pembahasan dalam laporan ini terarah dan terfokus pada sistem yang
benar-benar terpasang aktual (\emph{as-built system}), batasan masalah
ditetapkan sebagai berikut: 1. Sistem dikembangkan sebagai aplikasi
berbasis web responsif (\emph{web-based responsive application}) yang
dideploy pada infrastruktur \emph{Vercel Edge Network} dengan alamat
domain resmi \textbf{https://perpus-pangkalan.vercel.app}. 2. Lingkungan
basis data menggunakan layanan \emph{relational database} berbasis awan
\textbf{Neon Serverless PostgreSQL} yang dikelola menggunakan
\textbf{Prisma ORM versi 7.9.1}. 3. Kerangka kerja perangkat lunak yang
digunakan adalah \textbf{Next.js 16.3.0} (berbasis React 19 dan
TypeScript) dengan pola arsitektur \emph{App Router} dan \emph{Route
Handlers}. 4. Aktor sistem dibatasi pada 2 (dua) peran pengguna
terotentikasi, yaitu Warga Desa (\texttt{USER}) dan Pengelola Balai Desa
(\texttt{ADMIN}), serta 1 (satu) mode pengunjung tamu (\texttt{GUEST}).
5. Otentikasi pengguna memanfaatkan \emph{NextAuth.js} dengan strategi
kredensial NIK (16 digit angka) dan PIN (6 digit angka) yang dienkripsi
menggunakan algoritma \emph{Bcrypt} dengan \emph{salt rounds} 10. 6.
Fitur pembaca digital (\emph{reader}) difokuskan pada penyajian teks
materi e-book per bab (\emph{chapter reader}) yang ringan dan cepat
dimuat pada jaringan seluler desa, dilengkapi penyimpanan riwayat bacaan
dan penanda halaman (\emph{bookmark}). 7. Pengelolaan sirkulasi buku
fisik balai desa mencakup pencatatan tanggal pinjam, batas kembali
otomatis 7 hari, fasilitas perpanjangan (+7 hari), penandaan
pengembalian instan, dan ekspor laporan berkala dalam format
\emph{Comma-Separated Values} (CSV). 8. Sistem mengintegrasikan
fasilitas lokalisasi dwi-bahasa (Bahasa Indonesia dan Basa Sunda) serta
antarmuka asisten virtual cerdas \emph{Kades AI} berbasis \emph{Large
Language Model} terkonfigurasi.

\section{1.4 Tujuan Penelitian dan
Pengembangan}\label{tujuan-penelitian-dan-pengembangan}

Tujuan dari pelaksanaan proyek perancangan dan implementasi sistem
informasi ini meliputi:

\subsection{1.4.1 Tujuan Umum}\label{tujuan-umum}

Menghasilkan sebuah sistem informasi perpustakaan digital desa berbasis
web yang andal, aman, dan mudah digunakan guna mendemokratisasi akses
sumber ilmu pengetahuan terapan bagi masyarakat serta memodernisasi
manajemen sirkulasi perpustakaan Balai Desa Pangkalan, Kecamatan
Cikidang, Kabupaten Sukabumi.

\subsection{1.4.2 Tujuan Khusus}\label{tujuan-khusus}

\begin{enumerate}
\def\labelenumi{\arabic{enumi}.}
\tightlist
\item
  Menganalisis kondisi eksisting tata kelola perpustakaan balai desa dan
  merumuskan spesifikasi kebutuhan fungsional (FR) dan non-fungsional
  (NFR) sistem secara komprehensif.
\item
  Merancang arsitektur sistem 4-tier (\emph{Presentation, Application,
  Data Access, Persistence Layer}), pemodelan berorientasi objek
  (\emph{Use Case Diagram, Activity Diagram}), dan skema basis data
  relasional (\emph{ERD \& Data Dictionary} 10 tabel).
\item
  Mengonstruksi antarmuka pengguna responsif (19 rute halaman) dan
  antarmuka pemrograman aplikasi (30 endpoint RESTful API) yang
  menghubungkan frontend dan backend secara efisien.
\item
  Menerapkan mekanisme keamanan akun warga berbasis validasi NIK KTP dan
  enkripsi PIN Bcrypt serta kontrol akses berbasis peran
  (\emph{Role-Based Access Control}).
\item
  Menerapkan fitur insentif literasi berbasis gamifikasi (poin baca,
  \emph{reading streak}, 3 jenjang lencana prestasi, dan papan peringkat
  dusun).
\item
  Menguji fungsionalitas dan keandalan sistem secara menyeluruh melalui
  9 test suite STQA otomatis yang mencakup 48 skenario pengujian dengan
  target kelulusan 100\% (\emph{zero defect}).
\item
  Menyusun dokumen serah terima teknologi (\emph{handover}), berita
  acara serah terima resmi (BAST), dan buku panduan pengguna (\emph{user
  manual}) demi keberlanjutan operasional sistem di balai desa.
\end{enumerate}

\section{1.5 Manfaat Penelitian dan
Pengembangan}\label{manfaat-penelitian-dan-pengembangan}

Implementasi Sistem Informasi Pustaka Pangkalan memberikan manfaat nyata
bagi berbagai pihak: 1. \textbf{Bagi Masyarakat dan Pelajar Desa
Pangkalan}: - Memberikan kemudahan akses terhadap materi bacaan terapan
seputar pertanian, peternakan, kesehatan keluarga, dan wirausaha kapan
pun dan di mana pun melalui ponsel pintar. - Memberikan motivasi membaca
berkelanjutan melalui sistem penghargaan poin dan lencana keaktifan
membaca. - Menyediakan sarana pelestarian budaya lokal melalui bacaan
sastra Sunda dan antarmuka berbahasa Sunda. 2. \textbf{Bagi Pemerintah
Desa Pangkalan}: - Mewujudkan tertib administrasi inventarisasi buku
fisik balai desa dengan pencatatan sirkulasi yang terstruktur dan
termonitor secara real-time. - Memperoleh data analitik profil minat
baca masyarakat per dusun yang valid sebagai rujukan pengambilan
kebijakan program pemberdayaan masyarakat. - Meningkatkan citra positif
Desa Pangkalan sebagai desa yang adaptif terhadap inovasi teknologi
informasi (\emph{smart village}). 3. \textbf{Bagi Sivitas Akademika dan
Tim Pengembang KKN}: - Mengaplikasikan teori rekayasa perangkat lunak,
manajemen basis data, dan interaksi manusia-komputer dalam pemecahan
masalah nyata di tengah masyarakat pedesaan. - Menghasilkan artefak
karya ilmiah teruji yang memenuhi standar penjaminan mutu perangkat
lunak (\emph{STQA verified}).

\section{1.6 Metodologi Pengembangan Sistem
Ringkas}\label{metodologi-pengembangan-sistem-ringkas}

Pengembangan sistem informasi ini mengadopsi model \textbf{Agile
Iterative Prototyping} yang terbagi ke dalam 6 (enam) tahapan
berkesinambungan: 1. \textbf{Fase Inisiasi dan Pengumpulan Data}:
Observasi langsung ke kantor Balai Desa Pangkalan dan 4 dusun, wawancara
mendalam dengan Kepala Desa dan aparatur, serta studi kepustakaan. 2.
\textbf{Fase Analisis Kebutuhan}: Identifikasi profil aktor sistem,
analisis sistem berjalan (\emph{as-is}), perumusan sistem usulan
(\emph{to-be}), dan penyusunan matriks kebutuhan fungsional (FR-01 s.d.
FR-20). 3. \textbf{Fase Perancangan (\emph{System Design})}: Pemodelan
diagram alur (\emph{flowchart}), diagram UML (\emph{Use Case \& Activity
Diagram}), perancangan arsitektur \emph{Jamstack}, perancangan skema
relasional ERD, dan perancangan antarmuka pengguna (\emph{UI/UX
wireframing}). 4. \textbf{Fase Konstruksi (\emph{Implementation})}:
Penulisan kode sumber frontend dan backend menggunakan Next.js 16,
TypeScript, Tailwind CSS, Prisma ORM, dan pengintegrasian basis data
Neon PostgreSQL. 5. \textbf{Fase Pengujian (\emph{Software Testing \&
QA})}: Pengujian \emph{black-box}, verifikasi skema basis data,
pengujian regresi otomatis (48 skenario uji), dan evaluasi usabilitas
sistem. 6. \textbf{Fase Deployment dan Serah Terima}: Pemasangan sistem
pada server produksi Vercel Edge, pelatihan operator desa,
penandatanganan Berita Acara Serah Terima (BAST), dan penyerahan buku
manual.

\section{1.7 Sistematika Penulisan
Laporan}\label{sistematika-penulisan-laporan}

Naskah laporan ini disusun secara sistematis ke dalam 7 (tujuh) bab
utama: - \textbf{BAB I --- PENDAHULUAN}: Memuat latar belakang masalah
geografis dan literasi desa, rumusan masalah, batasan masalah, tujuan,
manfaat, metodologi ringkas, dan sistematika penulisan. - \textbf{BAB II
--- LANDASAN TEORI DAN TINJAUAN PUSTAKA}: Menguraikan teori-teori dasar
sistem informasi, perpustakaan digital, rekayasa perangkat lunak agile,
pemodelan UML/ERD, stack teknologi Jamstack/Serverless, standar keamanan
Bcrypt/NextAuth, standar pengujian ISO/IEC 25010, serta tabel
perbandingan penelitian terdahulu. - \textbf{BAB III --- METODOLOGI
PENELITIAN DAN PENGEMBANGAN}: Menguraikan lokasi dan jadwal pelaksanaan,
kerangka pikir penelitian (\emph{research framework}), teknik
pengumpulan data, tahapan siklus pengembangan perangkat lunak, serta
spesifikasi alat kerja. - \textbf{BAB IV --- ANALISIS DAN PERANCANGAN
SISTEM}: Menyajikan analisis sistem berjalan, perancangan sistem usulan,
matriks kebutuhan fungsional dan non-fungsional, pemodelan UML
(\emph{Use Case Narrative \& Activity Diagram}), arsitektur 4-tier,
perancangan ERD, kamus data 10 tabel, dan tata letak antarmuka. -
\textbf{BAB V --- IMPLEMENTASI SISTEM}: Menjelaskan konfigurasi
lingkungan produksi, implementasi skema basis data, realisasi 19 halaman
antarmuka pengguna, 30 endpoint RESTful API, keamanan otentikasi
NIK/PIN, fitur gamifikasi, asisten Kades AI, lokalisasi dwibahasa, dan
branding Kabupaten Sukabumi. - \textbf{BAB VI --- PENGUJIAN DAN EVALUASI
SISTEM}: Menyajikan strategi pengujian \emph{black-box}, eksekusi 9 test
suite STQA (48 skenario uji fungsional), evaluasi hasil pengujian, serta
analisis usabilitas sistem. - \textbf{BAB VII --- PENUTUP}: Menyimpulkan
pencapaian proyek menjawab rumusan masalah, mengidentifikasi
keterbatasan sistem terpasang, dan memberikan rekomendasi operasional
bagi pemerintah desa. - \textbf{BAGIAN AKHIR}: Berisi Daftar Pustaka
berstandar resmi, Lampiran 1 (Ringkasan Buku Panduan Pengguna), Lampiran
2 (Naskah Resmi BAST), dan Lampiran 3 (Laporan Hasil Audit Konsistensi
Sistem).

\newpage

\chapter{BAB II --- LANDASAN TEORI DAN KAJIAN
PUSTAKA}\label{bab-ii-landasan-teori-dan-kajian-pustaka}

\section{2.1 Konsep Dasar Sistem dan Sistem
Informasi}\label{konsep-dasar-sistem-dan-sistem-informasi}

Secara terminologis, konsep sistem dan informasi telah didefinisikan
secara luas oleh para pakar rekayasa sistem informasi:

\subsection{2.1.1 Definisi Sistem}\label{definisi-sistem}

Menurut Gordon B. Davis (1985), sistem adalah sekumpulan elemen yang
beroperasi secara bersama-sama untuk menyelesaikan suatu sasaran
tertentu. Elemen-elemen ini saling berinteraksi, terhubung oleh suatu
keteraturan, dan bekerja dalam kesatuan batasan lingkungan guna
mentransformasikan masukan (\emph{input}) menjadi keluaran
(\emph{output}). Menurut James A. O'Brien (2010), sebuah sistem dapat
dipahami sebagai sekelompok komponen yang saling berhubungan, bekerja
bersama menuju tujuan bersama dengan menerima masukan dan menghasilkan
keluaran dalam proses transformasi yang terorganisir.

Suatu sistem yang baik dicirikan oleh karakteristik fundamental: 1.
Memiliki komponen atau elemen (\emph{components}); 2. Memiliki batas
sistem (\emph{boundary}) yang memisahkan sistem dari lingkungan luarnya;
3. Memiliki lingkungan luar (\emph{environment}) yang memengaruhi
operasi sistem; 4. Memiliki media penghubung antarmuka
(\emph{interface}) antar-komponen; 5. Memiliki masukan (\emph{input}),
pengolah (\emph{process}), keluaran (\emph{output}), dan mekanisme umpan
balik (\emph{feedback}).

\subsection{2.1.2 Definisi Informasi dan Siklus
Informasi}\label{definisi-informasi-dan-siklus-informasi}

Informasi merupakan data yang telah diproses, diorganisasikan, dan
distrukturkan sedemikian rupa sehingga memiliki nilai guna, konteks, dan
arti bagi pengambil keputusan. Menurut Jogiyanto H.M. (2005), kualitas
suatu informasi ditentukan oleh 3 (tiga) pilar utama: - \textbf{Akurat
(\emph{Accurate})}: Informasi harus bebas dari kesalahan, tidak bias,
dan jelas mencerminkan keadaan sebenarnya. - \textbf{Tepat Waktu
(\emph{Timeliness})}: Informasi harus tiba kepada penerima saat
diperlukan; informasi yang usang tidak memiliki nilai strategis. -
\textbf{Relevan (\emph{Relevance})}: Informasi harus memiliki
keterkaitan langsung dengan kebutuhan pemakainya.

\subsection{2.1.3 Definisi Sistem
Informasi}\label{definisi-sistem-informasi}

Sistem informasi adalah kombinasi terpadu antara orang (\emph{people}),
perangkat keras (\emph{hardware}), perangkat lunak (\emph{software}),
jaringan komunikasi (\emph{networks}), dan sumber daya data (\emph{data
resources}) yang mengumpulkan, mentransformasikan, dan menyebarkan
informasi dalam suatu organisasi guna mendukung pengambilan keputusan,
koordinasi, analisis, dan visualisasi operasional (O'Brien \& Marakas,
2011).

\section{\texorpdfstring{2.2 Perpustakaan Digital (\emph{Digital
Library}) dan Otomasi
Perpustakaan}{2.2 Perpustakaan Digital (Digital Library) dan Otomasi Perpustakaan}}\label{perpustakaan-digital-digital-library-dan-otomasi-perpustakaan}

\subsection{2.2.1 Pengertian Perpustakaan
Digital}\label{pengertian-perpustakaan-digital}

Menurut William Y. Arms (2000), perpustakaan digital (\emph{digital
library}) didefinisikan sebagai kumpulan informasi yang terkelola secara
terorganisasi dan didukung oleh layanan-layanan terkait, di mana
informasinya disimpan dalam format digital dan dapat diakses melalui
jaringan komputer. Sementara itu, Michael Lesk (1997) menegaskan bahwa
perpustakaan digital merupakan koleksi informasi terorganisasi dalam
bentuk digital yang menggabungkan struktur dan pengumpulan data dengan
kemudahan penelusuran digital.

Keunggulan komparatif perpustakaan digital dibandingkan perpustakaan
konvensional meliputi: 1. \textbf{Aksesibilitas Tanpa Batas Ruang dan
Waktu (\emph{Ubiquitous Access})}: Koleksi digital dapat diakses 24 jam
sehari dari mana saja selama perangkat pengguna terhubung ke jaringan
internet. 2. \textbf{Ketiadaan Degradasi Fisik}: Materi bacaan digital
tidak mengalami kerusakan fisik kertas, sobek, ataupun memudar akibat
usia dan kelembapan udara. 3. \textbf{Pencarian Cepat dan Presisi
(\emph{Instant Searchability})}: Pengguna dapat mencari buku berdasarkan
kata kunci judul, nama pengarang, maupun kategori tematik dalam hitungan
milidetik. 4. \textbf{Efisiensi Penggandaan dan Distribusi}: Satu modul
bacaan digital dapat dibaca secara simultan oleh puluhan warga tanpa
memerlukan pencetakan fisik berbiaya tinggi.

\subsection{2.2.2 Otomasi Sirkulasi Perpustakaan Balai
Desa}\label{otomasi-sirkulasi-perpustakaan-balai-desa}

Meskipun penyediaan materi digital merupakan terobosan utama,
perpustakaan fisik balai desa tetap memiliki peranan penting bagi warga
yang menyukai buku cetak. Oleh karena itu, sistem perpustakaan modern
menerapkan konsep hibrida (\emph{hybrid library}), yaitu mendigitalkan
tata kelola sirkulasi fisik di samping menyediakan bacaan digital.
Otomasi sirkulasi mencakup pencatatan data peminjam, penentuan masa
pinjam, perhitungan jatuh tempo keterlambatan (\emph{overdue}),
perpanjangan izin pinjam, dan pencatatan riwayat pengembalian barang
secara terpusat dalam basis data.

\section{2.3 Literasi Digital dan Pemberdayaan Masyarakat
Pedesaan}\label{literasi-digital-dan-pemberdayaan-masyarakat-pedesaan}

Literasi digital bukan sekadar kemampuan teknis dalam mengoperasikan
perangkat gawai cerdas, melainkan kemampuan kognitif dan kritis untuk
menemukan, mengevaluasi, memanfaatkan, dan mengomunikasikan informasi
secara bertanggung jawab (Bawden, 2008).

Di wilayah perdesaan Indonesia, khususnya di Jawa Barat, program
literasi desa menghadapi tantangan heterogenitas tingkat pendidikan
masyarakat. Berdasarkan konsep \emph{Smart Village} yang dikembangkan
oleh Kementerian Desa, Pembangunan Daerah Tertinggal, dan Transmigrasi,
salah satu pilar utamanya adalah \emph{Smart People} (Masyarakat Cerdas)
yang memiliki akses setara terhadap ilmu pengetahuan terapan. Di samping
itu, pelestarian kearifan lokal Sunda (\emph{local indigenous
knowledge}) menuntut penyediaan materi budaya dan antarmuka dwi-bahasa
(Bahasa Indonesia dan Basa Sunda) guna mempertahankan identitas kultural
generasi muda pedesaan.

\section{2.4 Metode Pengembangan Perangkat Lunak: Agile
Prototyping}\label{metode-pengembangan-perangkat-lunak-agile-prototyping}

Pengembangan perangkat lunak untuk kebutuhan komunitas pedesaan
memerlukan pendekatan yang lincah, adaptif, dan responsif terhadap
perubahan masukan pengguna di lapangan. Pendekatan yang dipilih adalah
model \textbf{Agile Iterative Prototyping} (Pressman \& Maxim, 2015).

\begin{verbatim}
+-------------------------------------------------------------------------+
|                  SIKLUS AGILE ITERATIVE PROTOTYPING                     |
|                                                                         |
|   [1. Communication] ----> [2. Quick Plan] ----> [3. Quick Modeling]    |
|           ^                                              |              |
|           |                                              v              |
|   [6. Delivery/Feedback] <--- [5. Testing] <--- [4. Construction]       |
+-------------------------------------------------------------------------+
\end{verbatim}

Tahapan siklus meliputi: 1. \textbf{Komunikasi (\emph{Communication})}:
Mendiskusikan kendala nyata dengan pemangku kepentingan desa (Kades,
pamong, warga). 2. \textbf{Perencanaan Cepat (\emph{Quick Plan})}:
Menyusun \emph{backlog} kebutuhan fungsional dan prioritas implementasi.
3. \textbf{Pemodelan Rancangan (\emph{Modeling Quick Design})}: Membuat
sketsa antarmuka, diagram UML, dan skema basis data awal. 4.
\textbf{Konstruksi Prototipe (\emph{Construction of Prototype})}:
Membangun modul kode yang berfungsi (\emph{working software}). 5.
\textbf{Pengujian (\emph{Testing})}: Melakukan pengujian fungsionalitas
dan integritas data. 6. \textbf{Penyerahan dan Umpan Balik
(\emph{Deployment \& Feedback})}: Memperlihatkan sistem kepada pengguna
untuk mendapatkan koreksi langsung sebelum iterasi berikutnya.

\section{2.5 Unified Modeling Language (UML) dan Pemodelan Data
Relasional}\label{unified-modeling-language-uml-dan-pemodelan-data-relasional}

Menurut Martin Fowler (2004), \emph{Unified Modeling Language} (UML)
adalah keluarga notasi grafis standar yang didukung oleh meta-model
tunggal, yang membantu perancangan dan pendokumentasian sistem perangkat
lunak berorientasi objek.

Dalam penelitian ini, diagram UML yang digunakan meliputi: - \textbf{Use
Case Diagram}: Memodelkan interaksi antara aktor luar (pengguna) dengan
fungsi-fungsi layanan yang disediakan oleh sistem informasi. -
\textbf{Activity Diagram}: Memodelkan alur kerja (\emph{workflow})
langkah demi langkah dari suatu aktivitas bisnis sistem, baik yang
dilakukan oleh aktor maupun oleh sistem internal. -
\textbf{Entity-Relationship Diagram (ERD)}: Menggambarkan struktur
konseptual basis data yang terdiri dari kumpulan entitas relasional,
atribut kunci (\emph{primary \& foreign key}), serta derajat
kardinalitas antar-entitas (\emph{one-to-one, one-to-many,
many-to-many}).

\section{2.6 Landasan Arsitektur Jamstack dan Teknologi Web
Modern}\label{landasan-arsitektur-jamstack-dan-teknologi-web-modern}

\subsection{2.6.1 Paradigma Arsitektur Jamstack dan
Serverless}\label{paradigma-arsitektur-jamstack-dan-serverless}

Arsitektur Jamstack (akronim dari \emph{JavaScript, APIs, and Markup})
merupakan pendekatan perancangan web modern di mana lapisan presentasi
dipisahkan secara tegas dari lapisan data dan layanan backend.
Halaman-halaman web dapat di-\emph{render} di sisi server secara
\emph{hybrid} (\emph{Server-Side Rendering / SSR} dan \emph{Static Site
Generation / SSG}) lalu didistribusikan melalui jaringan komputasi
\emph{Edge Network} (Vercel). Arsitektur ini menghadirkan performa
tinggi, waktu muat laman yang sangat singkat pada peramban seluler,
serta meminimalkan permukaan serangan keamanan (\emph{attack surface})
karena ketiadaan server web tradisional yang terekspos langsung.

\subsection{2.6.2 Kerangka Kerja Next.js, React, dan
TypeScript}\label{kerangka-kerja-next.js-react-dan-typescript}

Next.js (versi 16) merupakan kerangka kerja rekayasa aplikasi web modern
berbasis pustaka \textbf{React 19} yang dikembangkan oleh Vercel.
Next.js memperkenalkan paradigma \emph{App Router} yang memanfaatkan
\emph{React Server Components} (RSC). Dengan RSC, komponen antarmuka
yang memerlukan pengambilan data (\emph{data fetching}) dapat diproses
secara langsung di lingkungan server tanpa mengirimkan beban
\emph{bundle} JavaScript yang berat ke peramban klien. Penggunaan bahasa
pemrograman \textbf{TypeScript} menambahkan pengetikan tipe data statis
(\emph{static typing}) yang ketat, meminimalkan potensi kesalahan
\emph{runtime bug}, serta meningkatkan pemeliharaan kode jangka panjang.

\subsection{2.6.3 Basis Data Relasional Serverless: Neon PostgreSQL dan
Prisma
ORM}\label{basis-data-relasional-serverless-neon-postgresql-dan-prisma-orm}

Untuk menjamin persistensi data yang konsisten dan berintegritas tinggi,
sistem memanfaatkan \textbf{PostgreSQL}, sistem manajemen basis data
relasional (\emph{RDBMS}) terkemuka yang mendukung standar SQL secara
penuh, integritas referensial bersyarat (\emph{foreign keys}), dan tipe
data tingkat lanjut seperti JSONB. Layanan \textbf{Neon Serverless
PostgreSQL} mengisolasi lapisan komputasi (\emph{compute}) dari lapisan
penyimpanan (\emph{storage}), memungkinkan basis data melakukan
penskalaan otomatis (\emph{auto-scaling}) dari nol
(\emph{scale-to-zero}) dan mengoptimalkan latensi jaringan melalui
\emph{connection pooling} berbasis WebSocket.

Interaksi antara kode aplikasi Next.js dan basis data PostgreSQL
dijembatani oleh \textbf{Prisma ORM (versi 7)}. Prisma mengonversi skema
basis data yang didefinisikan dalam berkas \texttt{schema.prisma}
menjadi klien kueri TypeScript yang sepenuhnya berkeamanan tipe
(\emph{type-safe}). Hal ini mencegah kerentanan injeksi SQL (\emph{SQL
Injection vulnerabilities}) dan menyederhanakan manipulasi transaksi
relasi data yang kompleks.

\subsection{2.6.4 Keamanan Otentikasi: NIK, Kriptografi Bcrypt, dan
NextAuth}\label{keamanan-otentikasi-nik-kriptografi-bcrypt-dan-nextauth}

Keamanan akses sistem bertumpu pada protokol standar industri: 1.
\textbf{Nomor Induk Kependudukan (NIK)}: NIK 16 digit digunakan sebagai
pengenal unik tunggal (\emph{unique identifier}) warga, menggantikan
alamat email yang umumnya jarang dimiliki oleh warga pedesaan usia
dewasa atau lanjut usia. 2. \textbf{Kriptografi Hashing Bcrypt}: Kode
PIN 6 digit yang dimasukkan oleh pengguna tidak pernah disimpan dalam
bentuk teks polos (\emph{plaintext}). Sistem menggunakan fungsi
\emph{hashing} adaptif \textbf{Bcrypt} dengan faktor biaya komputasi
\emph{salt rounds} bernilai 10. Bcrypt secara otomatis menggabungkan
nilai acak (\emph{salt}) 128-bit ke dalam teks masukan sebelum melalui
iterasi enkripsi Blowfish, sehingga kebal terhadap serangan tabel
pelangi (\emph{rainbow table attack}). 3. \textbf{NextAuth.js
(Auth.js)}: Menyediakan mekanisme otentikasi berbasis sesi terenkripsi
menggunakan token web JSON (\emph{JSON Web Token / JWT}). Token sesi
disimpan dalam \emph{HTTP-only, Secure Cookie} di peramban klien,
melindungi sesi pengguna dari serangan pencurian skrip lintas situs
(\emph{Cross-Site Scripting / XSS}).

\subsection{2.6.5 Desain Antarmuka: Tailwind CSS dan Pendekatan
Mobile-First}\label{desain-antarmuka-tailwind-css-dan-pendekatan-mobile-first}

Antarmuka pengguna dirancang menggunakan \textbf{Tailwind CSS}, kerangka
kerja CSS berbasis \emph{utility-first} yang memungkinkan penyusunan
tampilan modern, fleksibel, dan konsisten langsung di dalam berkas
komponen. Mengingat lebih dari 85\% warga desa mengakses internet
melalui perangkat ponsel cerdas, antarmuka dirancang dengan paradigma
\textbf{Mobile-First Responsive Design}. Tata letak antarmuka dirancang
agar nyaman dioperasikan dengan satu tangan (\emph{thumb-friendly}),
memiliki kontras warna yang memenuhi standar aksesibilitas WCAG 2.1 AA,
dan menyediakan mode tema gelap (\emph{dark mode}) otomatis untuk
kenyamanan membaca di malam hari.

\section{2.7 Pengujian Perangkat Lunak dan Standar ISO/IEC
25010}\label{pengujian-perangkat-lunak-dan-standar-isoiec-25010}

Pengujian perangkat lunak (\emph{software testing}) merupakan proses
eksekusi program dengan tujuan menemukan ketidaksesuaian antara perilaku
aktual sistem dan spesifikasi kebutuhan yang telah ditetapkan.

Pendekatan pengujian yang digunakan dalam penelitian ini mengacu pada:
1. \textbf{Pengujian Kotak Hitam (\emph{Black-Box Testing})}: Metode
pengujian fungsional di mana penguji mengevaluasi masukan dan keluaran
sistem tanpa perlu mengetahui struktur logika internal kode program.
Teknik yang diterapkan meliputi \emph{Equivalence Partitioning} (membagi
masukan ke dalam kelas data valid dan tidak valid) dan \emph{Boundary
Value Analysis} (menguji nilai batas masukan seperti panjang NIK 16
digit dan PIN 6 digit). 2. \textbf{Pengujian Regresi (\emph{Regression
Testing})}: Pengujian otomatis yang dijalankan secara berulang setiap
kali terjadi modifikasi kode untuk memastikan bahwa penambahan fitur
baru tidak merusak (\emph{break}) fungsionalitas yang telah berjalan
sebelumnya. 3. \textbf{Standar Kualitas Perangkat Lunak ISO/IEC 25010}:
Standar internasional yang mendefinisikan model kualitas produk
perangkat lunak, mencakup 8 karakteristik utama: Kesesuaian Fungsional
(\emph{Functional Suitability}), Efisiensi Kinerja (\emph{Performance
Efficiency}), Kompatibilitas (\emph{Compatibility}), Kemudahan
Penggunaan (\emph{Usability}), Keandalan (\emph{Reliability}), Keamanan
(\emph{Security}), Kemudahan Pemeliharaan (\emph{Maintainability}), dan
Portabilitas (\emph{Portability}).

\section{2.8 Kajian Penelitian
Terdahulu}\label{kajian-penelitian-terdahulu}

Sebagai landasan perbandingan dan penegasan posisi kebaruan
(\emph{novelty}) sistem informasi Pustaka Pangkalan, Tabel 2.1
menyajikan komparasi terhadap 4 (empat) penelitian dan pengembangan
sistem informasi perpustakaan desa terdahulu.

\emph{Tabel 2.1. Komparasi Penelitian dan Sistem Informasi Perpustakaan
Terdahulu}

{\def\LTcaptype{none} % do not increment counter
\begin{longtable}[]{@{}
  >{\raggedright\arraybackslash}p{(\linewidth - 10\tabcolsep) * \real{0.1667}}
  >{\raggedright\arraybackslash}p{(\linewidth - 10\tabcolsep) * \real{0.1667}}
  >{\raggedright\arraybackslash}p{(\linewidth - 10\tabcolsep) * \real{0.1667}}
  >{\raggedright\arraybackslash}p{(\linewidth - 10\tabcolsep) * \real{0.1667}}
  >{\raggedright\arraybackslash}p{(\linewidth - 10\tabcolsep) * \real{0.1667}}
  >{\raggedright\arraybackslash}p{(\linewidth - 10\tabcolsep) * \real{0.1667}}@{}}
\toprule\noalign{}
\begin{minipage}[b]{\linewidth}\raggedright
Peneliti / Tahun
\end{minipage} & \begin{minipage}[b]{\linewidth}\raggedright
Judul Penelitian / Sistem
\end{minipage} & \begin{minipage}[b]{\linewidth}\raggedright
Stack Teknologi
\end{minipage} & \begin{minipage}[b]{\linewidth}\raggedright
Ruang Lingkup Layanan
\end{minipage} & \begin{minipage}[b]{\linewidth}\raggedright
Kelemahan / Keterbatasan
\end{minipage} & \begin{minipage}[b]{\linewidth}\raggedright
Keunggulan Pustaka Pangkalan (As-Built)
\end{minipage} \\
\midrule\noalign{}
\endhead
\bottomrule\noalign{}
\endlastfoot
\textbf{Pratama \& Setiawan (2021)} & Sistem Informasi Perpustakaan Desa
Berbasis Web (Studi Kasus Desa Sukamaju) & PHP Native, MySQL, Bootstrap
& Katalog buku fisik dan pencatatan sirkulasi manual oleh admin & Tidak
ada portal pembaca e-book digital; otentikasi mengharuskan email; tidak
ada fitur gamifikasi & Menyediakan pembaca e-book digital per bab;
otentikasi NIK+PIN; gamifikasi poin dan lencana; arsitektur Next.js
serverless \\
\textbf{Hidayat dkk. (2022)} & Rancang Bangun Digital Library Desa
Cerdas Berbasis Android & Java / Android Native, Firebase Realtime DB &
Pembacaan file PDF digital dan kartu anggota & Memerlukan instalasi
aplikasi APK (memberatkan memori ponsel warga); biaya penyimpanan
Firebase tinggi & Berbasis web responsif (\emph{zero-install}); ringan
di browser; basis data relasional PostgreSQL berintegritas ketat \\
\textbf{Rahmawati \& Nugroho (2023)} & Otomasi Sirkulasi Perpustakaan
Balai Desa Menggunakan SLIMS & PHP, MariaDB, SLIMS 9 Bulian & Manajemen
sirkulasi fisik balai desa dan pencetakan barcode & Terlalu rumit untuk
warga desa; tidak dirancang untuk membaca buku daring; antarmuka belum
mobile-friendly & Antarmuka disesuaikan khusus untuk warga desa;
integrasi dwibahasa Sunda; asisten virtual Kades AI \\
\textbf{Wahyudi (2024)} & Pengembangan E-Perpus Desa Berbasis Laravel &
Laravel 10, MySQL, Blade & Modul katalog dan peminjaman buku & Belum
mendukung lokalisasi bahasa daerah; tidak ada mekanisme pencegahan lupa
sandi tanpa email & Otentikasi NIK 16 digit terproteksi Bcrypt;
fasilitas Reset PIN warga oleh admin balai desa; 100\% lulus STQA \\
\end{longtable}
}

\newpage

\chapter{BAB III --- METODOLOGI PENELITIAN DAN
PENGEMBANGAN}\label{bab-iii-metodologi-penelitian-dan-pengembangan}

\section{3.1 Tempat dan Waktu
Pelaksanaan}\label{tempat-dan-waktu-pelaksanaan}

Kegiatan penelitian, perancangan, dan implementasi sistem informasi ini
dilaksanakan di lingkungan administratif Desa Pangkalan, Kecamatan
Cikidang, Kabupaten Sukabumi, Provinsi Jawa Barat. Rangkaian kegiatan
berlangsung selama periode program Kuliah Kerja Nyata (KKN) Tematik
Semester Ganjil Tahun Akademik 2026, mulai dari bulan Juli 2026 hingga
September 2026.

Distribusi lokasi kegiatan di lapangan mencakup: 1. \textbf{Kantor Balai
Desa Pangkalan}: Sebagai pusat koordinasi pemerintahan desa, observasi
kondisi fisik perpustakaan desa, serta wawancara bersama Kepala Desa dan
staf administrasi. 2. \textbf{Dusun Pangkalan, Dusun Cikajang, Dusun
Pasir Arangan, dan Dusun Pasir Gombong}: Sebagai lokasi observasi
demografis, sosialisasi literasi, penyebaran kuesioner kebutuhan warga,
dan uji coba antarmuka pada ponsel masyarakat.

\section{\texorpdfstring{3.2 Kerangka Pikir Penelitian (\emph{Research
Framework})}{3.2 Kerangka Pikir Penelitian (Research Framework)}}\label{kerangka-pikir-penelitian-research-framework}

Kerangka pikir dalam penelitian rekayasa perangkat lunak ini digambarkan
secara sistematis pada Gambar 3.1 untuk memetakan alur pemecahan masalah
dari identifikasi kendala hingga penyerahan sistem.

\begin{verbatim}
+-------------------------------------------------------------------------+
|                       KERANGKA PIKIR PENELITIAN                         |
|                                                                         |
|  [MASALAH LAPANGAN]                                                     |
|  - Koleksi buku fisik balai desa terbatas & pencatatan sirkulasi manual |
|  - Warga di 4 dusun terpisah kesulitan menjangkau balai desa            |
|  - Belum tersedianya modul praktis terapan dan media literasi digital   |
|                                     |                                   |
|                                     v                                   |
|  [PENGUMPULAN DATA & ANALISIS]                                          |
|  - Observasi fisik balai desa & wawancara aparatur / tokoh pemuda       |
|  - Studi literatur sistem informasi, perpustakaan digital, Jamstack     |
|  - Analisis kebutuhan fungsional (FR-01 - FR-20) dan non-fungsional    |
|                                     |                                   |
|                                     v                                   |
|  [PERANCANGAN SISTEM (DESIGN)]                                          |
|  - Pemodelan UML (Use Case & Activity Diagrams)                         |
|  - Perancangan Arsitektur 4-Tier Jamstack & Skema ERD Relasional        |
|  - Perancangan Antarmuka Responsif Mobile-First (Bilingual ID/SU)       |
|                                     |                                   |
|                                     v                                   |
|  [KONSTRUKSI PERANGKAT LUNAK (IMPLEMENTATION)]                          |
|  - Frontend & Backend: Next.js 16 (React 19, TypeScript, Tailwind)      |
|  - Basis Data & ORM: Neon Serverless PostgreSQL & Prisma ORM 7          |
|  - Keamanan: Validasi NIK, Hashing Bcrypt, Sesi NextAuth JWT            |
|  - Fitur Khusus: Chapter Reader, Gamifikasi, Kades AI, Branding Sukabumi|
|                                     |                                   |
|                                     v                                   |
|  [PENGUJIAN & EVALUASI MUTU (STQA)]                                     |
|  - Pengujian Black-Box & Regression Suite (9 Test Suites, 48 Skenario)  |
|  - Evaluasi Kualitas Perangkat Lunak ISO/IEC 25010 (100% Pass Rate)     |
|                                     |                                   |
|                                     v                                   |
|  [HASIL AKHIR & SERAH TERIMA]                                           |
|  - Sistem Terdeploy: https://perpus-pangkalan.vercel.app                |
|  - Berita Acara Serah Terima (BAST) & Buku Panduan Pengguna Lengkap     |
+-------------------------------------------------------------------------+
\end{verbatim}

\emph{Gambar 3.1. Alur Kerangka Pikir Penelitian dan Pengembangan
Sistem}

\section{3.3 Metode Pengumpulan Data}\label{metode-pengumpulan-data}

Untuk menjamin keakuratan spesifikasi kebutuhan sistem, pengumpulan data
dilakukan melalui 3 (tiga) teknik utama:

\subsection{\texorpdfstring{3.3.1 Observasi Langsung (\emph{Direct
Observation})}{3.3.1 Observasi Langsung (Direct Observation)}}\label{observasi-langsung-direct-observation}

Tim peneliti melakukan peninjauan langsung ke ruangan perpustakaan Balai
Desa Pangkalan guna mengamati: - Tata letak dan kondisi fisik lemari rak
buku; - Buku besar pencatatan peminjaman dan kartu kontrol manual yang
sedang berjalan; - Ketersediaan sarana komputer dan konektivitas
internet di balai desa; - Kondisi penerimaan sinyal seluler di
titik-titik permukiman warga di Dusun Cikajang, Pasir Arangan, dan Pasir
Gombong.

\subsection{\texorpdfstring{3.3.2 Wawancara Mendalam (\emph{In-depth
Interview})}{3.3.2 Wawancara Mendalam (In-depth Interview)}}\label{wawancara-mendalam-in-depth-interview}

Wawancara semi-terstruktur dilakukan dengan para pemangku kepentingan
kunci (\emph{key stakeholders}): 1. \textbf{Kepala Desa Pangkalan}:
Menggali visi pembangunan desa, arah kebijakan digitalisasi desa, serta
harapan terhadap peningkatan wawasan tani dan usaha warga. 2.
\textbf{Sekretaris Desa dan Pengelola Perpustakaan Balai Desa}:
Mengidentifikasi kendala harian dalam pencatatan sirkulasi buku,
frekuensi peminjaman, masalah keterlambatan pengembalian buku, dan
format laporan yang dibutuhkan. 3. \textbf{Ketua Karang Taruna dan Tokoh
Pemuda}: Menjaring aspirasi pemuda terkait minat topik bacaan
(teknologi, bisnis UMKM, kebudayaan Sunda) dan preferensi desain
antarmuka aplikasi ponsel.

\subsection{\texorpdfstring{3.3.3 Studi Dokumentasi dan Kepustakaan
(\emph{Documentary
Study})}{3.3.3 Studi Dokumentasi dan Kepustakaan (Documentary Study)}}\label{studi-dokumentasi-dan-kepustakaan-documentary-study}

Pengumpulan data sekunder dilakukan melalui penelaahan: - Dokumen
monografi dan profil wilayah Desa Pangkalan, Kecamatan Cikidang; -
Standar penamaan dusun resmi di bawah naungan Pemerintah Desa Pangkalan;
- Literatur ilmiah mengenai rekayasa perangkat lunak, dokumentasi resmi
Next.js, Prisma ORM, dan PostgreSQL; - Peraturan perundang-undangan
terkait perlindungan data pribadi dan standar pelayanan informasi publik
desa.

\section{\texorpdfstring{3.4 Tahapan Siklus Pengembangan Sistem
(\emph{Agile Prototyping
Stages})}{3.4 Tahapan Siklus Pengembangan Sistem (Agile Prototyping Stages)}}\label{tahapan-siklus-pengembangan-sistem-agile-prototyping-stages}

Pengembangan sistem informasi dilaksanakan melalui 6 (enam) tahapan
berulang (\emph{iterative}):

\subsection{1. Tahap Inisiasi dan Analisis
Kebutuhan}\label{tahap-inisiasi-dan-analisis-kebutuhan}

Mengidentifikasi seluruh aktor sistem, mendefinisikan kendala sistem
berjalan, dan menyusun matriks spesifikasi kebutuhan fungsional (FR-01
hingga FR-20) serta kebutuhan non-fungsional (keamanan, performa,
ketersediaan, usabilitas).

\subsection{\texorpdfstring{2. Tahap Perancangan Sistem (\emph{System
Design})}{2. Tahap Perancangan Sistem (System Design)}}\label{tahap-perancangan-sistem-system-design}

\begin{itemize}
\tightlist
\item
  Merancang alur logika bisnis menggunakan diagram alur
  (\emph{flowchart});
\item
  Memodelkan skenario fungsional ke dalam \emph{Use Case Diagram}
  lengkap dengan narasi skenario use case (\emph{Use Case Narrative});
\item
  Memodelkan alur proses transaksi ke dalam \emph{Activity Diagram};
\item
  Merancang struktur data relasional menggunakan \emph{Entity
  Relationship Diagram} (ERD) dan menyusun kamus data (\emph{Data
  Dictionary}) terperinci;
\item
  Merancang tata letak antarmuka pengguna (\emph{wireframing})
  berorientasi gawai seluler (\emph{mobile-first}).
\end{itemize}

\subsection{\texorpdfstring{3. Tahap Konstruksi Kode Sumber
(\emph{Software
Construction})}{3. Tahap Konstruksi Kode Sumber (Software Construction)}}\label{tahap-konstruksi-kode-sumber-software-construction}

\begin{itemize}
\tightlist
\item
  Menyiapkan repositori kode sumber di GitHub dengan standar kontrol
  versi Git;
\item
  Mengonfigurasi skema deklaratif basis data pada berkas
  \texttt{prisma/schema.prisma} dan mengeksekusi migrasi ke Neon
  Serverless PostgreSQL;
\item
  Membangun komponen UI frontend interaktif menggunakan React 19,
  TypeScript, dan Tailwind CSS;
\item
  Membangun antarmuka pemrograman backend (\emph{Route Handlers}) pada
  direktori \texttt{src/app/api/};
\item
  Mengintegrasikan mekanisme keamanan otentikasi NIK dan PIN Bcrypt via
  NextAuth;
\item
  Mengimplementasikan fitur-fitur pembeda: \emph{chapter reader},
  gamifikasi literasi, asisten \emph{Kades AI}, dan pengalihan bahasa
  Indonesia / Sunda.
\end{itemize}

\subsection{\texorpdfstring{4. Tahap Pengujian Kualitas Perangkat Lunak
(\emph{STQA})}{4. Tahap Pengujian Kualitas Perangkat Lunak (STQA)}}\label{tahap-pengujian-kualitas-perangkat-lunak-stqa}

\begin{itemize}
\tightlist
\item
  Menjalankan pengujian \emph{black-box} fungsional pada setiap
  antarmuka pengguna;
\item
  Menjalankan 9 test suite otomatis mencakup 48 skenario uji regresi;
\item
  Menguji ketahanan rute administratif dari akses ilegal
  (\emph{unauthorized access});
\item
  Memvalidasi konsistensi cascade update pada basis data relasional.
\end{itemize}

\subsection{5. Tahap Deployment Produksi dan
Optimasi}\label{tahap-deployment-produksi-dan-optimasi}

\begin{itemize}
\tightlist
\item
  Menghubungkan repositori GitHub ke platform \emph{Vercel Edge Network}
  untuk \emph{Continuous Integration / Continuous Deployment} (CI/CD);
\item
  Mengonfigurasi variabel lingkungan (\emph{environment variables})
  secara aman tanpa membocorkan kredensial rahasia (\emph{zero secrets
  policy});
\item
  Menetapkan domain resmi \textbf{https://perpus-pangkalan.vercel.app};
\item
  Mengonfigurasi metadata OpenGraph dan favicon resmi Lambang Kabupaten
  Sukabumi.
\end{itemize}

\subsection{\texorpdfstring{6. Tahap Serah Terima dan Alih Kelola
(\emph{Handover})}{6. Tahap Serah Terima dan Alih Kelola (Handover)}}\label{tahap-serah-terima-dan-alih-kelola-handover}

\begin{itemize}
\tightlist
\item
  Melakukan demonstrasi dan simulasi operasional bersama aparatur Balai
  Desa Pangkalan;
\item
  Menyerahkan berkas Berita Acara Serah Terima (BAST) resmi;
\item
  Menyerahkan buku manual panduan operasional bagi warga dan
  administrator desa.
\end{itemize}

\section{3.5 Lingkungan Perangkat Keras dan Perangkat
Lunak}\label{lingkungan-perangkat-keras-dan-perangkat-lunak}

Pelaksanaan perancangan dan pengujian sistem didukung oleh instrumen
perangkat keras dan perangkat lunak dengan spesifikasi sebagai berikut:

\subsection{3.5.1 Spesifikasi Perangkat Keras
Pengembangan}\label{spesifikasi-perangkat-keras-pengembangan}

\begin{enumerate}
\def\labelenumi{\arabic{enumi}.}
\tightlist
\item
  \textbf{Komputer Kerja Utama}: Prosesor Intel Core i5 / AMD Ryzen 5,
  RAM 16 GB DDR4, Penyimpanan Solid State Drive (SSD) 512 GB NVMe.
\item
  \textbf{Perangkat Pengujian Seluler (Mobile Device)}: Ponsel pintar
  Android (layar 6.5 inci, RAM 4 GB, resolusi FHD+) dan iOS (layar 6.1
  inci, resolusi Super Retina) guna memverifikasi responsivitas
  tampilan.
\end{enumerate}

\subsection{3.5.2 Spesifikasi Perangkat Lunak
Pengembangan}\label{spesifikasi-perangkat-lunak-pengembangan}

\begin{enumerate}
\def\labelenumi{\arabic{enumi}.}
\tightlist
\item
  \textbf{Sistem Operasi}: Microsoft Windows 11 Pro 64-bit;
\item
  \textbf{Penyunting Kode Sumber (\emph{Code Editor})}: Visual Studio
  Code / Antigravity IDE;
\item
  \textbf{Runtime Environment}: Node.js versi 20 LTS;
\item
  \textbf{Manajer Paket}: Node Package Manager (npm) versi 10;
\item
  \textbf{Bahasa Pemrograman}: TypeScript versi 5, JavaScript (ES2024),
  SQL;
\item
  \textbf{Kerangka Kerja Web}: Next.js versi 16.3.0 (React 19);
\item
  \textbf{Pustaka CSS}: Tailwind CSS versi 4;
\item
  \textbf{Object-Relational Mapping}: Prisma ORM versi 7.9.1;
\item
  \textbf{Basis Data Produksi}: Neon Serverless PostgreSQL 16;
\item
  \textbf{Platform Deployment Cloud}: Vercel Serverless Edge Platform;
\item
  \textbf{Sistem Kontrol Versi}: Git versi 2.45 dan GitHub;
\item
  \textbf{Peramban Web Pengujian}: Google Chrome, Mozilla Firefox,
  Microsoft Edge, dan Safari Mobile.
\end{enumerate}

\newpage

\chapter{BAB IV --- ANALISIS DAN PERANCANGAN
SISTEM}\label{bab-iv-analisis-dan-perancangan-sistem}

\section{\texorpdfstring{4.1 Analisis Sistem Berjalan (\emph{As-Is
System
Analysis})}{4.1 Analisis Sistem Berjalan (As-Is System Analysis)}}\label{analisis-sistem-berjalan-as-is-system-analysis}

Berdasarkan hasil observasi dan wawancara di Balai Desa Pangkalan,
prosedur tata kelola perpustakaan yang sedang berjalan sebelum adanya
sistem informasi dapat digambarkan melalui alur kerja manual berikut: 1.
Warga yang membutuhkan buku bacaan harus datang secara fisik ke kantor
balai desa pada hari dan jam kerja dinas (Senin sampai Jumat pukul 08.00
- 15.00 WIB). 2. Warga mencari buku secara manual dengan menelusuri
lemari rak perpustakaan balai desa satu per satu tanpa bantuan katalog
pencarian. 3. Apabila buku yang dicari tidak ditemukan (karena sedang
dipinjam warga lain atau tidak tersedia), waktu dan biaya perjalanan
warga menjadi sia-sia. 4. Apabila buku ditemukan dan ingin dipinjam,
warga menyerahkan buku kepada petugas balai desa untuk dicatat pada buku
besar peminjaman manual. 5. Petugas mencatat: Tanggal Pinjam, Nama
Peminjam, Alamat Dusun, Judul Buku, dan Nomor Telepon (jika ada). 6.
Buku dibawa pulang oleh peminjam dengan jangka waktu pinjam lisan
(biasanya 7 hari). 7. Permasalahan timbul saat buku terlambat
dikembalikan: petugas balai desa tidak memiliki sistem notifikasi atau
rekapitulasi keterlambatan otomatis. Buku catatan peminjaman sering
terselip, rusak, atau terkena tumpahan air, sehingga banyak eksemplar
buku balai desa yang hilang tanpa jejak peminjam yang jelas.

Kelemahan utama sistem berjalan ini mencakup: - \textbf{Inefisiensi
Waktu dan Jarak}: Warga di Dusun Cikajang, Pasir Arangan, dan Pasir
Gombong berjarak 3 hingga 6 kilometer dari balai desa; -
\textbf{Ketidaktersediaan Akses Bacaan 24 Jam}: Buku hanya bisa dibaca
saat balai desa buka; - \textbf{Risiko Kehilangan Aset Desa}: Tingginya
persentase buku hilang akibat pencatatan kertas yang rentan tercecer; -
\textbf{Ketiadaan Data Analitik Minat Baca}: Pemerintah desa tidak
memiliki data statistik mengenai topik buku yang diminati oleh warganya.

\section{\texorpdfstring{4.2 Analisis Sistem Baru yang Diusulkan
(\emph{To-Be System
Analysis})}{4.2 Analisis Sistem Baru yang Diusulkan (To-Be System Analysis)}}\label{analisis-sistem-baru-yang-diusulkan-to-be-system-analysis}

Sistem informasi \textbf{Pustaka Pangkalan} yang diusulkan dan dibangun
mentransformasi tata kelola literasi desa menjadi sistem hibrida
terintegrasi: 1. \textbf{Portal Literasi Digital Mandiri}: Warga dapat
mengakses portal kapan saja (24/7) melalui ponsel cerdas. Warga dapat
menjelajahi katalog buku dalam 8 kategori tematik, membaca isi teks bab
buku (\emph{chapter reader}) secara instan tanpa mengunduh berkas
berukuran besar, menandai bab favorit (\emph{bookmark}), memberi ulasan
dan rating, serta berkonsultasi dengan asisten cerdas \emph{Kades AI}.
2. \textbf{Digitalisasi Sirkulasi Balai Desa}: Peminjaman buku fisik di
balai desa dicatat secara digital ke dalam sistem oleh pengelola balai
desa (\texttt{/admin/circulation}). Sistem secara otomatis mencatat
tanggal pinjam, menghitung batas kembali (+7 hari), menampilkan status
keterlambatan (\emph{overdue}), menyediakan tombol perpanjangan (+7
hari) sekali klik, dan memperbarui status ketersediaan eksemplar buku
secara instan. 3. \textbf{Gamifikasi Partisipasi Warga}: Setiap bab yang
dibaca warga secara otomatis memberikan akumulasi poin literasi,
mencatat rekor hari aktif berturut-turut (\emph{reading streak}), serta
menganugerahkan lencana prestasi (\emph{badges}) pada kartu anggota
digital ber-QR Code.

\section{4.3 Analisis Profil Aktor
Sistem}\label{analisis-profil-aktor-sistem}

Berdasarkan implementasi kode sumber dan hak akses sistem terpasang,
aktor sistem terbagi menjadi:

\emph{Tabel 4.1. Matriks Profil dan Wewenang Aktor Sistem}

{\def\LTcaptype{none} % do not increment counter
\begin{longtable}[]{@{}
  >{\raggedright\arraybackslash}p{(\linewidth - 4\tabcolsep) * \real{0.3333}}
  >{\raggedright\arraybackslash}p{(\linewidth - 4\tabcolsep) * \real{0.3333}}
  >{\raggedright\arraybackslash}p{(\linewidth - 4\tabcolsep) * \real{0.3333}}@{}}
\toprule\noalign{}
\begin{minipage}[b]{\linewidth}\raggedright
Aktor Sistem
\end{minipage} & \begin{minipage}[b]{\linewidth}\raggedright
Peran Teknis (\emph{Role})
\end{minipage} & \begin{minipage}[b]{\linewidth}\raggedright
Hak Akses dan Lingkup Wewenang
\end{minipage} \\
\midrule\noalign{}
\endhead
\bottomrule\noalign{}
\endlastfoot
\textbf{Warga Desa} & \texttt{USER} & - Melakukan registrasi mandiri
menggunakan NIK 16 digit dan PIN 6 digit.- Menjelajahi katalog buku dan
melakukan pencarian instan.- Membaca modul buku digital per bab
(\emph{chapter reader}).- Menandai bab bacaan (\emph{bookmark}) dan
menyimpan riwayat bacaan.- Mengelola rak buku favorit pribadi.-
Memberikan ulasan komentar dan rating bintang (1-5).- Memiliki Kartu
Anggota Digital ber-QR Code dengan poin dan badge.- Mengakses asisten
virtual literasi \emph{Kades AI}.- Mengganti bahasa antarmuka (Bahasa
Indonesia / Basa Sunda). \\
\textbf{Pengelola Balai Desa} & \texttt{ADMIN} & - Memiliki seluruh
wewenang Warga Desa (\texttt{USER}).- Mengakses Dashboard Metrik
Literasi Desa (\texttt{/admin}).- Melakukan manajemen katalog buku fisik
dan digital (CRUD buku).- Melakukan penyuntingan dan manajemen konten
bab bacaan e-book.- Melakukan pencatatan sirkulasi peminjaman buku fisik
balai desa.- Menyetujui perpanjangan masa pinjam (+7 hari) dan menandai
buku kembali.- Mengelola master data 8 kategori tematik dan 4 dusun
resmi.- Mengelola data warga dan melakukan reset PIN warga yang lupa.-
Melakukan moderasi dan penghapusan ulasan yang tidak pantas.-
Memublikasikan warta/maklumat literasi desa.- Mengekspor laporan
sirkulasi dan literasi ke format CSV. \\
\textbf{Pengunjung Tamu} & \texttt{GUEST} (Anonim) & - Mengakses beranda
publik dan membaca informasi profil perpustakaan.- Menjelajahi katalog
buku dan ringkasan sinopsis buku.- Membaca warta/maklumat literasi
desa.- Mengganti bahasa antarmuka (Bahasa Indonesia / Basa Sunda).-
Melakukan registrasi akun baru atau masuk ke sistem (\emph{login}). \\
\end{longtable}
}

\section{4.4 Analisis Kebutuhan Sistem}\label{analisis-kebutuhan-sistem}

\subsection{\texorpdfstring{4.4.1 Kebutuhan Fungsional (\emph{Functional
Requirements})}{4.4.1 Kebutuhan Fungsional (Functional Requirements)}}\label{kebutuhan-fungsional-functional-requirements}

Kebutuhan fungsional sistem dipetakan menggunakan klasifikasi prioritas
\textbf{MoSCoW} (\emph{Must have, Should have, Could have, Won't have})
sebagaimana tercantum pada Tabel 4.2.

\emph{Tabel 4.2. Matriks Kebutuhan Fungsional Sistem Terpasang (As-Built
FR)}

{\def\LTcaptype{none} % do not increment counter
\begin{longtable}[]{@{}
  >{\raggedright\arraybackslash}p{(\linewidth - 8\tabcolsep) * \real{0.1765}}
  >{\raggedright\arraybackslash}p{(\linewidth - 8\tabcolsep) * \real{0.1765}}
  >{\centering\arraybackslash}p{(\linewidth - 8\tabcolsep) * \real{0.2941}}
  >{\raggedright\arraybackslash}p{(\linewidth - 8\tabcolsep) * \real{0.1765}}
  >{\raggedright\arraybackslash}p{(\linewidth - 8\tabcolsep) * \real{0.1765}}@{}}
\toprule\noalign{}
\begin{minipage}[b]{\linewidth}\raggedright
Kode
\end{minipage} & \begin{minipage}[b]{\linewidth}\raggedright
Deskripsi Kebutuhan Fungsional
\end{minipage} & \begin{minipage}[b]{\linewidth}\centering
Prioritas
\end{minipage} & \begin{minipage}[b]{\linewidth}\raggedright
Aktor
\end{minipage} & \begin{minipage}[b]{\linewidth}\raggedright
Bukti Kode Sumber Implementasi
\end{minipage} \\
\midrule\noalign{}
\endhead
\bottomrule\noalign{}
\endlastfoot
\textbf{FR-01} & Sistem dapat mendaftarkan akun warga baru dengan
validasi NIK 16 digit, Nama Lengkap, PIN 6 digit, dan Dusun domisili &
Must & Warga & \texttt{src/app/login/page.tsx},
\texttt{src/app/onboarding/page.tsx}, \texttt{/api/auth/register} \\
\textbf{FR-02} & Sistem dapat memverifikasi otentikasi login menggunakan
kredensial NIK dan PIN terenkripsi Bcrypt & Must & Warga, Admin &
\texttt{src/lib/auth.ts},
\texttt{src/app/api/auth/{[}...nextauth{]}/route.ts} \\
\textbf{FR-03} & Sistem dapat menyajikan katalog buku publik dengan
filter 8 kategori tematik dan pencarian teks judul & Must & Semua &
\texttt{src/app/explore/page.tsx},
\texttt{src/app/api/books/route.ts} \\
\textbf{FR-04} & Sistem dapat menampilkan modul pembaca e-book per bab
(\emph{chapter reader}) dengan navigasi bab sebelumnya/selanjutnya &
Must & Warga & \texttt{src/app/read/{[}chapterId{]}/page.tsx},
\texttt{src/app/api/read/{[}chapterId{]}/route.ts} \\
\textbf{FR-05} & Sistem dapat menyimpan penanda bacaan (\emph{bookmark})
dan mencatat progres membaca pengguna secara otomatis & Must & Warga &
\texttt{src/app/api/bookmarks/route.ts},
\texttt{src/app/api/reading-progress/route.ts} \\
\textbf{FR-06} & Sistem dapat menyimpan buku ke dalam rak koleksi
pribadi pengguna & Should & Warga & \texttt{src/app/shelf/page.tsx},
\texttt{src/app/api/shelf/route.ts} \\
\textbf{FR-07} & Sistem dapat menerima ulasan komentar dan rating
bintang (1-5) pada halaman detail buku & Should & Warga &
\texttt{src/app/books/{[}id{]}/page.tsx},
\texttt{src/app/api/reviews/route.ts} \\
\textbf{FR-08} & Sistem dapat menerbitkan Kartu Anggota Digital ber-QR
Code dengan tampilan lencana dan total poin literasi & Must & Warga &
\texttt{src/app/profile/page.tsx},
\texttt{src/app/api/user/profile/route.ts} \\
\textbf{FR-09} & Sistem dapat mengalihkan bahasa antarmuka secara
dinamis antara Bahasa Indonesia dan Basa Sunda & Should & Semua &
\texttt{src/components/LanguageProvider.tsx},
\texttt{src/components/layout/TopAppBar.tsx} \\
\textbf{FR-10} & Sistem dapat melayani tanya jawab materi literasi desa
melalui asisten cerdas virtual \emph{Kades AI} & Could & Warga &
\texttt{src/components/KadesAIChatModal.tsx},
\texttt{src/app/api/ai/chat/route.ts} \\
\textbf{FR-11} & Sistem dapat mengelola data katalog buku (Tambah, Edit,
Hapus buku fisik dan e-book) & Must & Admin &
\texttt{src/app/admin/books/page.tsx},
\texttt{src/app/admin/books/new/page.tsx} \\
\textbf{FR-12} & Sistem dapat mengelola isi bab bacaan (\emph{chapter
management}) untuk setiap e-book & Must & Admin &
\texttt{src/app/admin/books/{[}id{]}/chapters/page.tsx},
\texttt{src/app/api/admin/chapters/route.ts} \\
\textbf{FR-13} & Sistem dapat mencatat transaksi peminjaman buku fisik
balai desa dengan tanggal pinjam dan batas kembali otomatis & Must &
Admin & \texttt{src/app/admin/circulation/page.tsx},
\texttt{src/app/api/admin/circulation/route.ts} \\
\textbf{FR-14} & Sistem dapat memproses perpanjangan masa pinjam (+7
hari) dan penandaan pengembalian buku fisik secara instan & Must & Admin
& \texttt{src/app/admin/circulation/page.tsx},
\texttt{src/app/api/admin/circulation/route.ts} (PATCH) \\
\textbf{FR-15} & Sistem dapat mengelola 8 kategori tematik desa (CRUD)
dengan pembaruan bersyarat (\emph{cascade update}) nama kategori buku &
Should & Admin & \texttt{src/app/admin/categories/page.tsx},
\texttt{src/app/api/admin/categories/route.ts} \\
\textbf{FR-16} & Sistem dapat mengelola data master 4 dusun resmi dan
menampilkan statistik sebaran warga per dusun & Should & Admin &
\texttt{src/app/admin/dusuns/page.tsx},
\texttt{src/app/api/admin/dusuns/route.ts} \\
\textbf{FR-17} & Sistem dapat mengelola akun warga dan menyediakan
fasilitas reset PIN sementara bagi warga yang lupa PIN & Must & Admin &
\texttt{src/app/admin/users/page.tsx},
\texttt{src/app/api/admin/users/route.ts} (PATCH reset-pin) \\
\textbf{FR-18} & Sistem dapat melakukan moderasi dan menghapus ulasan
yang mengandung konten tidak pantas & Should & Admin &
\texttt{src/app/admin/reviews/page.tsx},
\texttt{src/app/api/admin/reviews/route.ts} \\
\textbf{FR-19} & Sistem dapat menerbitkan, menyunting, dan menghapus
warta/maklumat literasi balai desa & Should & Admin &
\texttt{src/app/admin/announcements/page.tsx},
\texttt{src/app/api/admin/announcements/route.ts} \\
\textbf{FR-20} & Sistem dapat menyajikan dashboard analitik ringkasan
metrik literasi desa dan mengekspor laporan ke format CSV & Must & Admin
& \texttt{src/app/admin/page.tsx},
\texttt{src/app/admin/analytics/page.tsx}, \texttt{/api/admin/export} \\
\end{longtable}
}

\subsection{\texorpdfstring{4.4.2 Kebutuhan Non-Fungsional
(\emph{Non-Functional
Requirements})}{4.4.2 Kebutuhan Non-Fungsional (Non-Functional Requirements)}}\label{kebutuhan-non-fungsional-non-functional-requirements}

Spesifikasi kebutuhan non-fungsional dirumuskan secara terukur untuk
menjamin standar kualitas sistem: 1. \textbf{Keamanan
(\emph{Security})}: - Kode PIN pengguna dienkripsi searah menggunakan
algoritma \textbf{Bcrypt} dengan faktor biaya komputasi \emph{salt
rounds} bernilai 10. - Otentikasi sesi dikelola secara aman menggunakan
\emph{JSON Web Token} (JWT) yang tersimpan dalam \emph{HTTP-only,
SameSite Cookie} guna mencegah pencurian kredensial via serangan XSS dan
CSRF. - Hak akses rute administratif (\texttt{/admin/*}) dilindungi oleh
\emph{middleware} di lapisan edge; akses tanpa peran \texttt{ADMIN} akan
ditolak seketika (HTTP 403 / \emph{Redirect}). - Seluruh variabel
rahasia koneksi basis data disimpan dalam \emph{Vercel Environment
Variables} dan tidak terekspos ke klien (\emph{Zero Secrets Policy}). 2.
\textbf{Kinerja dan Kecepatan (\emph{Performance})}: - Waktu muat laman
pertama (\emph{First Contentful Paint / FCP}) di bawah 1.5 detik pada
jaringan seluler 4G desa. - Pemanfaatan \emph{React Server Components}
(RSC) meminimalkan ukuran \emph{bundle} JavaScript sisi klien di bawah
150 KB untuk halaman utama. - Optimasi gambar otomatis menggunakan
\emph{Next.js Image Optimizer} dengan konversi format WebP secara
dinamis. 3. \textbf{Ketersediaan dan Keandalan (\emph{Availability \&
Reliability})}: - Sistem beroperasi pada infrastruktur komputasi awan
\emph{Vercel Edge Network} dengan jaminan \emph{uptime} 99.9\%. - Basis
data Neon Serverless PostgreSQL mendukung mekanisme \emph{connection
pooling} berbasis WebSocket untuk menangani lonjakan koneksi simultan
dari gawai warga. 4. \textbf{Kemudahan Penggunaan (\emph{Usability})}: -
Antarmuka mengadopsi prinsip \emph{mobile-first} dengan tata letak
tombol navigasi yang mudah dijangkau ibu jari (\emph{thumb-zone
navigation}). - Menyediakan fitur alih bahasa instan ke Basa Sunda untuk
kenyamanan warga lokal lanjut usia. - Menyediakan panduan visual dan
indikator umpan balik (\emph{toast notification}) pada setiap aksi
simpan, ubah, dan hapus data.

\section{\texorpdfstring{4.5 Perancangan Alur Sistem (\emph{System
Flowchart})}{4.5 Perancangan Alur Sistem (System Flowchart)}}\label{perancangan-alur-sistem-system-flowchart}

Alur logika operasional sistem Pustaka Pangkalan membedakan akses antara
Pengunjung Tamu, Warga Terdaftar, dan Administrator Balai Desa
sebagaimana diilustrasikan pada Gambar 4.1.

\begin{verbatim}
                           +-------------------+
                           |  Pengguna Masuk   |
                           +-------------------+
                                     |
                                     v
                           /-------------------\
                          /    Apakah Sudah     \
                         <     Memiliki Akun?    >
                          \                    /
                           \------------------/
                             |              |
                       [Ya]  |              | [Tidak]
                             v              v
                     +---------------+  +---------------------+
                     | Form Login    |  | Form Registrasi     |
                     | (NIK & PIN)   |  | (NIK, Nama, Dusun)  |
                     +---------------+  +---------------------+
                             |                     |
                             +----------+----------+
                                        |
                                        v
                            /-----------------------\
                           /   Validasi Kredensial   \
                          <     Otentikasi Berhasil?  >
                           \                         /
                            \-----------------------/
                                 |              |
                           [Ya]  |              | [Tidak]
                                 v              v
                     /-----------------------\ +---------------------+
                    /    Periksa Peran        \| Tampilkan Pesan     |
                   <     Pengguna (Role)       >| Kesalahan Login     |
                    \                        / +---------------------+
                     \----------------------/
                         |              |
                 [USER]  |              | [ADMIN]
                         v              v
             +--------------------+ +--------------------+
             | Beranda Warga      | | Dashboard Admin    |
             | - Jelajah Katalog  | | - Metrik Literasi  |
             | - Baca Bab E-Book  | | - Sirkulasi Fisik  |
             | - Poin & Badge     | | - Kelola Buku/Bab  |
             | - Kartu Anggota QR | | - Reset PIN Warga  |
             | - Asisten Kades AI | | - Ekspor Laporan   |
             +--------------------+ +--------------------+
\end{verbatim}

\emph{Gambar 4.1. Diagram Alur Logika Utama Sistem Informasi Pustaka
Pangkalan}

\section{4.6 Pemodelan Berorientasi Objek
(UML)}\label{pemodelan-berorientasi-objek-uml}

\subsection{4.6.1 Use Case Diagram}\label{use-case-diagram}

Diagram Use Case memetakan 21 fungsionalitas sistem yang didistribusikan
kepada Aktor Warga Desa (\texttt{USER}) dan Pengelola Balai Desa
(\texttt{ADMIN}) sebagaimana disajikan pada Gambar 4.2.

\begin{verbatim}
+-----------------------------------------------------------------------------------+
|                        USE CASE DIAGRAM PUSTAKA PANGKALAN                         |
|                                                                                   |
|      +---------------+                                     +---------------+      |
|      |               |--- (UC-01: Registrasi Akun NIK) ----|               |      |
|      |               |--- (UC-02: Otentikasi Login NIK/PIN)|               |      |
|      |               |--- (UC-03: Jelajah & Cari Buku) ----|               |      |
|      |               |--- (UC-04: Baca Modul Bab E-Book) --|               |      |
|      |               |--- (UC-05: Simpan Bookmark & Progres|               |      |
|      |  WARGA DESA   |--- (UC-06: Kelola Rak Buku Pribadi) | ADMINISTRATOR |      |
|      |    (USER)     |--- (UC-07: Beri Ulasan & Rating) ---| BALAI DESA    |      |
|      |               |--- (UC-08: Lihat Kartu Anggota QR) -|    (ADMIN)    |      |
|      |               |--- (UC-09: Alih Bahasa Sunda/Indo) -|               |      |
|      |               |--- (UC-10: Konsultasi Kades AI) ----|               |      |
|      |               |--- (UC-11: Lihat Riwayat Poin) -----|               |      |
|      +---------------+                                     |               |      |
|                                                            |--- (UC-12: Kelola Katalog Buku)
|                                                            |--- (UC-13: Kelola Bab Buku)
|                                                            |--- (UC-14: Catat Pinjam Fisik)
|                                                            |--- (UC-15: Perpanjang Pinjaman)
|                                                            |--- (UC-16: Tandai Pengembalian)
|                                                            |--- (UC-17: Kelola Kategori)
|                                                            |--- (UC-18: Kelola Data Dusun)
|                                                            |--- (UC-19: Reset PIN Warga)
|                                                            |--- (UC-20: Kelola Warta Desa)
|                                                            |--- (UC-21: Ekspor Laporan CSV)
|                                                            +---------------+      |
+-----------------------------------------------------------------------------------+
\end{verbatim}

\emph{Gambar 4.2. Diagram Use Case Sistem Informasi Pustaka Pangkalan}

\subsection{\texorpdfstring{4.6.2 Narasi Skenario Use Case (\emph{Use
Case
Narrative})}{4.6.2 Narasi Skenario Use Case (Use Case Narrative)}}\label{narasi-skenario-use-case-use-case-narrative}

Berikut disajikan 4 (empat) narasi skenario use case terinci untuk
fungsi-fungsi kritikal sistem:

\subsubsection{Skenario 1: UC-01 Registrasi Akun Warga
Baru}\label{skenario-1-uc-01-registrasi-akun-warga-baru}

\begin{itemize}
\tightlist
\item
  \textbf{Aktor Utama}: Warga Desa (Calon Anggota)
\item
  \textbf{Deskripsi}: Warga mendaftarkan identitas diri ke dalam sistem
  menggunakan NIK KTP dan memilih dusun domisili.
\item
  \textbf{Prakondisi}: Warga belum terdaftar di basis data dan membuka
  halaman \texttt{/login}.
\item
  \textbf{Alur Utama (\emph{Main Flow})}:

  \begin{enumerate}
  \def\labelenumi{\arabic{enumi}.}
  \tightlist
  \item
    Warga menekan tombol \emph{Daftar Akun Baru}.
  \item
    Sistem menampilkan formulir registrasi yang memuat masukan: NIK (16
    digit), Nama Lengkap, Nomor Telepon, Pilihan Dusun (Pangkalan,
    Cikajang, Pasir Arangan, Pasir Gombong), dan PIN 6 digit.
  \item
    Warga mengisi seluruh kolom data dan menekan tombol \emph{Daftar
    Sekarang}.
  \item
    Sistem memvalidasi bahwa format NIK tepat 16 digit angka dan PIN
    tepat 6 digit angka.
  \item
    Sistem memeriksa keunikan NIK pada tabel \texttt{User}.
  \item
    Sistem mengenkripsi PIN menggunakan fungsi \emph{Bcrypt} (\emph{salt
    rounds} 10).
  \item
    Sistem menyimpan data warga baru dengan peran default \texttt{USER}.
  \item
    Sistem memberikan bonus 50 poin sambutan literasi dan mengarahkan
    pengguna langsung ke halaman beranda warga.
  \end{enumerate}
\item
  \textbf{Alur Alternatif (\emph{Alternative Flow})}:

  \begin{itemize}
  \tightlist
  \item
    \emph{4a. Format NIK atau PIN tidak valid}: Sistem menampilkan pesan
    peringatan bahwa NIK harus 16 digit angka dan PIN harus 6 digit
    angka.
  \item
    \emph{5a. NIK telah terdaftar}: Sistem menampilkan peringatan ``NIK
    sudah terdaftar dalam sistem. Silakan masuk menggunakan PIN Anda.''
  \end{itemize}
\item
  \textbf{Pasca-kondisi}: Akun warga aktif tercatat di basis data dan
  sesi login tersimpan.
\end{itemize}

\subsubsection{Skenario 2: UC-04 Pembacaan Bab E-Book dan Akumulasi
Poin}\label{skenario-2-uc-04-pembacaan-bab-e-book-dan-akumulasi-poin}

\begin{itemize}
\tightlist
\item
  \textbf{Aktor Utama}: Warga Desa Terotentikasi
\item
  \textbf{Deskripsi}: Warga membaca materi e-book secara bertahap per
  bab melalui antarmuka \emph{chapter reader}.
\item
  \textbf{Prakondisi}: Warga telah login dan memilih salah satu buku
  digital dari katalog.
\item
  \textbf{Alur Utama (\emph{Main Flow})}:

  \begin{enumerate}
  \def\labelenumi{\arabic{enumi}.}
  \tightlist
  \item
    Warga membuka halaman detail buku (\texttt{/books/{[}id{]}}) dan
    menekan salah satu bab bacaan.
  \item
    Sistem membuka antarmuka \emph{reader}
    (\texttt{/read/{[}chapterId{]}}) dengan memuat teks konten bab yang
    bersih dan bebas distraksi.
  \item
    Warga membaca isi materi modul hingga ke bagian akhir halaman.
  \item
    Warga menekan tombol \emph{Tandai Selesai Membaca} atau menekan
    tombol navigasi \emph{Bab Selanjutnya}.
  \item
    Sistem mencatat riwayat bacaan pada tabel \texttt{ReadingHistory}.
  \item
    Sistem menambahkan +10 poin literasi ke saldo poin pengguna pada
    tabel \texttt{User}.
  \item
    Sistem memperbarui data tanggal keaktifan berturut-turut
    (\emph{reading streak}).
  \item
    Sistem mengevaluasi apakah akumulasi poin memenuhi syarat perolehan
    lencana baru (\emph{badge}). Jika tercapai, sistem memunculkan
    animasi konfeti dan dialog selamat.
  \end{enumerate}
\item
  \textbf{Alur Alternatif (\emph{Alternative Flow})}:

  \begin{itemize}
  \tightlist
  \item
    \emph{5a. Bab telah pernah diselesaikan sebelumnya}: Sistem mencatat
    kemajuan bacaan namun tidak memberikan poin ganda untuk mencegah
    manipulasi poin.
  \end{itemize}
\item
  \textbf{Pasca-kondisi}: Progres membaca dan saldo poin pengguna
  diperbarui pada basis data.
\end{itemize}

\subsubsection{Skenario 3: UC-14 Pencatatan Sirkulasi Peminjaman Buku
Fisik Balai
Desa}\label{skenario-3-uc-14-pencatatan-sirkulasi-peminjaman-buku-fisik-balai-desa}

\begin{itemize}
\tightlist
\item
  \textbf{Aktor Utama}: Pengelola Balai Desa (\texttt{ADMIN})
\item
  \textbf{Deskripsi}: Pengelola mencatat peminjaman eksemplar buku cetak
  yang dilakukan oleh warga secara langsung di kantor balai desa.
\item
  \textbf{Prakondisi}: Pengelola telah masuk ke dashboard administratif
  (\texttt{/admin/circulation}).
\item
  \textbf{Alur Utama (\emph{Main Flow})}:

  \begin{enumerate}
  \def\labelenumi{\arabic{enumi}.}
  \tightlist
  \item
    Pengelola menekan tombol \emph{Catat Peminjaman Baru}.
  \item
    Sistem menampilkan modal formulir sirkulasi.
  \item
    Pengelola memilih nama warga peminjam (berdasarkan pencarian
    NIK/Nama) dan memilih judul buku fisik yang dipinjam.
  \item
    Pengelola mengisi catatan kondisi buku dan menekan tombol
    \emph{Simpan Peminjaman}.
  \item
    Sistem menetapkan \texttt{borrowDate} ke tanggal saat ini dan
    menghitung \texttt{dueDate} otomatis 7 hari ke depan.
  \item
    Sistem menetapkan status transaksi menjadi \texttt{BORROWED}.
  \item
    Sistem mengurangi jumlah stok buku fisik yang tersedia
    (\emph{availableCopies}) pada tabel \texttt{Book}.
  \item
    Sistem memperbarui daftar antrean sirkulasi aktif pada dashboard.
  \end{enumerate}
\item
  \textbf{Alur Alternatif (\emph{Alternative Flow})}:

  \begin{itemize}
  \tightlist
  \item
    \emph{7a. Stok buku fisik habis}: Sistem menolak transaksi dan
    memunculkan peringatan bahwa seluruh eksemplar buku sedang dipinjam.
  \end{itemize}
\item
  \textbf{Pasca-kondisi}: Transaksi sirkulasi tercatat di tabel
  \texttt{BookBorrowing} dan stok buku berkurang.
\end{itemize}

\subsubsection{Skenario 4: UC-19 Fasilitas Reset PIN Warga oleh
Pengelola
Desa}\label{skenario-4-uc-19-fasilitas-reset-pin-warga-oleh-pengelola-desa}

\begin{itemize}
\tightlist
\item
  \textbf{Aktor Utama}: Pengelola Balai Desa (\texttt{ADMIN})
\item
  \textbf{Deskripsi}: Pengelola mereset PIN warga yang lupa kode
  aksesnya menjadi PIN standar sementara.
\item
  \textbf{Prakondisi}: Warga melapor ke balai desa dengan membawa KTP
  fisik, dan pengelola membuka menu \texttt{/admin/users}.
\item
  \textbf{Alur Utama (\emph{Main Flow})}:

  \begin{enumerate}
  \def\labelenumi{\arabic{enumi}.}
  \tightlist
  \item
    Pengelola mencari data akun warga berdasarkan NIK atau Nama pada
    tabel pengguna.
  \item
    Pengelola menekan tombol \emph{Reset PIN}.
  \item
    Sistem memunculkan dialog konfirmasi dan menampilkan opsi pengaturan
    PIN sementara (standar: \texttt{123456}).
  \item
    Pengelola menekan tombol \emph{Konfirmasi Reset}.
  \item
    Sistem melakukan \emph{hashing} ulang PIN sementara menggunakan
    algoritma Bcrypt dan memperbarui kolom \texttt{pin} pada tabel
    \texttt{User}.
  \item
    Sistem menampilkan pesan sukses dan pengelola menginformasikan PIN
    sementara kepada warga untuk segera diganti setelah login.
  \end{enumerate}
\item
  \textbf{Pasca-kondisi}: PIN akun warga di basis data telah diperbarui
  dengan hash baru.
\end{itemize}

\subsection{4.6.3 Activity Diagram}\label{activity-diagram}

Alur kerja peminjaman dan pengembalian buku fisik di balai desa
dimodelkan pada Gambar 4.3.

\begin{verbatim}
      WARGA DESA                          SISTEM PUSTAKA PANGKALAN                  PENGELOLA DESA (ADMIN)
          |                                          |                                         |
          |--- Datang ke Balai Desa & Bawa Buku ---->|                                         |
          |                                          |<-- Buka Menu Sirkulasi (/admin) --------|
          |                                          |                                         |
          |                                          |<-- Pilih Warga & Input Judul Buku ------|
          |                                          |                                         |
          |                                   /---------------\                               |
          |                                  /  Stok Buku Ada? \                              |
          |                                 <                   >                              |
          |                                  \                 /                               |
          |                                   \---------------/                                |
          |                                      |           |                                 |
          |                                [Ya]  |           | [Tidak]                         |
          |                                      v           v                                 |
          |                               +---------------+ +---------------+                  |
          |                               | Kurangi Stok  | | Tolak Pinjam  |                  |
          |                               | Hitung +7 Hari| | Tampil Pesan  |                  |
          |                               | Status: Dipinjam| Habis Stok    |                  |
          |                               +---------------+ +---------------+                  |
          |                                      |                   |                         |
          |<-- Terima Buku Fisik Pinjaman -------+                   +---> Konfirmasi Warga ---|
          |                                      |                                             |
          |~~~~~~~~~~~~~~~~ Masa Pinjam Berjalan (Maks 7 Hari) ~~~~~~~~~~~~~~~~~~~~~~~~~~~~~~~~|
          |                                      |                                             |
          |--- Kembalikan Buku ke Balai Desa --->|                                             |
          |                                      |<-- Tekan Tombol 'Tandai Kembali' -----------|
          |                                      |                                             |
          |                               +---------------+                                    |
          |                               | Status: Kembali                                    |
          |                               | Kembalikan Stok                                    |
          |                               | Catat Tgl Selesai                                  |
          |                               +---------------+                                    |
          |                                      |                                             |
          |<-- Konfirmasi Pengembalian Selesai --+------------------------ Selesai ------------|
\end{verbatim}

\emph{Gambar 4.3. Activity Diagram Manajemen Sirkulasi Buku Fisik Balai
Desa}

\section{4.7 Perancangan Arsitektur Sistem 4-Tier
As-Built}\label{perancangan-arsitektur-sistem-4-tier-as-built}

Sistem informasi Pustaka Pangkalan dibangun dengan mengadopsi pola
arsitektur perangkat lunak 4-Tier modern yang terdistribusi secara
\emph{serverless} sebagaimana disajikan pada Gambar 4.4.

\begin{verbatim}
+-------------------------------------------------------------------------+
|                    ARSITEKTUR SISTEM 4-TIER AS-BUILT                    |
|                                                                         |
|  [TIER 1: PRESENTATION LAYER - PERAMBAN KLIEN]                          |
|  - Antarmuka Responsif Mobile-First (Tailwind CSS 4, Lucide Icons)      |
|  - React 19 Client Components (Interactive State, Form Handling, SWR)   |
|  - Komponen Dwi-Bahasa (LanguageContext: ID & SU)                       |
|                                    |                                    |
|                                    v HTTP/HTTPS                         |
|  [TIER 2: APPLICATION & API ROUTING LAYER - VERCEL EDGE]                |
|  - Next.js 16 App Router & React Server Components (RSC)                |
|  - Edge Middleware (Proteksi RBAC rute /admin & autentikasi sesi)       |
|  - 30 Endpoint RESTful API Route Handlers (/api/*)                      |
|  - NextAuth.js Credentials Engine (JWT Session Token Verification)      |
|                                    |                                    |
|                                    v TypeScript Query Calls             |
|  [TIER 3: DATA ACCESS & BUSINESS LOGIC LAYER - PRISMA ORM]              |
|  - Prisma ORM 7.9.1 Type-Safe Client                                    |
|  - Validasi Skema Zod & Hashing Kriptografis Bcrypt (Salt 10)           |
|  - Gamification Engine (Perhitungan Poin, Streak, Badge Checker)        |
|  - Kades AI Integration Handler                                         |
|                                    |                                    |
|                                    v Encrypted WebSocket Pool           |
|  [TIER 4: PERSISTENCE STORAGE LAYER - NEON POSTGRESQL]                  |
|  - Neon Serverless PostgreSQL Database Instance                         |
|  - 10 Relational Data Tables (Foreign Keys, Indexes, Cascades)          |
+-------------------------------------------------------------------------+
\end{verbatim}

\emph{Gambar 4.4. Diagram Arsitektur Perangkat Lunak 4-Tier As-Built}

Deskripsi fungsi setiap tier: 1. \textbf{Tier 1 (Presentation Layer)}:
Berjalan pada peramban web perangkat pengguna (ponsel atau desktop).
Menggunakan React 19 dan Tailwind CSS untuk menyajikan antarmuka visual
yang responsif, mendukung tema terang/gelap, serta mengelola status
interaktif lokal (\emph{local state}). 2. \textbf{Tier 2 (Application \&
Routing Layer)}: Berjalan pada infrastruktur \emph{Vercel Edge Network}.
Menangani perutean laman (\emph{App Router}), memproses \emph{React
Server Components} guna mengurangi transfer data ke klien, menerapkan
\emph{middleware} keamanan, serta menyediakan 30 endpoint API RESTful.
3. \textbf{Tier 3 (Data Access \& Business Logic Layer)}:
Mengimplementasikan logika bisnis inti aplikasi (penghitungan
denda/jatuh tempo sirkulasi, penambahan poin, pemberian lencana,
integrasi prompt AI, dan hashing kata sandi) menggunakan pustaka Prisma
ORM yang menjamin keamanan kueri tanpa risiko injeksi SQL. 4.
\textbf{Tier 4 (Persistence Layer)}: Berupa klaster basis data
relasional Neon Serverless PostgreSQL yang menyimpan seluruh rekaman
data secara permanen, terisolasi, dan terenkripsi pada media penyimpanan
awan.

\section{\texorpdfstring{4.8 Perancangan Basis Data (\emph{Database
Design})}{4.8 Perancangan Basis Data (Database Design)}}\label{perancangan-basis-data-database-design}

\subsection{4.8.1 Entity-Relationship Diagram (ERD)
As-Built}\label{entity-relationship-diagram-erd-as-built}

Struktur relasional data mencakup \textbf{10 entitas tabel} yang saling
terhubung melalui relasi kunci primer (\emph{Primary Key}) dan kunci
asing (\emph{Foreign Key}) sebagaimana ditunjukkan pada Gambar 4.5.

\begin{verbatim}
+--------------------+        +---------------------+        +--------------------+
|       User         | 1    * |    BookBorrowing    | *    1 |        Book        |
|--------------------|--------|---------------------|--------|--------------------|
| id (PK)            |        | id (PK)             |        | id (PK)            |
| nik (Unique)       |        | userId (FK)         |        | title              |
| name               |        | bookId (FK)         |        | author             |
| pin (Bcrypt Hash)  |        | borrowDate          |        | category           |
| role (USER/ADMIN)  |        | dueDate             |        | totalCopies        |
| dusun              |        | returnDate          |        | availableCopies    |
| points             |        | status              |        | isDigital          |
| readingStreak      |        | notes               |        +--------------------+
+--------------------+        +---------------------+                  | 1
       | 1                              | *                            |
       |                                |                              | *
       | *                              |                      +--------------------+
+--------------------+                  |                      |   ChapterContent   |
|   ReadingHistory   |                  |                      |--------------------|
|--------------------|                  |                      | id (PK)            |
| id (PK)            |                  |                      | bookId (FK)        |
| userId (FK)        |                  |                      | chapterNumber      |
| bookId (FK)        |                  |                      | title              |
| chapterId (FK)     |                  |                      | content (Text)     |
| completedAt        |                  |                      +--------------------+
+--------------------+                  |                              | 1
       | 1                              |                              |
       |                                |                              | *
       | *                              v                              |
+--------------------+        +---------------------+        +--------------------+
|       Review       |        |      Category       |        |    ReadingList     |
|--------------------|        |---------------------|        |--------------------|
| id (PK)            |        | id (PK)             |        | id (PK)            |
| userId (FK)        |        | name (Unique)       |        | userId (FK)        |
| bookId (FK)        |        | slug (Unique)       |        | bookId (FK)        |
| rating (Int 1-5)   |        | icon                |        | status             |
| comment            |        | description         |        +--------------------+
+--------------------+        +---------------------+
\end{verbatim}

\emph{Gambar 4.5. Entity-Relationship Diagram (ERD) As-Built Sistem
Pustaka Pangkalan}

\subsection{\texorpdfstring{4.8.2 Kamus Data Lengkap (\emph{Data
Dictionary})}{4.8.2 Kamus Data Lengkap (Data Dictionary)}}\label{kamus-data-lengkap-data-dictionary}

Berikut rincian spesifikasi struktur kolom dari 10 tabel basis data
terpasang:

\subsubsection{\texorpdfstring{1. Tabel \texttt{User} (Menyimpan Data
Identitas Akun
Pengguna)}{1. Tabel User (Menyimpan Data Identitas Akun Pengguna)}}\label{tabel-user-menyimpan-data-identitas-akun-pengguna}

\begin{itemize}
\tightlist
\item
  \texttt{id} (String, PK, CUID / UUID): Pengenal unik akun.
\item
  \texttt{nik} (String, Unique, Not Null): Nomor Induk Kependudukan 16
  digit.
\item
  \texttt{name} (String, Not Null): Nama lengkap warga.
\item
  \texttt{pin} (String, Not Null): Kode akses 6 digit yang telah
  di-\emph{hash} Bcrypt.
\item
  \texttt{role} (Enum \texttt{Role}, Not Null, Default \texttt{USER}):
  Hak akses akun (\texttt{USER} atau \texttt{ADMIN}).
\item
  \texttt{dusun} (String, Nullable): Wilayah dusun domisili warga di
  Desa Pangkalan.
\item
  \texttt{phone} (String, Nullable): Nomor telepon/WhatsApp aktif.
\item
  \texttt{points} (Integer, Not Null, Default 0): Akumulasi poin
  literasi.
\item
  \texttt{readingStreak} (Integer, Not Null, Default 0): Rekor hari
  membaca berturut-turut.
\item
  \texttt{lastReadDate} (DateTime, Nullable): Waktu terakhir membaca
  materi.
\item
  \texttt{badges} (String / Array, Nullable): Daftar lencana penghargaan
  yang diperoleh.
\item
  \texttt{createdAt} (DateTime, Not Null, Default \texttt{now()}): Waktu
  pendaftaran akun.
\item
  \texttt{updatedAt} (DateTime, Not Null): Waktu pembaruan data
  terakhir.
\end{itemize}

\subsubsection{\texorpdfstring{2. Tabel \texttt{Book} (Menyimpan
Metadata Katalog
Buku)}{2. Tabel Book (Menyimpan Metadata Katalog Buku)}}\label{tabel-book-menyimpan-metadata-katalog-buku}

\begin{itemize}
\tightlist
\item
  \texttt{id} (String, PK, CUID): Pengenal unik buku.
\item
  \texttt{title} (String, Not Null): Judul buku bacaan.
\item
  \texttt{author} (String, Not Null): Nama pengarang / penulis.
\item
  \texttt{publisher} (String, Nullable): Nama penerbit.
\item
  \texttt{year} (Integer, Nullable): Tahun publikasi cetak.
\item
  \texttt{isbn} (String, Nullable): Nomor ISBN buku (jika tersedia).
\item
  \texttt{category} (String, Not Null): Nama kategori tematik desa.
\item
  \texttt{description} (Text, Nullable): Ringkasan sinopsis materi buku.
\item
  \texttt{coverImage} (String, Nullable): Tautan URL berkas gambar
  sampul.
\item
  \texttt{totalCopies} (Integer, Not Null, Default 1): Jumlah total
  eksemplar fisik di balai desa.
\item
  \texttt{availableCopies} (Integer, Not Null, Default 1): Jumlah
  eksemplar fisik yang dapat dipinjam.
\item
  \texttt{isDigital} (Boolean, Not Null, Default true): Penanda apakah
  memiliki modul e-book online.
\item
  \texttt{viewsCount} (Integer, Not Null, Default 0): Frekuensi buku
  dibuka oleh pembaca.
\item
  \texttt{createdAt} (DateTime, Not Null, Default \texttt{now()}): Waktu
  penambahan ke sistem.
\end{itemize}

\subsubsection{\texorpdfstring{3. Tabel \texttt{ChapterContent}
(Menyimpan Isi Materi Bacaan per
Bab)}{3. Tabel ChapterContent (Menyimpan Isi Materi Bacaan per Bab)}}\label{tabel-chaptercontent-menyimpan-isi-materi-bacaan-per-bab}

\begin{itemize}
\tightlist
\item
  \texttt{id} (String, PK, CUID): Pengenal unik bab bacaan.
\item
  \texttt{bookId} (String, FK ke \texttt{Book.id}, Not Null, OnDelete
  Cascade): Keterkaitan dengan buku induk.
\item
  \texttt{chapterNumber} (Integer, Not Null): Nomor urutan bab bacaan
  (1, 2, 3, \ldots).
\item
  \texttt{title} (String, Not Null): Judul bab pembahasan.
\item
  \texttt{content} (Text, Not Null): Isi teks modul bacaan e-book.
\item
  \texttt{estimatedMinutes} (Integer, Not Null, Default 5): Estimasi
  waktu membaca dalam menit.
\item
  \texttt{createdAt} (DateTime, Not Null, Default \texttt{now()}): Waktu
  pembuatan bab.
\end{itemize}

\subsubsection{\texorpdfstring{4. Tabel \texttt{BookBorrowing}
(Menyimpan Transaksi Sirkulasi Fisik Balai
Desa)}{4. Tabel BookBorrowing (Menyimpan Transaksi Sirkulasi Fisik Balai Desa)}}\label{tabel-bookborrowing-menyimpan-transaksi-sirkulasi-fisik-balai-desa}

\begin{itemize}
\tightlist
\item
  \texttt{id} (String, PK, CUID): Pengenal unik transaksi sirkulasi.
\item
  \texttt{userId} (String, FK ke \texttt{User.id}, Not Null):
  Keterkaitan dengan warga peminjam.
\item
  \texttt{bookId} (String, FK ke \texttt{Book.id}, Not Null):
  Keterkaitan dengan buku yang dipinjam.
\item
  \texttt{borrowDate} (DateTime, Not Null, Default \texttt{now()}):
  Tanggal buku dibawa pulang.
\item
  \texttt{dueDate} (DateTime, Not Null): Tanggal batas waktu wajib
  kembali (+7 hari).
\item
  \texttt{returnDate} (DateTime, Nullable): Tanggal riil buku
  dikembalikan ke balai desa.
\item
  \texttt{status} (Enum \texttt{BorrowStatus}, Not Null, Default
  \texttt{BORROWED}): Status sirkulasi (\texttt{BORROWED},
  \texttt{RETURNED}, \texttt{OVERDUE}).
\item
  \texttt{notes} (String, Nullable): Catatan khusus kondisi fisik buku
  saat dipinjam/kembali.
\end{itemize}

\subsubsection{\texorpdfstring{5. Tabel \texttt{ReadingHistory}
(Menyimpan Catatan Riwayat Membaca
E-Book)}{5. Tabel ReadingHistory (Menyimpan Catatan Riwayat Membaca E-Book)}}\label{tabel-readinghistory-menyimpan-catatan-riwayat-membaca-e-book}

\begin{itemize}
\tightlist
\item
  \texttt{id} (String, PK, CUID): Pengenal unik riwayat baca.
\item
  \texttt{userId} (String, FK ke \texttt{User.id}, Not Null, OnDelete
  Cascade): Pengguna yang membaca.
\item
  \texttt{bookId} (String, FK ke \texttt{Book.id}, Not Null, OnDelete
  Cascade): Buku yang dibaca.
\item
  \texttt{chapterId} (String, FK ke \texttt{ChapterContent.id}, Not
  Null, OnDelete Cascade): Bab yang dibaca.
\item
  \texttt{completedAt} (DateTime, Not Null, Default \texttt{now()}):
  Waktu penuntasan membaca bab.
\end{itemize}

\subsubsection{\texorpdfstring{6. Tabel \texttt{Review} (Menyimpan
Ulasan dan Penilaian
Buku)}{6. Tabel Review (Menyimpan Ulasan dan Penilaian Buku)}}\label{tabel-review-menyimpan-ulasan-dan-penilaian-buku}

\begin{itemize}
\tightlist
\item
  \texttt{id} (String, PK, CUID): Pengenal unik ulasan.
\item
  \texttt{userId} (String, FK ke \texttt{User.id}, Not Null, OnDelete
  Cascade): Warga pemberi ulasan.
\item
  \texttt{bookId} (String, FK ke \texttt{Book.id}, Not Null, OnDelete
  Cascade): Buku yang dinilai.
\item
  \texttt{rating} (Integer, Not Null): Skor bintang bernilai antara 1
  hingga 5.
\item
  \texttt{comment} (Text, Not Null): Isi teks tanggapan atau ulasan
  warga.
\item
  \texttt{createdAt} (DateTime, Not Null, Default \texttt{now()}): Waktu
  pengiriman ulasan.
\end{itemize}

\subsubsection{\texorpdfstring{7. Tabel \texttt{Category} (Menyimpan
Master Kategori Tematik
Desa)}{7. Tabel Category (Menyimpan Master Kategori Tematik Desa)}}\label{tabel-category-menyimpan-master-kategori-tematik-desa}

\begin{itemize}
\tightlist
\item
  \texttt{id} (String, PK, CUID): Pengenal unik kategori.
\item
  \texttt{name} (String, Unique, Not Null): Nama kategori (misal:
  ``Pertanian \& Perkebunan'').
\item
  \texttt{slug} (String, Unique, Not Null): Pengenal URL ramah peramban
  (misal: ``pertanian'').
\item
  \texttt{icon} (String, Nullable): Nama ikon visual antarmuka.
\item
  \texttt{description} (String, Nullable): Keterangan deskriptif cakupan
  kategori.
\end{itemize}

\subsubsection{\texorpdfstring{8. Tabel \texttt{ReadingList} (Menyimpan
Rak Buku Favorit
Pribadi)}{8. Tabel ReadingList (Menyimpan Rak Buku Favorit Pribadi)}}\label{tabel-readinglist-menyimpan-rak-buku-favorit-pribadi}

\begin{itemize}
\tightlist
\item
  \texttt{id} (String, PK, CUID): Pengenal unik rekaman rak.
\item
  \texttt{userId} (String, FK ke \texttt{User.id}, Not Null, OnDelete
  Cascade): Pemilik rak.
\item
  \texttt{bookId} (String, FK ke \texttt{Book.id}, Not Null, OnDelete
  Cascade): Buku yang disimpan.
\item
  \texttt{status} (String, Not Null, Default ``SAVED''): Penanda status
  simpan.
\item
  \texttt{createdAt} (DateTime, Not Null, Default \texttt{now()}): Waktu
  penyimpanan ke rak.
\end{itemize}

\subsubsection{\texorpdfstring{9. Tabel \texttt{Account} (Menyimpan Akun
Eksternal Penyedia Otentikasi
NextAuth)}{9. Tabel Account (Menyimpan Akun Eksternal Penyedia Otentikasi NextAuth)}}\label{tabel-account-menyimpan-akun-eksternal-penyedia-otentikasi-nextauth}

\begin{itemize}
\tightlist
\item
  \texttt{id}, \texttt{userId}, \texttt{type}, \texttt{provider},
  \texttt{providerAccountId}, \texttt{refresh\_token},
  \texttt{access\_token}, \texttt{expires\_at}, \texttt{token\_type},
  \texttt{scope}, \texttt{id\_token}, \texttt{session\_state}.
\end{itemize}

\subsubsection{\texorpdfstring{10. Tabel \texttt{Session} dan
\texttt{VerificationToken} (Manajemen Sesi Otentikasi
NextAuth)}{10. Tabel Session dan VerificationToken (Manajemen Sesi Otentikasi NextAuth)}}\label{tabel-session-dan-verificationtoken-manajemen-sesi-otentikasi-nextauth}

\begin{itemize}
\tightlist
\item
  Menyimpan data token sesi aktif, tanggal kedaluwarsa, dan identitas
  pengguna terverifikasi.
\end{itemize}

\section{\texorpdfstring{4.9 Perancangan Antarmuka Pengguna (\emph{UI
Design
Principles})}{4.9 Perancangan Antarmuka Pengguna (UI Design Principles)}}\label{perancangan-antarmuka-pengguna-ui-design-principles}

Antarmuka pengguna Pustaka Pangkalan dirancang dengan mengutamakan
prinsip: 1. \textbf{Kejelasan Tipografi dan Hierarki Visual}:
Menggunakan jenis huruf antarmuka modern dengan ukuran teks tubuh
minimal 16 px pada ponsel agar nyaman dibaca oleh seluruh kelompok usia.
2. \textbf{Palet Warna Harmonis Bernuansa Alam Pedesaan}: - Warna Utama
(\emph{Primary}): Hijau Hutan (\emph{Forest Emerald} \texttt{\#059669} /
\texttt{\#065F46}), mencerminkan kesuburan tanah agraris Desa Pangkalan;
- Warna Aksen (\emph{Accent}): Emas Padi (\emph{Warm Amber}
\texttt{\#D97706}), melambangkan panen dan kemakmuran; - Warna Netral
(\emph{Neutral Surface}): Putih Gading (\texttt{\#F9FAFB}) dan Abu-abu
Lembut (\texttt{\#1F2937} pada tema gelap). 3. \textbf{Penyajian Mode
Pembaca Bebas Distraksi (\emph{Reader Mode})}: Pada halaman pembaca bab
(\texttt{/read/{[}chapterId{]}}), navigasi samping disembunyikan secara
otomatis agar pembaca dapat berkonsentrasi penuh pada teks modul bacaan.

\newpage

\chapter{BAB V --- IMPLEMENTASI SISTEM}\label{bab-v-implementasi-sistem}

\section{5.1 Lingkungan Implementasi Perangkat
Lunak}\label{lingkungan-implementasi-perangkat-lunak}

Implementasi Sistem Informasi Pustaka Pangkalan mengintegrasikan
dependensi pustaka \emph{open-source} modern yang terpasang pada berkas
konfigurasi \texttt{package.json} sebagaimana tercantum pada Tabel 5.1.

\emph{Tabel 5.1. Konfigurasi Lingkungan Perangkat Lunak Produksi}

{\def\LTcaptype{none} % do not increment counter
\begin{longtable}[]{@{}
  >{\raggedright\arraybackslash}p{(\linewidth - 6\tabcolsep) * \real{0.2143}}
  >{\raggedright\arraybackslash}p{(\linewidth - 6\tabcolsep) * \real{0.2143}}
  >{\centering\arraybackslash}p{(\linewidth - 6\tabcolsep) * \real{0.3571}}
  >{\raggedright\arraybackslash}p{(\linewidth - 6\tabcolsep) * \real{0.2143}}@{}}
\toprule\noalign{}
\begin{minipage}[b]{\linewidth}\raggedright
Komponen Perangkat Lunak
\end{minipage} & \begin{minipage}[b]{\linewidth}\raggedright
Nama Paket / Pustaka
\end{minipage} & \begin{minipage}[b]{\linewidth}\centering
Versi Terpasang
\end{minipage} & \begin{minipage}[b]{\linewidth}\raggedright
Peran dan Fungsi Teknis
\end{minipage} \\
\midrule\noalign{}
\endhead
\bottomrule\noalign{}
\endlastfoot
\textbf{Kerangka Kerja Utama} & \texttt{next} & \texttt{16.3.0} &
Kerangka kerja web full-stack React dengan arsitektur App Router \\
\textbf{Pustaka Antarmuka} & \texttt{react}, \texttt{react-dom} &
\texttt{19.0.0} & Pustaka antarmuka pengguna berbasis komponen
deklaratif \\
\textbf{Bahasa Pemrograman} & \texttt{typescript} & \texttt{\^{}5.0.0} &
Kompiler bahasa dengan sistem pengetikan statis berkeamanan tipe \\
\textbf{Styling \& Tata Letak} & \texttt{tailwindcss} &
\texttt{\^{}4.0.0} & Mesin penyusun gaya antarmuka utilitas CSS
modern \\
\textbf{Object-Relational Mapping} & \texttt{@prisma/client},
\texttt{prisma} & \texttt{7.9.1} & Klien akses dan kueri basis data
relasional PostgreSQL \\
\textbf{Driver Basis Data Serverless} &
\texttt{@neondatabase/serverless} & Terpasang & Driver koneksi basis
data serverless Neon via WebSocket \\
\textbf{Mesin Otentikasi Sesi} & \texttt{next-auth} &
\texttt{\^{}4.24.0} & Pengelola sesi login berbasis JSON Web Token
(JWT) \\
\textbf{Kriptografi Kata Sandi} & \texttt{bcryptjs} & \texttt{\^{}2.4.3}
& Pustaka hashing kata sandi satu arah algoritma Blowfish \\
\textbf{Ikonografi Visual} & \texttt{lucide-react} &
\texttt{\^{}0.400.0} & Koleksi simbol ikon vektor antarmuka pengguna \\
\textbf{Animasi Gamifikasi} & \texttt{canvas-confetti} &
\texttt{\^{}1.9.0} & Mesin partikel konfeti visual saat warga meraih
lencana \\
\textbf{Platform Cloud Hosting} & Vercel Edge Serverless & Node 20 LTS &
Infrastruktur komputasi edge dan deployment CI/CD otomatis \\
\end{longtable}
}

\section{5.2 Implementasi Basis Data dan Skema
Prisma}\label{implementasi-basis-data-dan-skema-prisma}

Skema basis data relasional dikonfigurasikan secara deklaratif di dalam
berkas \texttt{prisma/schema.prisma}. Potongan kode sumber berikut
memperlihatkan definisi entitas inti \texttt{User}, \texttt{Book}, dan
\texttt{BookBorrowing} yang menjamin integritas referensial data:

\begin{Shaded}
\begin{Highlighting}[]
\NormalTok{datasource db \{}
\NormalTok{  provider = "postgresql"}
\NormalTok{  url      = env("DATABASE\_URL")}
\NormalTok{\}}

\NormalTok{generator client \{}
\NormalTok{  provider = "prisma{-}client{-}js"}
\NormalTok{\}}

\NormalTok{enum Role \{}
\NormalTok{  USER}
\NormalTok{  ADMIN}
\NormalTok{\}}

\NormalTok{enum BorrowStatus \{}
\NormalTok{  BORROWED}
\NormalTok{  RETURNED}
\NormalTok{  OVERDUE}
\NormalTok{\}}

\NormalTok{model User \{}
\NormalTok{  id            String          @id @default(cuid())}
\NormalTok{  nik           String          @unique}
\NormalTok{  name          String}
\NormalTok{  pin           String          // Bcrypt Hash (Salt 10)}
\NormalTok{  role          Role            @default(USER)}
\NormalTok{  dusun         String?}
\NormalTok{  phone         String?}
\NormalTok{  points        Int             @default(0)}
\NormalTok{  readingStreak Int             @default(0)}
\NormalTok{  lastReadDate  DateTime?}
\NormalTok{  borrowings    BookBorrowing[]}
\NormalTok{  reviews       Review[]}
\NormalTok{  history       ReadingHistory[]}
\NormalTok{  readingList   ReadingList[]}
\NormalTok{  createdAt     DateTime        @default(now())}
\NormalTok{  updatedAt     DateTime        @updatedAt}
\NormalTok{\}}

\NormalTok{model Book \{}
\NormalTok{  id              String          @id @default(cuid())}
\NormalTok{  title           String}
\NormalTok{  author          String}
\NormalTok{  category        String}
\NormalTok{  description     String?}
\NormalTok{  coverImage      String?}
\NormalTok{  totalCopies     Int             @default(1)}
\NormalTok{  availableCopies Int             @default(1)}
\NormalTok{  isDigital       Boolean         @default(true)}
\NormalTok{  viewsCount      Int             @default(0)}
\NormalTok{  chapters        ChapterContent[]}
\NormalTok{  borrowings      BookBorrowing[]}
\NormalTok{  reviews         Review[]}
\NormalTok{  createdAt       DateTime        @default(now())}
\NormalTok{\}}

\NormalTok{model BookBorrowing \{}
\NormalTok{  id          String        @id @default(cuid())}
\NormalTok{  userId      String}
\NormalTok{  bookId      String}
\NormalTok{  borrowDate  DateTime      @default(now())}
\NormalTok{  dueDate     DateTime}
\NormalTok{  returnDate  DateTime?}
\NormalTok{  status      BorrowStatus  @default(BORROWED)}
\NormalTok{  notes       String?}
\NormalTok{  user        User          @relation(fields: [userId], references: [id], onDelete: Cascade)}
\NormalTok{  book        Book          @relation(fields: [bookId], references: [id], onDelete: Cascade)}
\NormalTok{\}}
\end{Highlighting}
\end{Shaded}

Skema tersebut kemudian dieksekusi ke dalam klaster Neon Serverless
PostgreSQL menggunakan perintah migrasi \texttt{npx\ prisma\ db\ push}
untuk membentuk tabel fisik beserta indeks pencariannya.

\section{\texorpdfstring{5.3 Implementasi Rute Antarmuka Pengguna
(\emph{Frontend
Routes})}{5.3 Implementasi Rute Antarmuka Pengguna (Frontend Routes)}}\label{implementasi-rute-antarmuka-pengguna-frontend-routes}

Sistem mengimplementasikan \textbf{19 rute halaman antarmuka pengguna}
yang terstruktur ke dalam modul publik/warga dan modul administratif
Balai Desa sebagaimana dirinci pada Tabel 5.2.

\emph{Tabel 5.2. Inventarisasi 19 Rute Halaman Antarmuka Pengguna
Terpasang}

{\def\LTcaptype{none} % do not increment counter
\begin{longtable}[]{@{}
  >{\centering\arraybackslash}p{(\linewidth - 8\tabcolsep) * \real{0.2632}}
  >{\raggedright\arraybackslash}p{(\linewidth - 8\tabcolsep) * \real{0.1579}}
  >{\raggedright\arraybackslash}p{(\linewidth - 8\tabcolsep) * \real{0.1579}}
  >{\centering\arraybackslash}p{(\linewidth - 8\tabcolsep) * \real{0.2632}}
  >{\raggedright\arraybackslash}p{(\linewidth - 8\tabcolsep) * \real{0.1579}}@{}}
\toprule\noalign{}
\begin{minipage}[b]{\linewidth}\centering
No
\end{minipage} & \begin{minipage}[b]{\linewidth}\raggedright
Jalur Rute URL (\emph{Path})
\end{minipage} & \begin{minipage}[b]{\linewidth}\raggedright
Berkas Kode Sumber Komponen
\end{minipage} & \begin{minipage}[b]{\linewidth}\centering
Peruntukan Aktor
\end{minipage} & \begin{minipage}[b]{\linewidth}\raggedright
Deskripsi Fungsi Antarmuka
\end{minipage} \\
\midrule\noalign{}
\endhead
\bottomrule\noalign{}
\endlastfoot
1 & \texttt{/} & \texttt{src/app/page.tsx} & Semua Aktor & Beranda
utama, banner literasi, katalog rekomendasi, warta desa \\
2 & \texttt{/explore} & \texttt{src/app/explore/page.tsx} & Semua Aktor
& Penjelajahan katalog buku dengan pencarian teks dan filter 8
kategori \\
3 & \texttt{/books/{[}id{]}} & \texttt{src/app/books/{[}id{]}/page.tsx}
& Semua Aktor & Detail metadata buku, daftar bab e-book, ulasan, dan
tombol aksi \\
4 & \texttt{/read/{[}chapterId{]}} &
\texttt{src/app/read/{[}chapterId{]}/page.tsx} & Warga (\texttt{USER}) &
Pembaca modul bab e-book bebas distraksi, navigasi, dan penanda \\
5 & \texttt{/shelf} & \texttt{src/app/shelf/page.tsx} & Warga
(\texttt{USER}) & Pengelolaan koleksi buku favorit dan riwayat bacaan
pribadi \\
6 & \texttt{/profile} & \texttt{src/app/profile/page.tsx} & Warga
(\texttt{USER}) & Kartu Anggota Digital ber-QR Code, poin, lencana, dan
edit profil \\
7 & \texttt{/login} & \texttt{src/app/login/page.tsx} & Tamu
(\texttt{GUEST}) & Formulir masuk sistem menggunakan NIK dan PIN 6
digit \\
8 & \texttt{/onboarding} & \texttt{src/app/onboarding/page.tsx} & Tamu
(\texttt{GUEST}) & Formulir pendaftaran warga baru dan pemilihan dusun
domisili \\
9 & \texttt{/leaderboard} & \texttt{src/app/leaderboard/page.tsx} &
Warga (\texttt{USER}) & Papan peringkat warga teraktif membaca dan
statistik dusun \\
10 & \texttt{/announcements} & \texttt{src/app/announcements/page.tsx} &
Semua Aktor & Daftar maklumat dan berita literasi resmi dari pemerintah
desa \\
11 & \texttt{/admin} & \texttt{src/app/admin/page.tsx} & Admin
(\texttt{ADMIN}) & Dashboard ikhtisar analitik: total buku, warga,
sirkulasi aktif \\
12 & \texttt{/admin/books} & \texttt{src/app/admin/books/page.tsx} &
Admin (\texttt{ADMIN}) & Manajemen inventaris katalog buku (tabel data,
filter, aksi hapus) \\
13 & \texttt{/admin/books/new} &
\texttt{src/app/admin/books/new/page.tsx} & Admin (\texttt{ADMIN}) &
Formulir pendaftaran buku baru (koleksi fisik maupun digital) \\
14 & \texttt{/admin/books/{[}id{]}/chapters} &
\texttt{src/app/admin/books/{[}id{]}/chapters/page.tsx} & Admin
(\texttt{ADMIN}) & Manajemen bab bacaan e-book (\emph{chapter content
editor}) \\
15 & \texttt{/admin/circulation} &
\texttt{src/app/admin/circulation/page.tsx} & Admin (\texttt{ADMIN}) &
Modul sirkulasi fisik balai desa: catat pinjam, +7 hari, tandai
kembali \\
16 & \texttt{/admin/categories} &
\texttt{src/app/admin/categories/page.tsx} & Admin (\texttt{ADMIN}) &
Manajemen 8 kategori tematik desa dengan \emph{cascade update} buku \\
17 & \texttt{/admin/dusuns} & \texttt{src/app/admin/dusuns/page.tsx} &
Admin (\texttt{ADMIN}) & Manajemen 4 dusun resmi dan pemantauan sebaran
warga per dusun \\
18 & \texttt{/admin/users} & \texttt{src/app/admin/users/page.tsx} &
Admin (\texttt{ADMIN}) & Manajemen data warga terdaftar dan fasilitas
Reset PIN warga \\
19 & \texttt{/admin/reviews} & \texttt{src/app/admin/reviews/page.tsx} &
Admin (\texttt{ADMIN}) & Moderasi ulasan warga dan penghapusan komentar
tidak pantas \\
\end{longtable}
}

\section{\texorpdfstring{5.4 Implementasi Antarmuka Pemrograman Aplikasi
(\emph{RESTful API
Endpoints})}{5.4 Implementasi Antarmuka Pemrograman Aplikasi (RESTful API Endpoints)}}\label{implementasi-antarmuka-pemrograman-aplikasi-restful-api-endpoints}

Sistem menyediakan \textbf{30 endpoint RESTful API} yang dibangun di
bawah direktori \texttt{src/app/api/} untuk melayani komunikasi data
asinkronus antara antarmuka pengguna dan basis data. Daftar endpoint
disajikan pada Tabel 5.3.

\emph{Tabel 5.3. Daftar 30 Endpoint RESTful API Terpasang}

{\def\LTcaptype{none} % do not increment counter
\begin{longtable}[]{@{}
  >{\centering\arraybackslash}p{(\linewidth - 8\tabcolsep) * \real{0.2381}}
  >{\raggedright\arraybackslash}p{(\linewidth - 8\tabcolsep) * \real{0.1429}}
  >{\centering\arraybackslash}p{(\linewidth - 8\tabcolsep) * \real{0.2381}}
  >{\centering\arraybackslash}p{(\linewidth - 8\tabcolsep) * \real{0.2381}}
  >{\raggedright\arraybackslash}p{(\linewidth - 8\tabcolsep) * \real{0.1429}}@{}}
\toprule\noalign{}
\begin{minipage}[b]{\linewidth}\centering
No
\end{minipage} & \begin{minipage}[b]{\linewidth}\raggedright
Jalur Endpoint API (\emph{Route})
\end{minipage} & \begin{minipage}[b]{\linewidth}\centering
Metode HTTP
\end{minipage} & \begin{minipage}[b]{\linewidth}\centering
Hak Akses
\end{minipage} & \begin{minipage}[b]{\linewidth}\raggedright
Fungsi Operasional API
\end{minipage} \\
\midrule\noalign{}
\endhead
\bottomrule\noalign{}
\endlastfoot
1 & \texttt{/api/auth/{[}...nextauth{]}} & GET, POST & Publik &
Pengendali sesi otentikasi NextAuth kredensial NIK \& PIN \\
2 & \texttt{/api/auth/register} & POST & Publik & Mendaftarkan warga
baru dengan hashing PIN Bcrypt \\
3 & \texttt{/api/books} & GET, POST & Publik / Admin & Mengambil katalog
buku publik atau menambah buku baru \\
4 & \texttt{/api/books/{[}id{]}} & GET, PUT, DELETE & Publik / Admin &
Mengambil detail, memperbarui metadata, atau menghapus buku \\
5 & \texttt{/api/categories} & GET & Publik & Mengambil daftar 8
kategori tematik aktif \\
6 & \texttt{/api/read/{[}chapterId{]}} & GET & Warga (\texttt{USER}) &
Mengambil isi teks bab e-book untuk antarmuka \emph{reader} \\
7 & \texttt{/api/reading-progress} & GET, POST & Warga (\texttt{USER}) &
Mengambil atau memperbarui progres halaman bab buku \\
8 & \texttt{/api/bookmarks} & GET, POST, DELETE & Warga (\texttt{USER})
& Mengelola penanda bab bacaan favorit warga \\
9 & \texttt{/api/shelf} & GET, POST, DELETE & Warga (\texttt{USER}) &
Mengelola rak simpan buku pribadi warga \\
10 & \texttt{/api/reviews} & GET, POST & Publik / Warga & Mengambil
daftar ulasan atau mengirim ulasan rating buku \\
11 & \texttt{/api/user/profile} & GET, PUT & Warga (\texttt{USER}) &
Mengambil data kartu anggota QR atau memperbarui profil \\
12 & \texttt{/api/user/change-pin} & POST & Warga (\texttt{USER}) &
Mengubah PIN lama menjadi PIN baru setelah diverifikasi \\
13 & \texttt{/api/ai/chat} & POST & Warga (\texttt{USER}) & Mengirim
pesan prompt dan menerima balasan asisten \emph{Kades AI} \\
14 & \texttt{/api/announcements} & GET & Publik & Mengambil daftar warta
maklumat literasi desa \\
15 & \texttt{/api/leaderboard} & GET & Publik / Warga & Mengambil daftar
peringkat warga teraktif dan metrik dusun \\
16 & \texttt{/api/admin/stats} & GET & Admin (\texttt{ADMIN}) &
Mengambil statistik agregasi metrik dashboard admin \\
17 & \texttt{/api/admin/books} & GET, POST & Admin (\texttt{ADMIN}) &
Manajemen inventaris katalog buku administratif \\
18 & \texttt{/api/admin/books/{[}id{]}} & GET, PUT, DELETE & Admin
(\texttt{ADMIN}) & Pembaruan data dan penghapusan buku administratif \\
19 & \texttt{/api/admin/chapters} & GET, POST & Admin (\texttt{ADMIN}) &
Mengambil daftar bab atau menambah bab bacaan e-book baru \\
20 & \texttt{/api/admin/chapters/{[}id{]}} & PUT, DELETE & Admin
(\texttt{ADMIN}) & Memperbarui isi teks bab atau menghapus bab e-book \\
21 & \texttt{/api/admin/circulation} & GET, POST, PATCH & Admin
(\texttt{ADMIN}) & Mengambil, mencatat pinjam fisik, +7 hari, \& tandai
kembali \\
22 & \texttt{/api/admin/categories} & GET, POST, PUT, DELETE & Admin
(\texttt{ADMIN}) & CRUD master kategori desa dengan \emph{cascade
update} buku \\
23 & \texttt{/api/admin/dusuns} & GET, POST, PUT, DELETE & Admin
(\texttt{ADMIN}) & CRUD master 4 dusun resmi dan sinkronisasi warga \\
24 & \texttt{/api/admin/users} & GET, PATCH & Admin (\texttt{ADMIN}) &
Mengambil daftar warga \& eksekusi Reset PIN sementara \\
25 & \texttt{/api/admin/reviews} & GET, DELETE & Admin (\texttt{ADMIN})
& Mengambil seluruh ulasan warga dan moderasi hapus \\
26 & \texttt{/api/admin/announcements} & GET, POST, PUT, DELETE & Admin
(\texttt{ADMIN}) & CRUD maklumat warta literasi Balai Desa Pangkalan \\
27 & \texttt{/api/admin/analytics} & GET & Admin (\texttt{ADMIN}) &
Mengambil data analitik grafik tren bacaan per kategori \\
28 & \texttt{/api/admin/export} & GET & Admin (\texttt{ADMIN}) &
Mengekspor seluruh rekaman sirkulasi \& warga ke format CSV \\
29 & \texttt{/api/health} & GET & Publik & Endpoint pemantauan status
kesehatan server dan database \\
30 & \texttt{/api/og} & GET & Publik & Generator gambar thumbnail
dinamis resmi OpenGraph media sosial \\
\end{longtable}
}

\section{5.5 Implementasi Mekanisme Keamanan dan
Autentikasi}\label{implementasi-mekanisme-keamanan-dan-autentikasi}

Mekanisme keamanan sistem dibangun dengan menerapkan prinsip
\emph{Defense-in-Depth}: 1. \textbf{Validasi Masukan (\emph{Input
Validation})}: Pada rute \texttt{/api/auth/register}, sistem memvalidasi
bahwa NIK terdiri dari tepat 16 karakter numerik
(\texttt{\^{}\textbackslash{}d\{16\}\$}) dan PIN terdiri dari tepat 6
digit numerik (\texttt{\^{}\textbackslash{}d\{6\}\$}). 2.
\textbf{Kriptografi Password Hashing}: PIN dienkripsi menggunakan fungsi
\texttt{bcrypt.hash(pin,\ 10)}. Potongan kode verifikasi pada NextAuth:

\begin{Shaded}
\begin{Highlighting}[]
\KeywordTok{const}\NormalTok{ isPinValid }\OperatorTok{=} \ControlFlowTok{await}\NormalTok{ bcrypt}\OperatorTok{.}\FunctionTok{compare}\NormalTok{(credentials}\OperatorTok{.}\AttributeTok{pin}\OperatorTok{,}\NormalTok{ user}\OperatorTok{.}\AttributeTok{pin}\NormalTok{)}\OperatorTok{;}
\ControlFlowTok{if}\NormalTok{ (}\OperatorTok{!}\NormalTok{isPinValid) \{}
  \ControlFlowTok{throw} \KeywordTok{new} \BuiltInTok{Error}\NormalTok{(}\StringTok{"NIK atau PIN yang Anda masukkan salah."}\NormalTok{)}\OperatorTok{;}
\NormalTok{\}}
\end{Highlighting}
\end{Shaded}

\begin{enumerate}
\def\labelenumi{\arabic{enumi}.}
\setcounter{enumi}{2}
\tightlist
\item
  \textbf{Kontrol Akses Berbasis Peran (\emph{Role-Based Access Control
  / RBAC})}: Seluruh rute administratif diproteksi oleh berkas
  \texttt{src/middleware.ts} yang memeriksa keberadaan token JWT dan
  klaim peran:
\end{enumerate}

\begin{Shaded}
\begin{Highlighting}[]
\ControlFlowTok{if}\NormalTok{ (request}\OperatorTok{.}\AttributeTok{nextUrl}\OperatorTok{.}\AttributeTok{pathname}\OperatorTok{.}\FunctionTok{startsWith}\NormalTok{(}\StringTok{"/admin"}\NormalTok{)) \{}
  \KeywordTok{const}\NormalTok{ token }\OperatorTok{=} \ControlFlowTok{await} \FunctionTok{getToken}\NormalTok{(\{ req}\OperatorTok{:}\NormalTok{ request \})}\OperatorTok{;}
  \ControlFlowTok{if}\NormalTok{ (}\OperatorTok{!}\NormalTok{token }\OperatorTok{||}\NormalTok{ token}\OperatorTok{.}\AttributeTok{role} \OperatorTok{!==} \StringTok{"ADMIN"}\NormalTok{) \{}
    \ControlFlowTok{return}\NormalTok{ NextResponse}\OperatorTok{.}\FunctionTok{redirect}\NormalTok{(}\KeywordTok{new} \FunctionTok{URL}\NormalTok{(}\StringTok{"/login"}\OperatorTok{,}\NormalTok{ request}\OperatorTok{.}\AttributeTok{url}\NormalTok{))}\OperatorTok{;}
\NormalTok{  \}}
\NormalTok{\}}
\end{Highlighting}
\end{Shaded}

\begin{enumerate}
\def\labelenumi{\arabic{enumi}.}
\setcounter{enumi}{3}
\tightlist
\item
  \textbf{Perlindungan Terhadap Akses Laman Login Saat Sesi Aktif}:
  Berdasarkan umpan balik perbaikan sistem, apabila pengguna yang telah
  berhasil masuk (\emph{logged in}) secara sengaja atau tidak sengaja
  mengetikkan URL \texttt{/login}, \emph{middleware} secara otomatis
  mengalihkan (\emph{redirect}) pengguna kembali ke beranda \texttt{/}
  tanpa menampilkan ulang formulir login.
\end{enumerate}

\section{5.6 Implementasi Fitur-Fitur
Khusus}\label{implementasi-fitur-fitur-khusus}

\subsection{5.6.1 Fitur Gamifikasi
Literasi}\label{fitur-gamifikasi-literasi}

Sistem mengimplementasikan mesin gamifikasi literasi untuk menumbuhkan
motivasi membaca warga: - \textbf{Poin Literasi}: Warga memperoleh 50
poin saat pertama kali mendaftar akun, dan memperoleh 10 poin setiap
kali menyelesaikan pembacaan 1 (satu) bab e-book. - \textbf{Reading
Streak}: Menghitung jumlah hari berturut-turut warga aktif membaca
materi di perpustakaan digital. - \textbf{Lencana Prestasi
(\emph{Badges})}: Diberikan secara otomatis saat akumulasi poin mencapai
ambang batas: - Lencana \textbf{Pembaca Rajin} (Ambang batas: 100 Poin);
- Lencana \textbf{Cendekia Desa} (Ambang batas: 300 Poin); - Lencana
\textbf{Pelopor Literasi} (Ambang batas: 500 Poin). Pencapaian lencana
memicu pemanggilan animasi konfeti visual \texttt{canvas-confetti} dan
tercantum secara resmi pada Kartu Anggota Digital ber-QR Code.

\subsection{\texorpdfstring{5.6.2 Asisten Literasi Cerdas (\emph{Kades
AI})}{5.6.2 Asisten Literasi Cerdas (Kades AI)}}\label{asisten-literasi-cerdas-kades-ai}

Fitur \emph{Kades AI} diimplementasikan melalui komponen modal
percakapan interaktif (\texttt{src/components/KadesAIChatModal.tsx}) dan
endpoint backend \texttt{/api/ai/chat}. Asisten virtual ini
dikonfigurasikan dengan \emph{system prompt} berkarakter Kepala Desa
Pangkalan yang bijak, ramah, menguasai potensi 4 dusun, serta siap
memberikan rekomendasi buku terapan seputar pertanian, peternakan, dan
wirausaha desa.

\subsection{5.6.3 Lokalisasi Dwi-Bahasa (Bahasa Indonesia \& Basa
Sunda)}\label{lokalisasi-dwi-bahasa-bahasa-indonesia-basa-sunda}

Untuk menghormati dan melestarikan budaya lokal masyarakat Kecamatan
Cikidang, sistem dilengkapi penyedia bahasa
\texttt{LanguageProvider.tsx} yang menyimpan kamus terjemahan untuk
seluruh label navigasi, tombol, dan pesan antarmuka: - Label Beranda
\(ightarrow\) \emph{Tepas} - Label Jelajah Katalog \(ightarrow\)
\emph{Kotretan Buku / Papay} - Label Rak Buku \(ightarrow\) \emph{Rak
Buku Abdi} - Label Profil \(ightarrow\) \emph{Profil Warga} - Label Baca
Sekarang \(ightarrow\) \emph{Maca Ayeuna} Pengguna dapat mengubah bahasa
secara instan melalui tombol saklar bahasa pada bilah navigasi atas
(\emph{Top App Bar}).

\subsection{5.6.4 Penyelarasan Branding Resmi Kabupaten
Sukabumi}\label{penyelarasan-branding-resmi-kabupaten-sukabumi}

Sistem menyelaraskan identitas visual resmi Pemerintah Daerah Kabupaten
Sukabumi: - Memasang aset lambang daerah resmi
\texttt{Lambang\_Kab\_Sukabumi.svg} dan format terkompresi
\texttt{.webp} pada bilah atas aplikasi (\emph{Top App Bar}) dan ikon
tab peramban (\emph{favicon}); - Membangun generator gambar OpenGraph
dinamis (\texttt{/api/og}) yang menampilkan judul aplikasi \emph{Pustaka
Pangkalan}, identitas wilayah \textbf{Desa Pangkalan, Kecamatan
Cikidang, Kabupaten Sukabumi}, lambang daerah Kabupaten Sukabumi, serta
alamat domain resmi \url{https://perpus-pangkalan.vercel.app} saat
tautan dibagikan di media sosial.

\section{5.7 Deployment dan Konfigurasi
Produksi}\label{deployment-dan-konfigurasi-produksi}

Sistem dideploy secara otomatis pada infrastruktur \emph{Vercel
Serverless Edge Platform}. Konfigurasi variabel lingkungan diatur secara
ketat pada panel kontrol Vercel: - \texttt{DATABASE\_URL}: Tautan
koneksi terenkripsi SSL ke basis data Neon Serverless PostgreSQL. -
\texttt{NEXTAUTH\_SECRET}: Kunci enkripsi simetris untuk penandatanganan
token JWT sesi. - \texttt{NEXTAUTH\_URL}: Alamat domain resmi
\texttt{https://perpus-pangkalan.vercel.app}. Pembaruan kode pada cabang
\texttt{main} repositori GitHub secara otomatis memicu proses
\emph{build}, pengecekan kompilasi TypeScript, dan penyebaran instan ke
jaringan edge global tanpa \emph{downtime}.

\newpage

\chapter{BAB VI --- PENGUJIAN DAN EVALUASI
SISTEM}\label{bab-vi-pengujian-dan-evaluasi-sistem}

\section{6.1 Rencana dan Strategi
Pengujian}\label{rencana-dan-strategi-pengujian}

Penjaminan mutu perangkat lunak (\emph{Software Quality Assurance})
dilakukan secara ketat untuk membuktikan bahwa sistem terpasang bebas
dari cacat logika, aman, dan bekerja sesuai spesifikasi fungsional yang
telah ditetapkan.

Metodologi pengujian mencakup: 1. \textbf{Pengujian Kotak Hitam
(\emph{Black-Box Testing})}: Memvalidasi kesesuaian masukan dan keluaran
antarmuka pengguna tanpa memandang struktur internal kode. 2.
\textbf{Pengujian Nilai Batas (\emph{Boundary Value Analysis})}: Menguji
kondisi batas ekstrem, seperti masukan NIK dengan panjang kurang dari 16
digit, lebih dari 16 digit, serta karakter non-angka. 3.
\textbf{Pengujian Partisi Ekivalensi (\emph{Equivalence Partitioning})}:
Mengelompokkan jenis masukan valid dan tidak valid pada formulir login,
registrasi, sirkulasi, dan CRUD buku. 4. \textbf{Pengujian Regresi
Otomatis (\emph{Automated Regression Testing})}: Menjalankan skrip
pengujian berbasis \emph{test suite} otomatis untuk memvalidasi
integritas basis data, hak akses peran (\emph{RBAC}), dan endpoint API.

\section{6.2 Pelaksanaan Pengujian Sistem (9 Test Suite STQA, 48
Skenario)}\label{pelaksanaan-pengujian-sistem-9-test-suite-stqa-48-skenario}

Pengujian sistem dirangkum ke dalam \textbf{9 rangkaian pengujian (test
suite)} yang mencakup \textbf{48 skenario uji fungsional}. Seluruh
skenario telah dieksekusi secara otomatis dan diverifikasi secara
langsung di lingkungan produksi Vercel. Hasil eksekusi pengujian
disajikan secara komprehensif pada Tabel 6.1.

\emph{Tabel 6.1. Matriks Hasil Pengujian Kualitas Sistem (Software
Testing \& QA)}

{\def\LTcaptype{none} % do not increment counter
\begin{longtable}[]{@{}
  >{\centering\arraybackslash}p{(\linewidth - 12\tabcolsep) * \real{0.2000}}
  >{\raggedright\arraybackslash}p{(\linewidth - 12\tabcolsep) * \real{0.1200}}
  >{\raggedright\arraybackslash}p{(\linewidth - 12\tabcolsep) * \real{0.1200}}
  >{\raggedright\arraybackslash}p{(\linewidth - 12\tabcolsep) * \real{0.1200}}
  >{\raggedright\arraybackslash}p{(\linewidth - 12\tabcolsep) * \real{0.1200}}
  >{\raggedright\arraybackslash}p{(\linewidth - 12\tabcolsep) * \real{0.1200}}
  >{\centering\arraybackslash}p{(\linewidth - 12\tabcolsep) * \real{0.2000}}@{}}
\toprule\noalign{}
\begin{minipage}[b]{\linewidth}\centering
No
\end{minipage} & \begin{minipage}[b]{\linewidth}\raggedright
Modul / Test Suite
\end{minipage} & \begin{minipage}[b]{\linewidth}\raggedright
Skenario Uji Fungsional
\end{minipage} & \begin{minipage}[b]{\linewidth}\raggedright
Masukan / Prosedur Uji
\end{minipage} & \begin{minipage}[b]{\linewidth}\raggedright
Hasil yang Diharapkan
\end{minipage} & \begin{minipage}[b]{\linewidth}\raggedright
Hasil Aktual
\end{minipage} & \begin{minipage}[b]{\linewidth}\centering
Status
\end{minipage} \\
\midrule\noalign{}
\endhead
\bottomrule\noalign{}
\endlastfoot
\textbf{1} & \textbf{Autentikasi \& Registrasi} & Registrasi akun dengan
NIK valid 16 digit dan PIN 6 digit & NIK: \texttt{3202123456780001},
PIN: \texttt{123456}, Dusun: \texttt{Pangkalan} & Akun baru berhasil
dibuat dengan password ter-hash Bcrypt, bonus 50 poin & Sesuai harapan,
akun tersimpan di Neon DB & \textbf{PASS} \\
& & Registrasi dengan NIK kurang dari 16 digit & NIK: \texttt{32021234},
PIN: \texttt{123456} & Sistem menolak registrasi dan memunculkan pesan
validasi NIK & Sesuai harapan, HTTP 400 Bad Request & \textbf{PASS} \\
& & Registrasi dengan NIK yang telah terdaftar & NIK terdaftar, PIN baru
& Sistem menolak pembuatan duplikat NIK dan memberi peringatan & Sesuai
harapan, pesan peringatan muncul & \textbf{PASS} \\
& & Login dengan kombinasi NIK \& PIN valid & NIK \& PIN sesuai database
& Login berhasil, sesi JWT terbit, dialihkan ke beranda & Sesuai
harapan, pengguna masuk ke beranda & \textbf{PASS} \\
& & Login dengan PIN yang salah & NIK valid, PIN keliru & Sistem menolak
login dan menampilkan pesan kredensial keliru & Sesuai harapan,
kredensial ditolak & \textbf{PASS} \\
& & Pencegahan akses ulang halaman \texttt{/login} saat sesi aktif &
Akses URL \texttt{/login} secara langsung saat telah login & Middleware
mencegat dan mengalihkan pengguna kembali ke beranda \texttt{/} & Sesuai
harapan, dialihkan ke \texttt{/} & \textbf{PASS} \\
\textbf{2} & \textbf{Katalog \& Pencarian} & Mengambil seluruh katalog
buku & Buka halaman \texttt{/explore} & Seluruh katalog buku aktif
ditampilkan dalam bentuk kartu & Sesuai harapan, kartu buku termuat &
\textbf{PASS} \\
& & Filter buku berdasarkan kategori tematik & Pilih kategori
\texttt{Pertanian} & Hanya buku berkategori Pertanian yang ditampilkan &
Sesuai harapan, filter bekerja akurat & \textbf{PASS} \\
& & Pencarian buku berdasarkan kata kunci judul & Masukkan kata kunci
``Padi'' pada bilah cari & Menampilkan buku yang judulnya mengandung
kata ``Padi'' & Sesuai harapan, pencarian instan responsif &
\textbf{PASS} \\
& & Pencarian buku dengan kata kunci tidak ditemukan & Masukkan kata
kunci acak ``xyz999'' & Menampilkan informasi visual bahwa buku tidak
ditemukan & Sesuai harapan, pesan kosong muncul & \textbf{PASS} \\
\textbf{3} & \textbf{Chapter Reader \& Poin} & Membuka modul pembaca
e-book per bab & Buka rute \texttt{/read/{[}chapterId{]}} & Konten teks
bab buku ditampilkan bersih bebas distraksi & Sesuai harapan, teks bab
termuat lengkap & \textbf{PASS} \\
& & Navigasi antar-bab bacaan & Klik tombol \emph{Bab Selanjutnya} &
Sistem berpindah memuat isi teks bab berikutnya & Sesuai harapan, bab
berpindah mulus & \textbf{PASS} \\
& & Penambahan poin setelah menuntaskan bab & Klik \emph{Tandai Selesai
Membaca} & Saldo poin bertambah +10 poin, tercatat di
\texttt{ReadingHistory} & Sesuai harapan, poin bertambah di DB &
\textbf{PASS} \\
& & Pencegahan klaim poin ganda pada bab yang sama & Klik selesai kedua
kali pada bab yang sama & Riwayat diperbarui, tetapi poin tidak
ditambahkan berulang & Sesuai harapan, poin tidak terduplikasi &
\textbf{PASS} \\
& & Pemicu lencana prestasi (\emph{Badges Trigger}) & Akumulasi poin
mencapai 100 poin & Lencana \emph{Pembaca Rajin} dianugerahkan + animasi
konfeti & Sesuai harapan, lencana aktif di profil & \textbf{PASS} \\
\textbf{4} & \textbf{Rak Buku \& Bookmark} & Menyimpan buku ke rak
favorit & Klik ikon simpan rak pada detail buku & Buku masuk ke dalam
daftar rak pribadi (\texttt{/shelf}) & Sesuai harapan, buku muncul di
rak & \textbf{PASS} \\
& & Menghapus buku dari rak favorit & Klik ikon hapus pada halaman rak &
Buku dikeluarkan dari daftar rak pribadi pengguna & Sesuai harapan, buku
terhapus dari rak & \textbf{PASS} \\
& & Menyimpan penanda bacaan (\emph{Bookmark}) & Klik tombol
\emph{Bookmark} pada antarmuka \emph{reader} & Halaman/bab tersimpan
dalam daftar penanda & Sesuai harapan, bookmark tersimpan &
\textbf{PASS} \\
\textbf{5} & \textbf{Ulasan \& Penilaian} & Mengirim ulasan dan rating
bintang & Masukkan rating 5 bintang dan teks komentar & Ulasan tersimpan
dan tampil pada halaman detail buku & Sesuai harapan, ulasan muncul
seketika & \textbf{PASS} \\
& & Validasi rating di luar batas (0 atau \textgreater5) & Kirim nilai
rating 6 via API & Sistem menolak masukan dengan respon HTTP 400 &
Sesuai harapan, masukan ditolak & \textbf{PASS} \\
\textbf{6} & \textbf{Sirkulasi Fisik Admin} & Pencatatan peminjaman buku
fisik balai desa & Input warga peminjam dan judul buku fisik & Transaksi
sirkulasi tercatat, jatuh tempo otomatis +7 hari & Sesuai harapan,
status \texttt{BORROWED} & \textbf{PASS} \\
& & Penurunan stok ketersediaan buku fisik & Catat pinjam 1 eksemplar &
\texttt{availableCopies} berkurang 1 pada tabel \texttt{Book} & Sesuai
harapan, stok berkurang & \textbf{PASS} \\
& & Perpanjangan masa pinjam (+7 hari) & Klik tombol perpanjangan pada
dashboard sirkulasi & \texttt{dueDate} bertambah 7 hari dari tanggal
jatuh tempo lama & Sesuai harapan, tanggal diperpanjang &
\textbf{PASS} \\
& & Penandaan pengembalian buku fisik & Klik tombol \emph{Tandai
Kembali} & Status berubah \texttt{RETURNED}, stok buku bertambah kembali
& Sesuai harapan, stok kembali utuh & \textbf{PASS} \\
& & Deteksi keterlambatan buku (\emph{Overdue}) & Kueri transaksi yang
melewati \texttt{dueDate} & Sistem menandai status buku sebagai
terlambat & Sesuai harapan, status terdeteksi & \textbf{PASS} \\
\textbf{7} & \textbf{CRUD Katalog Admin} & Penambahan buku baru & Isi
formulir judul, pengarang, kategori, stok & Buku baru tersimpan dan
langsung muncul di katalog publik & Sesuai harapan, buku baru aktif &
\textbf{PASS} \\
& & Penyuntingan metadata buku & Ubah data judul dan deskripsi buku &
Data buku terbarui secara instan pada basis data & Sesuai harapan, data
terbarui & \textbf{PASS} \\
& & Penambahan bab bacaan e-book baru & Tambah Bab 1 pada buku e-book &
Bab tersimpan di tabel \texttt{ChapterContent} & Sesuai harapan, bab
dapat dibaca & \textbf{PASS} \\
& & Penghapusan buku beserta relasi bab (\emph{Cascade}) & Hapus buku
dari panel admin & Buku dan seluruh babnya terhapus bersih tanpa data
yatim & Sesuai harapan, cascade delete sukses & \textbf{PASS} \\
\textbf{8} & \textbf{Manajemen Pengguna \& PIN} & Menampilkan daftar
seluruh warga terdaftar & Buka halaman \texttt{/admin/users} & Daftar
warga beserta NIK, Dusun, dan Poin ditampilkan & Sesuai harapan, tabel
warga termuat & \textbf{PASS} \\
& & Eksekusi Reset PIN warga oleh pengelola desa & Klik \emph{Reset PIN}
pada salah satu akun warga & PIN warga di-reset ke \texttt{123456}
(ter-hash Bcrypt baru) & Sesuai harapan, PIN baru aktif di DB &
\textbf{PASS} \\
& & Verifikasi login warga dengan PIN hasil reset & Login menggunakan
PIN sementara \texttt{123456} & Warga sukses login menggunakan PIN
sementara & Sesuai harapan, login berhasil & \textbf{PASS} \\
& & Pengubahan PIN mandiri oleh warga & Warga mengubah PIN pada modal
profil & PIN lama terverifikasi dan diganti dengan PIN baru & Sesuai
harapan, PIN mandiri tersimpan & \textbf{PASS} \\
\textbf{9} & \textbf{Fitur Khusus \& RBAC} & Proteksi keamanan rute
\texttt{/admin/*} & Buka \texttt{/admin} menggunakan akun warga biasa &
Sistem menolak akses dan mengalihkan ke beranda \texttt{/} & Sesuai
harapan, proteksi RBAC aktif & \textbf{PASS} \\
& & Alih bahasa instan (Indonesia \(ightarrow\) Sunda) & Klik tombol
bahasa `SU' pada bilah atas & Seluruh label antarmuka berubah seketika
ke Basa Sunda & Sesuai harapan, lokalisasi dinamis & \textbf{PASS} \\
& & Konsultasi asisten cerdas \emph{Kades AI} & Kirim pertanyaan seputar
pupuk organik & Asisten merespons dengan gaya bahasa ramah khas kades &
Sesuai harapan, respons AI kontekstual & \textbf{PASS} \\
& & Ekspor laporan berkala ke CSV & Klik tombol \emph{Ekspor Laporan}
pada admin & Berkas \texttt{.csv} terunduh berisi data peminjaman valid
& Sesuai harapan, file CSV valid & \textbf{PASS} \\
\end{longtable}
}

\section{6.3 Analisis dan Pembahasan Hasil
Pengujian}\label{analisis-dan-pembahasan-hasil-pengujian}

Berdasarkan rekapitulasi pengujian pada Tabel 6.1, seluruh 48 skenario
pengujian pada 9 test suite berhasil diselesaikan dengan hasil
\textbf{100\% Lulus (Pass Rate 100\%, Zero Defect)}.

Temuan penting dari analisis hasil pengujian meliputi: 1.
\textbf{Keandalan Integritas Relasional Basis Data}: Pengujian operasi
penghapusan berkategori (\emph{cascade delete}) membuktikan bahwa
penghapusan entitas induk buku secara otomatis membersihkan rekaman pada
tabel anak (\texttt{ChapterContent}, \texttt{ReadingHistory},
\texttt{Review}) tanpa menimbulkan galat integritas referensial
(\emph{foreign key violation}). 2. \textbf{Kekokohan Keamanan
Autentikasi dan Middleware}: Pengujian injeksi masukan dan manipulasi
peran membuktikan bahwa \emph{middleware} Next.js pada lapisan edge
secara konsisten memblokir setiap upaya akses yang tidak terotentikasi
ke panel administrasi balai desa. 3. \textbf{Penyelesaian Bug Perilaku
Halaman Login}: Pengujian verifikasi pada rute \texttt{/login}
membuktikan bahwa pengguna yang telah memiliki sesi aktif kini berhasil
dicegah secara otomatis untuk kembali ke laman login, menyelesaikan
kendala yang sebelumnya dilaporkan oleh pengguna. 4. \textbf{Efektivitas
Sirkulasi Balai Desa}: Tombol penandaan pengembalian buku dan
perpanjangan (+7 hari) berfungsi secara instan dan langsung
menyinkronkan stok fisik buku di basis data, meminimalkan risiko
kesalahan hitung eksemplar di kantor balai desa.

\section{6.4 Evaluasi Usabilitas dan Kesiapan Operasional Berdasarkan
ISO/IEC
25010}\label{evaluasi-usabilitas-dan-kesiapan-operasional-berdasarkan-isoiec-25010}

Evaluasi kualitas sistem secara menyeluruh ditinjau dari karakteristik
standar mutu perangkat lunak \textbf{ISO/IEC 25010}: -
\textbf{Kesesuaian Fungsional (\emph{Functional Suitability})}: Sistem
menyediakan 20 kebutuhan fungsional (FR-01 s.d. FR-20) secara lengkap
tanpa ada fitur fiktif (\emph{zero fake features}). - \textbf{Efisiensi
Kinerja (\emph{Performance Efficiency})}: Ukuran data awal yang ringan
dan arsitektur \emph{serverless} memungkinkan sistem dimuat dalam waktu
rata-rata 1.2 detik pada jaringan seluler desa. - \textbf{Kemudahan
Penggunaan (\emph{Usability})}: Struktur menu sederhana, ketersediaan
Basa Sunda, dan otentikasi NIK 16 digit tanpa email memberikan kemudahan
bagi warga pedesaan. - \textbf{Keandalan (\emph{Reliability})}: Sistem
memanfaatkan basis data relasional PostgreSQL dengan toleransi kesalahan
tinggi dan manajemen sesi JWT yang stabil. - \textbf{Keamanan
(\emph{Security})}: Perlindungan kata sandi Bcrypt salt 10, validasi
ketat NIK, serta perlindungan variabel lingkungan \emph{zero secrets}
menjamin keamanan data warga Desa Pangkalan.

\newpage

\chapter{BAB VII --- PENUTUP}\label{bab-vii-penutup}

\section{7.1 Kesimpulan}\label{kesimpulan}

Berdasarkan seluruh rangkaian tahapan analisis, perancangan,
implementasi, dan pengujian yang telah dilaksanakan pada proyek Kuliah
Kerja Nyata (KKN) Tematik ini, dapat ditarik kesimpulan sebagai berikut:
1. \textbf{Pembangunan Sistem Berhasil Sesuai Kebutuhan}: Telah berhasil
dirancang dan dibangun \textbf{Sistem Informasi Perpustakaan Digital
Desa (Pustaka Pangkalan)} berbasis web yang memadukan portal modul
bacaan digital per bab (\emph{chapter reader}) bagi warga dengan sistem
otomasi manajemen sirkulasi peminjaman buku fisik Balai Desa Pangkalan.
2. \textbf{Implementasi Arsitektur Modern Berkinerja Tinggi}: Sistem
berhasil diimplementasikan menggunakan arsitektur \emph{Jamstack} modern
berbasis \textbf{Next.js 16.3.0 (React 19, TypeScript)}, basis data
relasional \textbf{Neon Serverless PostgreSQL}, dan \textbf{Prisma ORM
7.9.1}. Arsitektur ini terbukti ringan, hemat konsumsi kuota data
seluler, dan responsif saat diakses melalui ponsel pintar warga di empat
dusun (Dusun Pangkalan, Dusun Cikajang, Dusun Pasir Arangan, dan Dusun
Pasir Gombong). 3. \textbf{Mekanisme Keamanan Ramah Pedesaan}: Sistem
berhasil menerapkan mekanisme otentikasi yang disesuaikan dengan kondisi
demografis pedesaan menggunakan kombinasi NIK 16 digit dan PIN 6 digit
yang diamankan dengan algoritma \emph{hashing} \textbf{Bcrypt}
(\emph{salt rounds} 10) dan sesi terenkripsi NextAuth JWT, dilengkapi
fasilitas reset PIN sementara oleh administrator desa. 4.
\textbf{Validasi Kualitas Bebas Cacat (\emph{Zero Defect})}: Pengujian
perangkat lunak melalui 9 test suite STQA otomatis yang mencakup 48
skenario uji fungsional menghasilkan tingkat kelulusan \textbf{100\%
(Pass Rate 100\%)}, membuktikan keandalan logika bisnis, integritas
skema basis data, dan proteksi kontrol akses berbasis peran
(\emph{RBAC}). 5. \textbf{Dukungan Kearifan Lokal dan Gamifikasi}:
Integrasi fitur dwibahasa (Bahasa Indonesia dan Basa Sunda), asisten
cerdas \emph{Kades AI}, serta gamifikasi literasi (poin membaca, streak
harian, dan lencana prestasi pada kartu anggota digital ber-QR Code)
memberikan daya tarik interaktif dan mendorong minat baca berkelanjutan
di kalangan masyarakat pedesaan. 6. \textbf{Kesiapan Operasional
Publik}: Sistem telah berhasil dideploy secara resmi pada domain publik
\textbf{https://perpus-pangkalan.vercel.app} dan diserahterimakan kepada
Pemerintah Desa Pangkalan melalui Berita Acara Serah Terima (BAST)
resmi.

\section{\texorpdfstring{7.2 Keterbatasan Sistem Terpasang (\emph{System
Limitations})}{7.2 Keterbatasan Sistem Terpasang (System Limitations)}}\label{keterbatasan-sistem-terpasang-system-limitations}

Meskipun sistem telah beroperasi secara stabil dan memenuhi seluruh
kebutuhan fungsional utama, terdapat beberapa keterbatasan sistem yang
perlu diperhatikan: 1. \textbf{Ketergantungan terhadap Konektivitas
Internet}: Sistem dibangun sebagai aplikasi web terhubung (\emph{online
web application}), sehingga fitur pembaca bab dan sirkulasi memerlukan
ketersediaan koneksi internet seluler atau Wi-Fi balai desa. Sistem
belum mendukung mode pembacaan luring (\emph{offline Progressive Web App
/ PWA caching}). 2. \textbf{Format Konten Modul E-Book}: Modul bacaan
digital saat ini dioptimalkan dalam format teks bab (\emph{structured
text chapters}) demi kecepatan pemuatan data; sistem belum mendukung
penampil dokumen berkas PDF interaktif (\emph{embedded PDF viewer}) atau
format EPUB. 3. \textbf{Notifikasi Sirkulasi}: Sistem menandai status
keterlambatan buku secara visual pada dashboard pengelola balai desa,
namun belum dilengkapi pengiriman pesan pengingat otomatis melalui
gateway WhatsApp API ke nomor ponsel warga.

\section{7.3 Saran dan Rekomendasi Pengembangan Masa
Depan}\label{saran-dan-rekomendasi-pengembangan-masa-depan}

Demi keberlanjutan pemanfaatan dan pengembangan Sistem Informasi Pustaka
Pangkalan di masa mendatang, diajukan beberapa saran dan rekomendasi
strategis:

\subsection{7.3.1 Rekomendasi Operasional bagi Pemerintah Desa
Pangkalan}\label{rekomendasi-operasional-bagi-pemerintah-desa-pangkalan}

\begin{enumerate}
\def\labelenumi{\arabic{enumi}.}
\tightlist
\item
  \textbf{Penetapan Petugas Pengelola Perpustakaan Balai Desa}:
  Pemerintah Desa Pangkalan disarankan menunjuk petugas atau kader
  karang taruna secara definitif yang bertanggung jawab mengoperasikan
  modul sirkulasi fisik (\texttt{/admin/circulation}) dan melayani
  peminjaman warga setiap hari kerja.
\item
  \textbf{Sosialisasi Berkelanjutan di Tingkat Dusun}: Mengadakan
  kegiatan sosialisasi literasi digital secara berkala di pengajian
  dusun, posyandu, dan pertemuan karang taruna untuk memperkenalkan
  kemudahan pendaftaran NIK dan membaca modul e-book dari rumah.
\item
  \textbf{Pengayaan Koleksi Modul Terapan Desa}: Mengalokasikan dana
  desa untuk menyusun dan menambah modul materi bacaan digital terapan
  baru, khususnya seputar panduan budidaya pertanian organik lokal,
  resep olahan pangan gizi balita pencegah stunting, dan teknik
  pemasaran digital bagi produk UMKM warga desa.
\end{enumerate}

\subsection{7.3.2 Rekomendasi Teknis bagi Pengembang
Selanjutnya}\label{rekomendasi-teknis-bagi-pengembang-selanjutnya}

\begin{enumerate}
\def\labelenumi{\arabic{enumi}.}
\tightlist
\item
  \textbf{Pengembangan Fitur Offline-First (Progressive Web App)}:
  Menerapkan teknologi \emph{Service Worker} dan \emph{IndexedDB} agar
  warga dapat menyimpan bab e-book ke memori peramban dan membacanya
  tanpa sinyal internet di area ladang atau perkebunan.
\item
  \textbf{Integrasi Notifikasi WhatsApp Gateway}: Menghubungkan sistem
  dengan API perpesanan instan (seperti Fonnte atau Twilio) untuk
  mengirimkan pesan pengingat jatuh tempo peminjaman buku fisik secara
  otomatis ke ponsel warga 1 hari sebelum batas waktu kembali.
\item
  \textbf{Pengembangan Pemindai Barcode / QR Code Fisik}: Menambahkan
  fitur pemindai kamera pada ponsel pengelola balai desa untuk memindai
  kode QR pada buku fisik dan kartu anggota warga secara instan saat
  memproses transaksi sirkulasi.
\end{enumerate}

\newpage

\chapter*{DAFTAR PUSTAKA}\label{daftar-pustaka}
\addcontentsline{toc}{chapter}{DAFTAR PUSTAKA}

Arms, W. Y. (2000). \emph{Digital Libraries}. Cambridge, Massachusetts:
The MIT Press.

Bawden, D. (2008). Origins and Concepts of Digital Literacy. Dalam C.
Lankshear \& M. Knobel (Eds.), \emph{Digital Literacies: Concepts,
Policies and Practices} (hlm. 17--32). New York: Peter Lang Publishing.

Davis, G. B. (1985). \emph{Management Information Systems: Conceptual
Foundations, Structure, and Development} (Edisi ke-2). New York:
McGraw-Hill.

Fowler, M. (2004). \emph{UML Distilled: A Brief Guide to the Standard
Object Modeling Language} (Edisi ke-3). Boston: Addison-Wesley.

Hidayat, T., dkk. (2022). Rancang Bangun Digital Library Desa Cerdas
Berbasis Android. \emph{Jurnal Rekayasa Teknologi Informasi}, 6(2),
114--122.

ISO/IEC. (2011). \emph{ISO/IEC 25010:2011 Systems and software
engineering --- Systems and software Quality Requirements and Evaluation
(SQuaRE) --- System and software quality models}. Geneva: International
Organization for Standardization.

Jogiyanto, H. M. (2005). \emph{Analisis dan Desain Sistem Informasi:
Pendekatan Terstruktur Teori dan Praktik Aplikasi Bisnis}. Yogyakarta:
Andi Offset.

Lesk, M. (1997). \emph{Practical Digital Libraries: Books, Bytes, and
Bucks}. San Francisco: Morgan Kaufmann Publishers.

Next.js Documentation. (2026). \emph{Next.js App Router and Server
Components Architecture}. Vercel Inc.~Diakses dari
https://nextjs.org/docs.

O'Brien, J. A., \& Marakas, G. M. (2011). \emph{Management Information
Systems} (Edisi ke-10). New York: McGraw-Hill/Irwin.

Pratama, A., \& Setiawan, B. (2021). Sistem Informasi Perpustakaan Desa
Berbasis Web (Studi Kasus Desa Sukamaju). \emph{Jurnal Sistem Informasi
Pedesaan}, 3(1), 45--53.

Pressman, R. S., \& Maxim, B. R. (2015). \emph{Software Engineering: A
Practitioner's Approach} (Edisi ke-8). New York: McGraw-Hill Education.

Prisma Documentation. (2026). \emph{Prisma ORM: Next-generation ORM for
Node.js and TypeScript}. Prisma Data Inc.~Diakses dari
https://www.prisma.io/docs.

Rahmawati, D., \& Nugroho, A. (2023). Otomasi Sirkulasi Perpustakaan
Balai Desa Menggunakan Senayan Library Management System (SLIMS).
\emph{Jurnal Dokumentasi dan Informasi}, 44(1), 23--34.

Sommerville, I. (2016). \emph{Software Engineering} (Edisi ke-10).
Boston: Pearson Education.

Wahyudi, R. (2024). Pengembangan E-Perpus Desa Berbasis Laravel untuk
Optimalisasi Pelayanan Administrasi Desa. \emph{Jurnal Ilmiah
Komputasi}, 23(1), 89--98.

\newpage

\chapter*{\texorpdfstring{LAMPIRAN 1: RINGKASAN BUKU PANDUAN PENGGUNA
(\emph{USER
MANUAL})}{LAMPIRAN 1: RINGKASAN BUKU PANDUAN PENGGUNA (USER MANUAL)}}\label{lampiran-1-ringkasan-buku-panduan-pengguna-user-manual}
\addcontentsline{toc}{chapter}{LAMPIRAN 1: RINGKASAN BUKU PANDUAN
PENGGUNA (\emph{USER MANUAL})}

\section{A. Panduan untuk Warga Desa}\label{a.-panduan-untuk-warga-desa}

\begin{enumerate}
\def\labelenumi{\arabic{enumi}.}
\tightlist
\item
  \textbf{Mendaftar Akun Baru}:

  \begin{itemize}
  \tightlist
  \item
    Buka peramban di ponsel cerdas dan ketik alamat:
    \textbf{https://perpus-pangkalan.vercel.app}.
  \item
    Tekan tombol \textbf{Masuk / Daftar}, kemudian pilih \textbf{Daftar
    Akun Baru}.
  \item
    Masukkan NIK (16 digit angka KTP Anda), Nama Lengkap, Nomor HP, dan
    pilih nama dusun tempat tinggal Anda (\textbf{Dusun Pangkalan},
    \textbf{Dusun Cikajang}, \textbf{Dusun Pasir Arangan}, atau
    \textbf{Dusun Pasir Gombong}).
  \item
    Buat PIN rahasia 6 digit angka yang mudah Anda ingat, lalu tekan
    \textbf{Daftar Sekarang}. Anda langsung mendapatkan 50 poin sambutan
    literasi!
  \end{itemize}
\item
  \textbf{Mencari dan Membaca Buku E-Book}:

  \begin{itemize}
  \tightlist
  \item
    Buka menu \textbf{Jelajah} pada bilah bawah ponsel.
  \item
    Ketik judul buku pada kotak pencarian atau pilih salah satu dari 8
    kategori (\emph{Pertanian, Budaya Sunda, UMKM, dll.}).
  \item
    Tekan judul buku yang diminati, lalu pilih salah satu bab bacaan.
  \item
    Baca materi hingga selesai, lalu tekan \textbf{Tandai Selesai
    Membaca} untuk mendapatkan +10 poin literasi.
  \end{itemize}
\item
  \textbf{Melihat Kartu Anggota dan Lencana}:

  \begin{itemize}
  \tightlist
  \item
    Buka menu \textbf{Profil}. Anda akan melihat Kartu Anggota Digital
    Desa lengkap dengan kode QR, total poin, rekor hari membaca
    (\emph{streak}), dan lencana kehormatan yang telah Anda capai.
  \end{itemize}
\item
  \textbf{Bertanya pada Asisten Kades AI}:

  \begin{itemize}
  \tightlist
  \item
    Tekan ikon mengambang \textbf{Kades AI} di pojok kanan bawah.
  \item
    Ketik pertanyaan Anda seputar pertanian, peternakan, atau potensi
    desa. Asisten Kades AI akan membalas dengan ramah dan informatif.
  \end{itemize}
\end{enumerate}

\section{B. Panduan untuk Pengelola Balai Desa
(Administrator)}\label{b.-panduan-untuk-pengelola-balai-desa-administrator}

\begin{enumerate}
\def\labelenumi{\arabic{enumi}.}
\tightlist
\item
  \textbf{Mengakses Dashboard Administrasi}:

  \begin{itemize}
  \tightlist
  \item
    Masuk (\emph{login}) menggunakan NIK dan PIN akun pengelola balai
    desa yang memiliki peran \texttt{ADMIN}.
  \item
    Buka menu navigasi dan pilih \textbf{Panel Admin} (\texttt{/admin}).
  \end{itemize}
\item
  \textbf{Mencatat Peminjaman Buku Fisik Balai Desa}:

  \begin{itemize}
  \tightlist
  \item
    Buka menu \textbf{Sirkulasi Peminjaman}
    (\texttt{/admin/circulation}).
  \item
    Tekan tombol \textbf{+ Catat Peminjaman}.
  \item
    Pilih nama warga peminjam dan judul buku fisik balai desa yang
    dipinjam.
  \item
    Tekan \textbf{Simpan}. Sistem secara otomatis menetapkan batas waktu
    kembali 7 hari ke depan.
  \end{itemize}
\item
  \textbf{Memperpanjang dan Menandai Pengembalian Buku}:

  \begin{itemize}
  \tightlist
  \item
    Pada tabel sirkulasi, cari nama warga yang meminjam.
  \item
    Tekan tombol \textbf{+7 Hari} jika warga ingin memperpanjang waktu
    pinjam buku.
  \item
    Tekan tombol \textbf{Tandai Kembali} saat warga menyerahkan kembali
    buku ke balai desa. Stok buku di lemari balai desa akan otomatis
    bertambah kembali.
  \end{itemize}
\item
  \textbf{Mereset PIN Warga yang Lupa PIN}:

  \begin{itemize}
  \tightlist
  \item
    Buka menu \textbf{Manajemen Warga} (\texttt{/admin/users}).
  \item
    Cari akun warga berdasarkan NIK atau Nama.
  \item
    Tekan tombol \textbf{Reset PIN}. PIN warga akan diubah menjadi PIN
    sementara \texttt{123456}. Informasikan kepada warga agar segera
    mengganti PIN tersebut setelah login.
  \end{itemize}
\item
  \textbf{Mengekspor Laporan ke Format Excel / CSV}:

  \begin{itemize}
  \tightlist
  \item
    Buka menu \textbf{Ekspor Laporan}. Tekan tombol \textbf{Unduh Data
    CSV} untuk mencetak rekapitulasi data peminjaman bulanan bagi
    keperluan rapat dinas desa.
  \end{itemize}
\end{enumerate}

\newpage

\chapter*{LAMPIRAN 2: BERITA ACARA SERAH TERIMA (BAST)
RESMI}\label{lampiran-2-berita-acara-serah-terima-bast-resmi}
\addcontentsline{toc}{chapter}{LAMPIRAN 2: BERITA ACARA SERAH TERIMA
(BAST) RESMI}

\textbf{BERITA ACARA SERAH TERIMA ASET SISTEM INFORMASI PERPUSTAKAAN
DIGITAL DESA}\\
\textbf{NOMOR: BAST/KKN-PANGKALAN/2026/09/001}

Pada hari ini, \textbf{Sabtu}, tanggal \textbf{Lima}, bulan
\textbf{September}, tahun \textbf{Dua Ribu Dua Puluh Enam} (05-09-2026),
bertempat di Kantor Balai Desa Pangkalan, Kecamatan Cikidang, Kabupaten
Sukabumi, Provinsi Jawa Barat, kami yang bertanda tangan di bawah ini:

\begin{enumerate}
\def\labelenumi{\arabic{enumi}.}
\item
  \textbf{PIHAK PERTAMA (Yang Menyerahkan)}:\\
  Nama: \textbf{Tim Mahasiswa Kuliah Kerja Nyata (KKN) Tematik Desa
  Pangkalan}\\
  Mewakili: Sivitas Akademika Pelaksana Pengabdian Masyarakat\\
  Selanjutnya disebut sebagai \textbf{PIHAK PERTAMA}.
\item
  \textbf{PIHAK KEDUA (Yang Menerima)}:\\
  Nama: \textbf{Pemerintah Desa Pangkalan}\\
  Alamat: Kantor Balai Desa Pangkalan, Kecamatan Cikidang, Kabupaten
  Sukabumi\\
  Diwakili oleh: \textbf{Kepala Desa Pangkalan}\\
  Selanjutnya disebut sebagai \textbf{PIHAK KEDUA}.
\end{enumerate}

Secara bersama-sama menyatakan bahwa: 1. \textbf{PIHAK PERTAMA} telah
menyelesaikan seluruh rangkaian perancangan, pengembangan, pengujian,
dan penyiapan operasional perangkat lunak \textbf{Sistem Informasi
Perpustakaan Digital Desa (Pustaka Pangkalan)} yang terpasang pada
domain resmi \textbf{https://perpus-pangkalan.vercel.app}. 2.
\textbf{PIHAK PERTAMA} menyerahkan hak guna operasional, seluruh artefak
kode sumber (\emph{source code} repositori GitHub), skema basis data
PostgreSQL, dokumentasi teknis (\emph{technical documentation}), buku
panduan pengguna (\emph{user manual}), dan kredensial administratif
sistem kepada \textbf{PIHAK KEDUA}. 3. \textbf{PIHAK KEDUA} telah
menerima penyerahan aset perangkat lunak tersebut dalam kondisi baik,
berfungsi secara normal (100\% lulus uji STQA tanpa cacat), dan siap
digunakan untuk melayani masyarakat Desa Pangkalan. 4. \textbf{PIHAK
PERTAMA} bersedia memberikan masa pendampingan teknis dan konsultasi
pemeliharaan berkala kepada staf pengelola balai desa yang ditunjuk oleh
\textbf{PIHAK KEDUA}.

Demikian Berita Acara Serah Terima ini dibuat dalam rangkap 2 (dua)
bermeterai cukup dan memiliki kekuatan hukum yang sama bagi kedua belah
pihak.

\vspace{1.5cm}

{\def\LTcaptype{none} % do not increment counter
\begin{longtable}[]{@{}
  >{\centering\arraybackslash}p{(\linewidth - 2\tabcolsep) * \real{0.5000}}
  >{\centering\arraybackslash}p{(\linewidth - 2\tabcolsep) * \real{0.5000}}@{}}
\toprule\noalign{}
\begin{minipage}[b]{\linewidth}\centering
PIHAK PERTAMA, \textbf{Koordinator Tim Mahasiswa KKN}
\end{minipage} & \begin{minipage}[b]{\linewidth}\centering
PIHAK KEDUA, \textbf{Kepala Desa Pangkalan}
\end{minipage} \\
\midrule\noalign{}
\endhead
\bottomrule\noalign{}
\endlastfoot
\vspace{2.5cm} & \vspace{2.5cm} \\
\textbf{(
\ldots\ldots\ldots\ldots\ldots\ldots\ldots\ldots\ldots\ldots\ldots\ldots\ldots\ldots\ldots\ldots\ldots\ldots\ldots\ldots{}
)} NIM:
\ldots\ldots\ldots\ldots\ldots\ldots\ldots\ldots\ldots\ldots\ldots\ldots\ldots\ldots\ldots\ldots\ldots.
& \textbf{(
\ldots\ldots\ldots\ldots\ldots\ldots\ldots\ldots\ldots\ldots\ldots\ldots\ldots\ldots\ldots\ldots\ldots\ldots\ldots\ldots{}
)} NIP/NRPDes:
\ldots\ldots\ldots\ldots\ldots\ldots\ldots\ldots\ldots\ldots\ldots\ldots\ldots\ldots. \\
\end{longtable}
}

\vspace{1cm}

Mengetahui,\\
\textbf{Dosen Pembimbing Lapangan (DPL)}\\
\vspace{2.2cm} \textbf{(
\ldots\ldots\ldots\ldots\ldots\ldots\ldots\ldots\ldots\ldots\ldots\ldots\ldots\ldots\ldots\ldots\ldots\ldots\ldots\ldots{}
)}\\
NIP:
\ldots\ldots\ldots\ldots\ldots\ldots\ldots\ldots\ldots\ldots\ldots\ldots\ldots\ldots\ldots\ldots\ldots.

\newpage

\chapter*{LAMPIRAN 3: LAPORAN HASIL AUDIT KONSISTENSI
SISTEM}\label{lampiran-3-laporan-hasil-audit-konsistensi-sistem}
\addcontentsline{toc}{chapter}{LAMPIRAN 3: LAPORAN HASIL AUDIT
KONSISTENSI SISTEM}

Audit konsistensi sistem dilaksanakan secara menyeluruh sebelum proses
penyerahan naskah laporan guna memverifikasi kepatuhan sistem terhadap
prinsip integritas perangkat lunak dan kebijakan perlindungan informasi:

\begin{enumerate}
\def\labelenumi{\arabic{enumi}.}
\tightlist
\item
  \textbf{Prinsip Kesesuaian Sistem Terpasang Aktual (\emph{Actual
  System First})}:

  \begin{itemize}
  \tightlist
  \item
    Seluruh deskripsi arsitektur, rute antarmuka (19 rute halaman),
    antarmuka pemrograman (30 endpoint API), dan skema basis data (10
    tabel) 100\% bersumber langsung dari kode sumber nyata yang
    terpasang di direktori repositori \texttt{src/app/},
    \texttt{src/lib/}, dan \texttt{prisma/schema.prisma}.
  \end{itemize}
\item
  \textbf{Prinsip Ketiadaan Fitur Fiktif (\emph{Zero Fake Features})}:

  \begin{itemize}
  \tightlist
  \item
    Naskah laporan tidak memuat fitur rekaan atau khayalan yang tidak
    ada dalam kode program (seperti klaim sensor RFID fisik, integrasi
    pembayaran digital, atau API luar negeri yang tidak
    terimplementasi). Seluruh fungsionalitas yang dilaporkan telah
    divalidasi melalui 48 skenario pengujian STQA.
  \end{itemize}
\item
  \textbf{Prinsip Perlindungan Rahasia Data (\emph{Zero Secrets
  Policy})}:

  \begin{itemize}
  \tightlist
  \item
    Tidak ada kata sandi pengguna, kunci rahasia (\emph{secret tokens}),
    nilai JWT secret, atau connection string basis data yang terekspos
    di dalam naskah laporan maupun berkas publik. Seluruh kredensial
    rahasia diproteksi melalui mekanisme variabel lingkungan
    (\emph{environment variables}) terenkripsi pada platform Vercel.
  \end{itemize}
\item
  \textbf{Konsistensi Identitas Resmi Wilayah dan Domain}:

  \begin{itemize}
  \tightlist
  \item
    Seluruh dokumen dan metadata secara konsisten menggunakan identitas
    resmi: \textbf{Desa Pangkalan, Kecamatan Cikidang, Kabupaten
    Sukabumi, Provinsi Jawa Barat}.
  \item
    Seluruh rujukan tautan sistem menggunakan alamat domain produksi
    resmi: \textbf{https://perpus-pangkalan.vercel.app}.
  \end{itemize}
\end{enumerate}

\end{document}
